%!TeX root=MemoriaTFG.tex

\chapter{Resultados}
\label{cap:resultados}

\section{Reproducción de PanDerm: sondeo lineal sobre diez datasets}
\label{sec:rep-panderm}

La tabla~\ref{tab:lp-panderm} reporta el sondeo lineal con PanDerm
Base y Large sobre los diez datasets externos contrastados, conforme
al protocolo descrito en secci\'on~\ref{sec:lp}. Las cinco métricas
reportadas son \emph{accuracy} (Acc), \emph{balanced accuracy}
(BAcc), AUROC \emph{one-vs-rest}, F1 ponderada (W-F1) y coeficiente
kappa de Cohen.

\begin{table}[H]
\centering
\footnotesize
\setlength{\tabcolsep}{4pt}
\renewcommand{\arraystretch}{0.95}
\caption{Sondeo lineal con PanDerm Base y Large sobre diez datasets
externos. El asterisco (\textsuperscript{$\ast$}) marca datasets con
solapamiento documentado frente a Derm1M (secci\'on~\ref{sec:auditoria-resultados}).
Mejores valores por dataset y métrica en negrita.}
\label{tab:lp-panderm}
\begin{tabular}{llrrrrr}
\hline
Dataset       & Modelo & Acc            & BAcc           & AUROC          & W-F1           & Kappa \\
\hline
HAM10000      & Base   & $0{,}853$      & $0{,}381$      & $0{,}892$      & $0{,}829$      & $0{,}709$ \\
HAM10000      & Large  & $\mathbf{0{,}888}$ & $\mathbf{0{,}575}$ & $\mathbf{0{,}954}$ & $\mathbf{0{,}880}$ & $\mathbf{0{,}814}$ \\
BCN20000      & Base   & $0{,}659$      & $0{,}292$      & $0{,}855$      & $0{,}615$      & $0{,}447$ \\
BCN20000      & Large  & $\mathbf{0{,}702}$ & $\mathbf{0{,}382}$ & $\mathbf{0{,}903}$ & $\mathbf{0{,}676}$ & $\mathbf{0{,}527}$ \\
PAD-UFES-20   & Base   & $0{,}725$      & $0{,}510$      & $0{,}913$      & $0{,}692$      & $0{,}369$ \\
PAD-UFES-20   & Large  & $\mathbf{0{,}772}$ & $\mathbf{0{,}642}$ & $\mathbf{0{,}949}$ & $\mathbf{0{,}760}$ & $\mathbf{0{,}527}$ \\
Dermnet$^\ast$& Base   & $0{,}450$      & $0{,}381$      & $0{,}888$      & $0{,}432$      & $0{,}437$ \\
Dermnet$^\ast$& Large  & $\mathbf{0{,}550}$ & $\mathbf{0{,}495}$ & $\mathbf{0{,}931}$ & $\mathbf{0{,}540}$ & $\mathbf{0{,}555}$ \\
WSI patches   & Base   & $0{,}781$      & $0{,}640$      & $0{,}976$      & $0{,}774$      & $0{,}842$ \\
WSI patches   & Large  & $\mathbf{0{,}868}$ & $\mathbf{0{,}792}$ & $\mathbf{0{,}991}$ & $\mathbf{0{,}871}$ & $\mathbf{0{,}896}$ \\
DDI           & Base   & $\mathbf{0{,}847}$ & $0{,}714$      & $0{,}827$      & $0{,}836$      & $0{,}482$ \\
DDI           & Large  & $\mathbf{0{,}847}$ & $\mathbf{0{,}764}$ & $\mathbf{0{,}860}$ & $\mathbf{0{,}846}$ & $\mathbf{0{,}535}$ \\
Derm7pt clín. & Base   & $0{,}762$      & $0{,}675$      & $0{,}808$      & $0{,}736$      & $0{,}397$ \\
Derm7pt clín. & Large  & $\mathbf{0{,}798}$ & $\mathbf{0{,}732}$ & $\mathbf{0{,}869}$ & $\mathbf{0{,}784}$ & $\mathbf{0{,}506}$ \\
Derm7pt dermo & Base   & $0{,}818$      & $0{,}749$      & $0{,}841$      & $0{,}811$      & $0{,}528$ \\
Derm7pt dermo & Large  & $\mathbf{0{,}836}$ & $\mathbf{0{,}766}$ & $\mathbf{0{,}858}$ & $\mathbf{0{,}829}$ & $\mathbf{0{,}570}$ \\
HIBA$^\ast$   & Base   & $0{,}883$      & $0{,}595$      & $0{,}894$      & $0{,}853$      & $0{,}275$ \\
HIBA$^\ast$   & Large  & $\mathbf{0{,}904}$ & $\mathbf{0{,}701}$ & $\mathbf{0{,}942}$ & $\mathbf{0{,}892}$ & $\mathbf{0{,}494}$ \\
MSKCC$^\ast$  & Base   & $0{,}721$      & $\mathbf{0{,}654}$ & $0{,}746$      & $0{,}720$      & $0{,}309$ \\
MSKCC$^\ast$  & Large  & $\mathbf{0{,}746}$ & $0{,}644$      & $\mathbf{0{,}751}$ & $\mathbf{0{,}732}$ & $\mathbf{0{,}316}$ \\
\hline
\emph{Media}  & Base   & $0{,}750$      & $0{,}559$      & $0{,}865$      & $0{,}730$      & --- \\
\emph{Media}  & Large  & $\mathbf{0{,}791}$ & $\mathbf{0{,}649}$ & $\mathbf{0{,}901}$ & $\mathbf{0{,}781}$ & --- \\
\hline
\end{tabular}
\end{table}

PanDerm Large supera a PanDerm Base en accuracy, AUROC y F1
ponderada sobre los diez datasets, con mejora media absoluta de
$+4{,}1$\,pp en accuracy, $+9{,}0$\,pp en BAcc y
$+3{,}6$\,pp en AUROC. La única excepción es la BAcc en MSKCC,
donde Base obtiene $0{,}654$ frente a $0{,}644$ de Large
($-1{,}0$\,pp). La mejora absoluta es mayor sobre Dermnet
($+10{,}0$\,pp accuracy), WSI~patches ($+8{,}7$\,pp) y
PAD-UFES-20 ($+4{,}7$\,pp), datasets con mayor número de clases o
mayor heterogeneidad de adquisición. La modalidad con mejor
rendimiento absoluto es la histopatología (WSI~patches: Acc
$0{,}868$, AUROC $0{,}991$, kappa $0{,}896$). Los tres datasets
marcados con asterisco (Dermnet, HIBA, MSKCC) presentan solapamiento
con Derm1M; su interpretación se discute en
secci\'on~\ref{sec:auditoria-resultados}.

\begin{figure}[H]
\centering
\includegraphics[width=0.92\textwidth]{fig_lp_per_dataset.png}
\caption{Sondeo lineal con PanDerm Base y Large sobre los diez
datasets externos contrastados. Diferencia absoluta en BAcc
($+9{,}0$\,pp de promedio) y AUROC ($+3{,}6$\,pp de
promedio) a favor de PanDerm Large.}
\label{fig:lp-per-dataset}
\end{figure}

\section{Benchmark comparativo entre familias}
\label{sec:benchmark}

La tabla~\ref{tab:benchmark} compara cinco modelos representativos
de tres familias (encoder específico de dominio, encoder generalista
auto-supervisado y arquitecturas convolucionales preentrenadas en
ImageNet) sobre HAM10000, PAD-UFES-20 y DDI bajo el mismo protocolo
de sondeo lineal.

\begin{table}[H]
\centering
\footnotesize
\setlength{\tabcolsep}{4pt}
\renewcommand{\arraystretch}{0.95}
\caption{Benchmark comparativo bajo sondeo lineal sobre HAM10000,
PAD-UFES-20 y DDI. Mejor valor por columna en negrita.}
\label{tab:benchmark}
\begin{tabular}{lrrrrrrr}
\hline
Modelo                  & Tipo          & Params & HAM Acc       & HAM AUROC     & PAD Acc       & PAD AUROC     & DDI AUROC \\
\hline
ConvNeXt-Large          & ConvNet sup.  & $197$\,M & $0{,}883$ & $\mathbf{0{,}967}$ & $0{,}668$ & $0{,}894$ & $0{,}774$ \\
EfficientNetV2-Large    & ConvNet sup.  & $118$\,M & $0{,}860$ & $0{,}950$ & $0{,}690$ & $0{,}895$ & $0{,}766$ \\
DINOv2 ViT-L/14         & SSL general   & $307$\,M & $0{,}744$ & $0{,}847$ & ---       & ---       & --- \\
PanDerm Large           & SSL derm.     & $307$\,M & $0{,}888$ & $0{,}954$ & $\mathbf{0{,}772}$ & $\mathbf{0{,}949}$ & $\mathbf{0{,}860}$ \\
SigLIP-Large SO400M     & Vis-leng.\ gen.& $400$\,M & $\mathbf{0{,}900}$ & $0{,}971$ & $0{,}718$ & $0{,}903$ & $0{,}793$ \\
\hline
\end{tabular}
\end{table}

La tabla~\ref{tab:gap-cnn} cuantifica la ventaja de PanDerm Large
sobre la mejor CNN supervisada en cada dataset.

\begin{table}[H]
\centering
\footnotesize
\setlength{\tabcolsep}{4pt}
\caption{PanDerm Large frente a la mejor CNN supervisada (la mejor de
ConvNeXt-L y EfficientNetV2-L para cada métrica).}
\label{tab:gap-cnn}
\begin{tabular}{lrrrrrr}
\hline
Dataset     & CNN Acc       & PanDerm Acc   & $\Delta$Acc (pp) & CNN AUROC     & PanDerm AUROC & $\Delta$AUROC (pp) \\
\hline
HAM10000    & $0{,}883$     & $0{,}888$     & $+0{,}5$         & $0{,}967$     & $0{,}954$     & $-1{,}3$ \\
PAD-UFES-20 & $0{,}690$     & $0{,}772$     & $+8{,}2$         & $0{,}895$     & $0{,}949$     & $+5{,}4$ \\
DDI         & $0{,}818$     & $0{,}847$     & $+2{,}9$         & $0{,}774$     & $0{,}860$     & $+8{,}6$ \\
\hline
\end{tabular}
\end{table}

PanDerm Large supera a la mejor CNN supervisada en accuracy sobre
los tres datasets, con margen creciente conforme aumenta la
dificultad: $+0{,}5$\,pp en HAM10000 (dermatoscopia
estandarizada), $+8{,}2$\,pp en PAD-UFES-20 (fotografía clínica
con dispositivo móvil) y $+2{,}9$\,pp en DDI (fotografía clínica
con representación equilibrada de fototipos). En AUROC el
patrón es similar excepto en HAM10000, donde ConvNeXt-Large
($0{,}967$) supera a PanDerm Large por $1{,}3$\,pp. SigLIP-Large
SO400M, modelo vision-lenguaje generalista, alcanza la mejor
accuracy y AUROC en HAM10000 ($0{,}900 / 0{,}971$), comportamiento
consistente con la sobrerrepresentación de dermatoscopia
estandarizada en su preentrenamiento sobre web. PanDerm Large
domina en PAD-UFES-20 y DDI, datasets clínicos con mayor distancia
respecto a los corpus web generalistas.

\begin{figure}[H]
\centering
\includegraphics[width=0.88\textwidth]{fig_lp_comparison.png}
\caption{Benchmark comparativo bajo sondeo lineal sobre HAM10000,
PAD-UFES-20 y DDI. ConvNeXt-Large supera a PanDerm Large en AUROC
sobre HAM10000 ($-1{,}3$\,pp); PanDerm Large domina sobre
PAD-UFES-20 ($+8{,}2$\,pp Acc) y DDI ($+8{,}6$\,pp AUROC).}
\label{fig:lp-benchmark}
\end{figure}

\section{Ajuste fino supervisado}
\label{sec:ft-results}

La tabla~\ref{tab:ft-panderm} compara el sondeo lineal y el ajuste
fino supervisado con PanDerm Base sobre cuatro datasets, con la
configuración descrita en secci\'on~\ref{sec:ft}. Sobre HAM10000 se aplica
adicionalmente \emph{test-time augmentation} con cinco augmentaciones
deterministas.

\begin{table}[H]
\centering
\footnotesize
\setlength{\tabcolsep}{4pt}
\renewcommand{\arraystretch}{0.95}
\caption{Sondeo lineal (LP) frente a ajuste fino (FT) con PanDerm
Base. $\Delta$BAcc en puntos porcentuales respecto al LP.}
\label{tab:ft-panderm}
\begin{tabular}{llrrrrr}
\hline
Dataset      & Modo & Acc           & BAcc          & AUROC         & W-F1          & $\Delta$BAcc \\
\hline
HAM10000     & LP   & $0{,}853$     & $0{,}381$     & $0{,}892$     & $0{,}829$     & --- \\
HAM10000     & FT   & $\mathbf{0{,}920}$ & $\mathbf{0{,}852}$ & $\mathbf{0{,}978}$ & $\mathbf{0{,}923}$ & $+47{,}1$ \\
PAD-UFES-20  & LP   & $0{,}725$     & $0{,}510$     & $0{,}913$     & $0{,}692$     & --- \\
PAD-UFES-20  & FT   & $\mathbf{0{,}755}$ & $\mathbf{0{,}704}$ & $\mathbf{0{,}935}$ & $\mathbf{0{,}757}$ & $+19{,}4$ \\
DDI          & LP   & $\mathbf{0{,}847}$ & $\mathbf{0{,}714}$ & $\mathbf{0{,}827}$ & $\mathbf{0{,}836}$ & --- \\
DDI          & FT   & $0{,}774$     & $0{,}693$     & $0{,}682$     & $0{,}780$     & $-2{,}1$ \\
MSKCC        & LP   & $0{,}721$     & $0{,}654$     & $\mathbf{0{,}746}$ & $0{,}720$     & --- \\
MSKCC        & FT   & $\mathbf{0{,}731}$ & $\mathbf{0{,}684}$ & $0{,}719$     & $\mathbf{0{,}735}$ & $+3{,}0$ \\
\hline
\end{tabular}
\end{table}

El ajuste fino aumenta sustancialmente la BAcc en HAM10000
($+47{,}1$\,pp) y PAD-UFES-20 ($+19{,}4$\,pp). El incremento se
atribuye al aprendizaje de las clases minoritarias (melanoma,
\emph{akiec}, \emph{df}, \emph{vasc} en HAM10000), favorecido por
el \emph{weighted random sampler} y los regularizadores Mixup y
CutMix. Sobre DDI, con $N_{\mathrm{train}} = 427$, el sondeo lineal
supera al ajuste fino en todas las métricas, comportamiento
compatible con sobreajuste del encoder ajustado pese a las
augmentaciones. Sobre MSKCC ($N_{\mathrm{train}} = 6\,189$), la
mejora de BAcc es marginal ($+3{,}0$\,pp) y la AUROC desciende
ligeramente. Estos dos casos delimitan el régimen de aplicabilidad
del ajuste fino: por debajo de aproximadamente $500$ imágenes
de entrenamiento, el sondeo lineal sobre PanDerm Large constituye
una alternativa preferible.

\begin{figure}[H]
\centering
\includegraphics[width=0.85\textwidth]{fig_ft_comparison.png}
\caption{Comparación entre sondeo lineal y ajuste fino con
PanDerm Base sobre cuatro datasets. El ajuste fino aporta
mejoras sustanciales sobre HAM10000 y PAD-UFES-20
($+47{,}1$\,pp y $+19{,}4$\,pp BAcc respectivamente) pero
degrada el rendimiento sobre DDI ($N_{\mathrm{train}} = 427$),
caso compatible con sobreajuste pese a las augmentaciones.}
\label{fig:ft-comparison}
\end{figure}

\section{Eficiencia en etiquetas}
\label{sec:label-eff-results}

La tabla~\ref{tab:label-eff} reporta el rendimiento de PanDerm Large
mediante sondeo lineal con submuestreos estratificados al
$1\%$, $5\%$, $10\%$, $20\%$, $50\%$ y $100\%$
del conjunto de entrenamiento sobre HAM10000 y PAD-UFES-20.

\begin{table}[H]
\centering
\footnotesize
\setlength{\tabcolsep}{4pt}
\renewcommand{\arraystretch}{0.95}
\caption{Eficiencia en etiquetas (sondeo lineal con PanDerm Large).
$N_{\mathrm{train}}$: tamaño efectivo del subconjunto de
entrenamiento. Mejor valor por columna en negrita.}
\label{tab:label-eff}
\begin{tabular}{lrrrrr}
\hline
$\%$ entrenam.  & $N_{\mathrm{train}}$ & Acc           & BAcc          & AUROC         & W-F1 \\
\hline
\multicolumn{6}{l}{\textit{HAM10000}} \\
$1\%$           & $82$                 & $0{,}850$     & $0{,}406$     & $0{,}918$     & $0{,}828$ \\
$5\%$           & $410$                & $0{,}870$     & $0{,}481$     & $0{,}932$     & $0{,}855$ \\
$10\%$          & $820$                & $0{,}875$     & $0{,}480$     & $0{,}939$     & $0{,}863$ \\
$20\%$          & $1\,641$             & $0{,}882$     & $0{,}547$     & $0{,}951$     & $0{,}873$ \\
$50\%$          & $4\,103$             & $\mathbf{0{,}894}$ & $\mathbf{0{,}589}$ & $0{,}952$ & $\mathbf{0{,}887}$ \\
$100\%$         & $8\,207$             & $0{,}888$     & $0{,}575$     & $\mathbf{0{,}954}$ & $0{,}880$ \\
\hline
\multicolumn{6}{l}{\textit{PAD-UFES-20}} \\
$1\%$           & $14$                 & $0{,}740$     & $0{,}599$     & $0{,}907$     & $0{,}727$ \\
$5\%$           & $74$                 & $0{,}751$     & $0{,}632$     & $0{,}919$     & $0{,}742$ \\
$10\%$          & $149$                & $0{,}738$     & $0{,}593$     & $0{,}923$     & $0{,}724$ \\
$20\%$          & $298$                & $0{,}748$     & $0{,}587$     & $0{,}931$     & $0{,}728$ \\
$50\%$          & $746$                & $0{,}757$     & $0{,}611$     & $0{,}941$     & $0{,}742$ \\
$100\%$         & $1\,493$             & $\mathbf{0{,}772}$ & $\mathbf{0{,}642}$ & $\mathbf{0{,}949}$ & $\mathbf{0{,}760}$ \\
\hline
\end{tabular}
\end{table}

Sobre HAM10000, con $1\%$ del conjunto de entrenamiento
($82$ imágenes), PanDerm Large recupera $0{,}850$ accuracy y
$0{,}918$ AUROC, $3{,}8$\,pp y $3{,}6$\,pp respectivamente
por debajo del rendimiento con el conjunto completo. La curva de
AUROC satura tempranamente: con $5\%$ del conjunto
($410$ imágenes) se alcanza $0{,}932$ AUROC, equivalente al
$97{,}7\%$ del rendimiento máximo. La BAcc, en cambio, sigue
escalando con el volumen y alcanza su máximo en $50\%$
($0{,}589$). Sobre PAD-UFES-20, con $14$ imágenes de
entrenamiento, el modelo alcanza $0{,}740$ accuracy y $0{,}907$
AUROC, $95{,}5\%$ del rendimiento de AUROC con el conjunto
completo. El comportamiento es consistente con encoders
preentrenados sobre corpus masivos del dominio: las representaciones
son ya informativas y el clasificador lineal converge con un
número reducido de ejemplos.

\begin{figure}[H]
\centering
\includegraphics[width=0.88\textwidth]{fig_label_efficiency.png}
\caption{Eficiencia en etiquetas con PanDerm Large sobre HAM10000
y PAD-UFES-20. La curva AUROC satura tempranamente: con
$1\%$ del conjunto de entrenamiento se recupera más del
$95\%$ del rendimiento alcanzable con el conjunto completo
en ambos datasets.}
\label{fig:label-efficiency}
\end{figure}

\section{Segmentación binaria de lesión}
\label{sec:seg-results}

La tabla~\ref{tab:seg} reporta los resultados de segmentación
binaria sobre ISIC2018 con SAM2.1-Large adaptado mediante LoRA,
junto a tres referencias: SAM2 sin adaptación
(\emph{zero-shot}), \emph{Convolutional AutoEncoder} con backbone
PanDerm Large entrenado desde cero (CAE-seg) y la cifra publicada
por Yan et~al.~\cite{panderm2025}.

\begin{table}[H]
\centering
\footnotesize
\setlength{\tabcolsep}{4pt}
\renewcommand{\arraystretch}{0.95}
\caption{Segmentación binaria de lesión: Dice de test sobre tres
datasets dermatoscópicos. Modelos entrenados sobre ISIC2018 y
evaluados sin reentrenamiento sobre ISIC2017 y PH2.}
\label{tab:seg}
\begin{tabular}{lrrr}
\hline
Configuración                         & ISIC2018 Dice & ISIC2017 Dice & PH2 Dice \\
\hline
SAM2 \emph{zero-shot} (bbox de GT)    & ---           & $0{,}797$     & $0{,}886$ \\
CAE-seg (PanDerm-L) FT ISIC2018       & $0{,}894$     & $0{,}879$     & $0{,}924$ \\
PanDerm \cite{panderm2025} (ref.)     & $0{,}921$     & ---           & --- \\
\textbf{SAM2.1-L + LoRA FT ISIC2018}  & $\mathbf{0{,}947}$ & $\mathbf{0{,}945}$ & $\mathbf{0{,}960}$ \\
\hline
\end{tabular}
\end{table}

La tabla~\ref{tab:seg-detalle} amplía la información con la
desviación estándar y el IoU sobre los datasets de transferencia.

\begin{table}[H]
\centering
\footnotesize
\setlength{\tabcolsep}{4pt}
\caption{Generalización fuera de dominio: Dice e IoU (media $\pm$
desv.~estándar) sobre ISIC2017 ($N = 600$) y PH2 ($N = 200$).}
\label{tab:seg-detalle}
\begin{tabular}{llrr}
\hline
Modelo                              & Dataset      & Dice              & IoU \\
\hline
SAM2 \emph{zero-shot} (bbox)        & ISIC2017     & $0{,}797 \pm 0{,}158$ & $0{,}686 \pm 0{,}183$ \\
SAM2 \emph{zero-shot} (bbox)        & PH2          & $0{,}886 \pm 0{,}199$ & $0{,}830 \pm 0{,}198$ \\
SAM2.1-L + LoRA FT ISIC2018         & ISIC2017     & $0{,}945 \pm 0{,}046$ & $0{,}898 \pm 0{,}065$ \\
SAM2.1-L + LoRA FT ISIC2018         & PH2          & $0{,}960 \pm 0{,}022$ & $0{,}925 \pm 0{,}039$ \\
CAE-seg (PanDerm-L) FT ISIC2018     & ISIC2017     & $0{,}879 \pm 0{,}139$ & $0{,}803 \pm 0{,}166$ \\
CAE-seg (PanDerm-L) FT ISIC2018     & PH2          & $0{,}924 \pm 0{,}066$ & $0{,}865 \pm 0{,}098$ \\
\hline
\end{tabular}
\end{table}

SAM2.1-L con adaptación LoRA alcanza Dice $0{,}947$ sobre
ISIC2018, superando en $+2{,}6$\,pp la cifra publicada por
Yan et~al.~\cite{panderm2025} ($0{,}921$) y en $+5{,}3$\,pp al
CAE-seg con backbone PanDerm-L. La transferencia fuera de dominio es
mínima: Dice desciende sólo $0{,}002$ sobre ISIC2017
($0{,}947 \to 0{,}945$) y, sobre PH2, asciende hasta
$0{,}960$, comportamiento explicable por la mayor limpieza
visual del dataset PH2. El paso de SAM2 sin adaptación a SAM2.1-L
con LoRA introduce mejoras de $+14{,}8$\,pp Dice sobre ISIC2017
($0{,}797 \to 0{,}945$) y $+7{,}4$\,pp sobre PH2
($0{,}886 \to 0{,}960$), cuantificando el aporte específico de la
adaptación de dominio.

\begin{figure}[H]
\centering
\includegraphics[width=0.88\textwidth]{fig_seg_generalization.png}
\caption{Segmentación binaria de lesión. SAM2.1-Large con
adaptación LoRA alcanza Dice $0{,}947$ sobre ISIC2018,
$+2{,}6$\,pp sobre el resultado publicado por
Yan~et~al.~\cite{panderm2025}, con transferencia OOD
prácticamente nula sobre ISIC2017 ($0{,}945$) y PH2
($0{,}960$).}
\label{fig:seg-generalization}
\end{figure}

\section{Clasificación zero-shot multimodal}
\label{sec:zs-results}

La tabla~\ref{tab:zs-ham} reporta los resultados de clasificación
zero-shot multimodal sobre HAM10000 ($7$ clases) con tres
configuraciones de \emph{prompts} y dos tokenizadores distintos.

\begin{table}[H]
\centering
\footnotesize
\setlength{\tabcolsep}{4pt}
\renewcommand{\arraystretch}{0.95}
\caption{Zero-shot multimodal sobre HAM10000 ($N_{\mathrm{test}} = 1\,232$,
$7$ clases). Mejor valor por columna en negrita.}
\label{tab:zs-ham}
\begin{tabular}{llllrr}
\hline
Modelo                    & Tokenizador  & \emph{Prompts}     & Acc       & AUROC     & F1-macro \\
\hline
DermLIP v1/v2             & PubMedBERT   & Genéricos          & $0{,}086$ & $0{,}366$ & $0{,}023$ \\
DermLIP v1/v2             & PubMedBERT   & Derm1M A3          & ---       & $0{,}516$ & --- \\
DermLIP original          & GPT-2/CLIP   & Derm1M A3          & $\mathbf{0{,}427}$ & $\mathbf{0{,}854}$ & --- \\
\hline
\end{tabular}
\end{table}

El paso de \emph{prompts} genéricos en inglés con tokenizador
PubMedBERT al conjunto A3 de Derm1M con tokenizador GPT-2/CLIP
eleva la AUROC desde $0{,}366$ hasta $0{,}854$
($+48{,}8$\,pp). Los \emph{prompts} genéricos con tokenizador
PubMedBERT obtienen AUROC$<0{,}5$, valor compatible con proyección
en un espacio no alineado con el del encoder visual de la variante
v1/v2. El cambio del conjunto A3 sin cambio de tokenizador eleva
AUROC a $0{,}516$, mientras que el cambio de tokenizador
manteniendo A3 alcanza $0{,}854$, lo que identifica la
compatibilidad del tokenizador con el espacio de \emph{embeddings}
como el factor crítico.

La tabla~\ref{tab:zs-binaria} reporta la evaluación zero-shot binaria
melanoma/no-melanoma sobre DDI, MSKCC, HIBA y Derm7pt clínico con
DermLIP v1/v2 y \emph{prompts} genéricos.

\begin{table}[H]
\centering
\footnotesize
\setlength{\tabcolsep}{4pt}
\caption{Zero-shot binario melanoma/no-melanoma con DermLIP v1/v2
(\emph{prompts} genéricos).}
\label{tab:zs-binaria}
\begin{tabular}{lrrrrr}
\hline
Dataset           & $N_{\mathrm{test}}$ & Acc       & BAcc      & AUROC     & F1-macro \\
\hline
DDI               & $137$               & $0{,}766$ & $0{,}499$ & $0{,}436$ & $0{,}463$ \\
MSKCC             & $1\,664$            & $0{,}716$ & $0{,}500$ & $0{,}430$ & $0{,}417$ \\
HIBA              & $334$               & $0{,}859$ & $0{,}497$ & $0{,}502$ & $0{,}462$ \\
Derm7pt clín.     & $168$               & $0{,}589$ & $0{,}472$ & $0{,}337$ & $0{,}439$ \\
\hline
\end{tabular}
\end{table}

En la formulación binaria, las cifras de BAcc próximas a
$0{,}500$ sobre los cuatro datasets indican que la rama
contrastiva opera al nivel del azar, comportamiento esperable dada
la sensibilidad documentada al tokenizador. El detalle del análisis
extendido sobre sensibilidad a \emph{prompts} y tokenizador se
desarrolla en el Anexo~\ref{cap:anexog}.

\begin{figure}[H]
\centering
\includegraphics[width=0.78\textwidth]{fig_zs_prompts.png}
\caption{Sensibilidad zero-shot al tokenizador y a la
formulación del \emph{prompt} sobre HAM10000. La AUROC oscila
entre $0{,}366$ (DermLIP v1/v2 + PubMedBERT + prompts
genéricos) y $0{,}854$ (DermLIP original +
GPT-2/CLIP + prompts A3 de Derm1M), un rango de
$48{,}8$\,pp para variaciones que no afectan al encoder
visual.}
\label{fig:zs-prompts}
\end{figure}

\section{Comparación con modelos de lenguaje multimodales}
\label{sec:llm-results}

La tabla~\ref{tab:llm-comparativa} reúne la evaluación de seis
modelos de lenguaje multimodales sobre HAM10000, PAD-UFES-20 y DDI,
junto a los modelos de referencia obtenidos en
secci\'on~\ref{sec:rep-panderm} y secci\'on~\ref{sec:ft-results}.

\begin{table}[H]
\centering
\footnotesize
\setlength{\tabcolsep}{4pt}
\renewcommand{\arraystretch}{0.95}
\caption{Comparación de paradigmas sobre HAM10000, PAD-UFES-20 y DDI
($N_{\mathrm{test}} = 1\,232 / 461 / 137$). FT: ajuste fino, LP:
sondeo lineal, ZS: zero-shot.}
\label{tab:llm-comparativa}
\begin{tabular}{llrrrrrr}
\hline
Modelo                 & Modo  & HAM Acc       & HAM BAcc      & PAD Acc       & PAD BAcc      & DDI Acc       & DDI BAcc \\
\hline
PanDerm Base (TTA HAM) & FT    & $\mathbf{0{,}920}$ & $\mathbf{0{,}852}$ & $0{,}755$     & $\mathbf{0{,}704}$ & $0{,}774$     & $0{,}693$ \\
PanDerm Large          & LP    & $0{,}888$     & $0{,}575$     & $\mathbf{0{,}772}$ & $0{,}642$     & $\mathbf{0{,}847}$ & $\mathbf{0{,}764}$ \\
SigLIP-Large           & LP    & $0{,}900$     & ---           & $0{,}718$     & ---           & $0{,}789$     & --- \\
MedGemma 27B + LoRA    & FT    & $0{,}802$     & $0{,}270$     & $0{,}430$     & $0{,}347$     & $0{,}657$     & $0{,}618$ \\
MedGemma 27B           & ZS    & $0{,}665$     & $0{,}243$     & $0{,}447$     & $0{,}342$     & $0{,}474$     & $0{,}566$ \\
GPT-4o                 & ZS    & $0{,}485$     & $0{,}465$     & $0{,}553$     & $0{,}523$     & $0{,}723$     & $0{,}648$ \\
GPT-4o-mini            & ZS    & $0{,}256$     & $0{,}319$     & $0{,}390$     & $0{,}435$     & $0{,}737$     & $0{,}556$ \\
Gemini 2.5 Pro         & ZS    & $0{,}403$     & ---           & $0{,}477$     & ---           & $0{,}441$     & --- \\
Gemini 2.5 Flash       & ZS    & $0{,}380$     & ---           & $0{,}486$     & ---           & $0{,}555$     & --- \\
DermLIP                & ZS    & $0{,}239$     & $0{,}581$     & $0{,}601$     & ---           & $0{,}329$     & --- \\
BLIP-2 (Flan-T5-XL)    & ZS    & $0{,}015$     & $0{,}145$     & $0{,}017$     & $0{,}167$     & $0{,}796$     & $0{,}542$ \\
InstructBLIP (Flan-T5-XL)& ZS  & $0{,}024$     & $0{,}123$     & $0{,}171$     & $0{,}224$     & $0{,}788$     & $0{,}500$ \\
\hline
\end{tabular}
\end{table}

La jerarquía de rendimiento sobre HAM10000 se ordena de mayor a
menor en: PanDerm Base FT con TTA ($0{,}920$ Acc) $>$
SigLIP-Large LP ($0{,}900$) $>$ PanDerm Large LP ($0{,}888$) $>$
MedGemma~27B con LoRA ($0{,}802$) $>$ MedGemma~27B ZS
($0{,}665$) $>$
GPT-4o ZS ($0{,}485$) $>$ Gemini~2.5~Pro ZS ($0{,}403$) $>$
Gemini~2.5~Flash ZS ($0{,}380$) $>$ GPT-4o-mini ZS ($0{,}256$)
$>$ DermLIP ZS ($0{,}239$) $\gg$ InstructBLIP/BLIP-2 ZS
($\sim 0{,}02$). La adaptación con LoRA sobre MedGemma~27B sobre
el \emph{split} de entrenamiento de HAM10000 produce mejoras de
$+13{,}7$\,pp Acc sobre HAM10000 y $+18{,}3$\,pp Acc sobre
DDI respecto a la variante \emph{zero-shot}, sin mejora apreciable
sobre PAD-UFES-20 ($-1{,}7$\,pp). El coste agregado de la
evaluación con APIs comerciales es de aproximadamente
\$$4{,}07$ entre GPT-4o ($\sim\$2{,}71$) y
GPT-4o-mini ($\sim\$1{,}36$). MedGemma~27B opera localmente en
BF16 sobre los $128$\,GB de memoria unificada de la DGX Spark
sin cuantización adicional, con latencia media de $1{,}49$
segundos por imagen.

La tabla~\ref{tab:gpt4o-perclass} muestra el desglose por clase del
modelo GPT-4o sobre HAM10000, ilustrando patrones de error
asimétricos: la clase \emph{actinic keratosis} obtiene
\emph{recall} $0{,}000$ y \emph{dermatofibroma} recibe
$284$ predicciones positivas frente a $8$ muestras reales,
en línea con un sesgo hacia categorías menos restringidas en el
\emph{prompt} clínico.

\begin{table}[H]
\centering
\footnotesize
\setlength{\tabcolsep}{4pt}
\caption{Desglose de \emph{recall} por clase sobre HAM10000 con
GPT-4o zero-shot. Predicciones: número total de imágenes asignadas a
cada clase por el modelo.}
\label{tab:gpt4o-perclass}
\begin{tabular}{lrrr}
\hline
Clase                       & $N_{\mathrm{test}}$ & Predicciones GPT-4o & Recall \\
\hline
actinic keratosis           & $35$    & $2$     & $0{,}000$ \\
basal cell carcinoma        & $44$    & ---     & $0{,}477$ \\
seborrheic keratosis        & $107$   & ---     & $0{,}308$ \\
dermatofibroma              & $8$     & $284$   & $0{,}750$ \\
melanoma                    & $70$    & $299$   & $0{,}686$ \\
melanocytic nevus           & $951$   & ---     & $0{,}505$ \\
vascular lesion             & $17$    & ---     & $0{,}529$ \\
\hline
\end{tabular}
\end{table}

\begin{figure}[H]
\centering
\includegraphics[width=0.92\textwidth]{fig_llm_comparison.png}
\caption{Comparativa de paradigmas sobre HAM10000. Gradiente
de $43{,}5$\,pp accuracy entre el mejor especialista
ajustado (PanDerm Base FT con TTA, $0{,}920$) y el mejor LLM
multimodal generalista \emph{zero-shot} (GPT-4o,
$0{,}485$). BLIP-2 e InstructBLIP descienden hasta accuracy
$\sim 0{,}02$.}
\label{fig:llm-comparison}
\end{figure}

\section{Equidad por fototipo cutáneo}
\label{sec:equidad-results}

La tabla~\ref{tab:equidad-global} reporta el rendimiento global de
cinco encoders sobre Fitzpatrick17k completo ($N = 16\,577$) en
las dos formulaciones consideradas: clasificación a $114$
patologías (grano fino del dataset) y clasificación a $3$ clases
de malignidad (benigno / no~neoplásico / maligno).

\begin{table}[H]
\centering
\footnotesize
\setlength{\tabcolsep}{4pt}
\renewcommand{\arraystretch}{0.95}
\caption{Equidad sobre Fitzpatrick17k completo. Rendimiento global
en las dos formulaciones. Mejor valor por columna en negrita.}
\label{tab:equidad-global}
\begin{tabular}{lrrrrrrrr}
\hline
                          & \multicolumn{4}{c}{$114$ patologías} & \multicolumn{4}{c}{$3$ clases malignidad} \\
\cline{2-5}\cline{6-9}
Modelo                    & Acc           & BAcc          & AUROC         & W-F1          & Acc           & BAcc          & AUROC         & W-F1 \\
\hline
PanDerm Large             & $\mathbf{0{,}581}$ & $0{,}549$     & $0{,}968$     & $\mathbf{0{,}576}$ & $0{,}835$     & $0{,}699$     & $0{,}908$     & $0{,}828$ \\
PanDerm Base              & $0{,}500$     & $0{,}478$     & $0{,}955$     & $0{,}491$     & $0{,}819$     & $0{,}657$     & $0{,}882$     & $0{,}807$ \\
DermLIP v2 (visual)       & $0{,}572$     & $\mathbf{0{,}549}$ & $\mathbf{0{,}971}$ & $0{,}569$ & $\mathbf{0{,}852}$ & $\mathbf{0{,}721}$ & $\mathbf{0{,}909}$ & $\mathbf{0{,}846}$ \\
DINOv2 ViT-L/14           & $0{,}519$     & $0{,}485$     & $0{,}947$     & $0{,}516$     & $0{,}809$     & $0{,}656$     & $0{,}854$     & $0{,}801$ \\
BiomedCLIP                & $0{,}334$     & $0{,}321$     & $0{,}910$     & $0{,}323$     & $0{,}799$     & $0{,}590$     & $0{,}830$     & $0{,}776$ \\
\hline
\end{tabular}
\end{table}

La tabla~\ref{tab:equidad-fototipo} reporta la BAcc por fototipo
Fitzpatrick sobre la formulación de tres clases.

\begin{table}[H]
\centering
\footnotesize
\setlength{\tabcolsep}{4pt}
\renewcommand{\arraystretch}{0.95}
\caption{BAcc por subgrupo Fitzpatrick (3 clases de malignidad) y
\emph{gap} entre fototipos extremos. Tamaños muestrales por
subgrupo: $299 / 470 / 325 / 261 / 169 / 73$ para FST~I a~VI.}
\label{tab:equidad-fototipo}
\begin{tabular}{lrrrrrrr}
\hline
Modelo                    & FST I         & FST II        & FST III       & FST IV        & FST V         & FST VI        & Gap I-VI \\
\hline
PanDerm Large             & $0{,}723$     & $0{,}708$     & $0{,}743$     & $0{,}669$     & $0{,}568$     & $0{,}506$     & $+0{,}217$ \\
PanDerm Base              & $0{,}653$     & $0{,}693$     & $0{,}669$     & $0{,}630$     & $0{,}585$     & $0{,}384$     & $+0{,}269$ \\
DermLIP v2 (visual)       & $0{,}716$     & $0{,}730$     & $0{,}765$     & $0{,}710$     & $0{,}693$     & $0{,}434$     & $+0{,}281$ \\
DINOv2 ViT-L/14           & $0{,}673$     & $0{,}671$     & $0{,}679$     & $0{,}654$     & $0{,}553$     & $0{,}368$     & $+0{,}306$ \\
BiomedCLIP                & $0{,}569$     & $0{,}602$     & $0{,}606$     & $0{,}566$     & $0{,}590$     & $\mathbf{0{,}445}$ & $\mathbf{+0{,}124}$ \\
\hline
\end{tabular}
\end{table}

Los cinco modelos presentan \emph{gap}~I-VI positivo, indicando
rendimiento inferior sobre fototipos elevados. DermLIP v2 obtiene la
mejor BAcc global ($0{,}721$ sobre $3$ clases) y la mejor AUROC
global ($0{,}909$), pero su \emph{gap}~I-VI ($+0{,}281$) es
superior al de PanDerm Large ($+0{,}217$). BiomedCLIP presenta el
\emph{gap} menor ($+0{,}124$) pero al precio de una BAcc global
inferior ($0{,}590$), lo que ilustra el carácter no decisional
del \emph{gap} aislado: un modelo uniformemente peor entre
subgrupos puede presentar \emph{gap} bajo sin ser equitativo en
sentido sustantivo. La AUROC de PanDerm Large sobre la formulación
de $3$ clases es $0{,}908$ global, con \emph{gap} AUROC entre
fototipos extremos de $0{,}089$ (FST~IV: $0{,}919$ vs FST~VI:
$0{,}830$).

La tabla~\ref{tab:cross-eval} amplía la evaluación con un
experimento cruzado \emph{Light}~$\leftrightarrow$~\emph{Dark}, donde
\emph{Light} agrupa FST~I-III y \emph{Dark} agrupa FST~IV-VI.

\begin{table}[H]
\centering
\footnotesize
\setlength{\tabcolsep}{4pt}
\caption{Evaluación cruzada \emph{Light}~$\leftrightarrow$~\emph{Dark}
(BAcc, formulación de $3$ clases). \emph{L}: \emph{Light};
\emph{D}: \emph{Dark}; \emph{All}: ambos grupos.}
\label{tab:cross-eval}
\begin{tabular}{lrrrrrr}
\hline
Modelo                    & L$\to$L       & L$\to$D       & D$\to$D       & D$\to$L       & All$\to$L     & All$\to$D \\
\hline
PanDerm Large             & $0{,}699$     & $0{,}589$     & $0{,}701$     & $0{,}667$     & $0{,}726$     & $0{,}627$ \\
PanDerm Base              & $0{,}673$     & $0{,}557$     & $0{,}621$     & $0{,}629$     & $0{,}673$     & $0{,}594$ \\
DermLIP v2 (visual)       & $\mathbf{0{,}729}$ & $\mathbf{0{,}651}$ & $0{,}667$ & $\mathbf{0{,}683}$ & $\mathbf{0{,}745}$ & $\mathbf{0{,}682}$ \\
DINOv2 ViT-L/14           & $0{,}655$     & $0{,}523$     & $0{,}625$     & $0{,}569$     & $0{,}682$     & $0{,}589$ \\
BiomedCLIP                & $0{,}595$     & $0{,}532$     & $0{,}571$     & $0{,}551$     & $0{,}594$     & $0{,}565$ \\
\hline
\end{tabular}
\end{table}

La caída de rendimiento entrenando sólo con fototipos claros y
evaluando sobre fototipos oscuros (L$\to$D vs L$\to$L) es mayor en
DINOv2 ($\Delta = -0{,}132$\,pp BAcc) y menor en BiomedCLIP
($\Delta = -0{,}063$\,pp). DermLIP v2 mantiene la menor caída
Dark$\to$Light vs Dark$\to$Dark ($\Delta = +0{,}016$\,pp), lo que
sugiere mayor invariancia de sus representaciones al fototipo cutáneo.

\begin{figure}[H]
\centering
\includegraphics[width=0.95\textwidth]{figura_fairness_5models.png}
\caption{Equidad por fototipo Fitzpatrick (formulación de tres
clases de malignidad) sobre los cinco encoders evaluados. Los
cinco presentan \emph{gap}~I-VI positivo (entre $+0{,}124$ y
$+0{,}306$\,pp BAcc). Paradoja BiomedCLIP: \emph{gap} mínimo
asociado a peor BAcc global, lo que demuestra el carácter no
decisional de la métrica aislada.}
\label{fig:fairness-5models}
\end{figure}

\begin{figure}[H]
\centering
\includegraphics[width=0.85\textwidth]{fig_cross_eval.png}
\caption{Evaluación cruzada
\emph{Light}~$\leftrightarrow$~\emph{Dark} sobre Fitzpatrick17k
(BAcc en formulación de tres clases). La caída
L$\to$D vs L$\to$L cuantifica la pérdida de rendimiento al
entrenar sólo con fototipos claros y evaluar sobre fototipos
oscuros, máxima en DINOv2 ($\Delta = -0{,}132$\,pp).}
\label{fig:cross-eval}
\end{figure}

\section{Auditoría de solapamiento con Derm1M}
\label{sec:auditoria-resultados}

La auditoría por hash MD5 entre los \emph{splits} de test
contrastados y la fracción accesible del corpus Derm1M identifica
tres datasets con solapamiento documentado: Dermnet, HIBA y
MSKCC. \textbf{[NO ESPECIFICADO: completar las cardinalidades
exactas de la intersección por hash MD5 cuando se ejecute el script
de auditoría. Los tres datasets afectados aparecen marcados con
asterisco en las tablas~\ref{tab:lp-panderm} y~\ref{tab:benchmark}.]}

La presencia de solapamiento no invalida las cifras reportadas en
secci\'on~\ref{sec:rep-panderm} para PanDerm, DINOv2, ConvNeXt-Large,
EfficientNetV2-Large y SigLIP-Large (modelos cuyo preentrenamiento
no incluye Derm1M), pero condiciona la interpretación de cualquier
cifra zero-shot reportada en secci\'on~\ref{sec:zs-results} para la familia
DermLIP sobre los datasets afectados.

\section{Ensemble de seguridad para detección de melanoma}
\label{sec:ensemble-results}

Como protocolo auxiliar, se evalúa la combinación de tres
clasificadores independientes sobre el \emph{split} de test de
HAM10000 ($N = 1\,232$ imágenes, $70$ melanomas) en una
configuración orientada a maximizar el \emph{recall} de melanoma:
M1 (PanDerm Large \emph{fine-tuned} sobre HAM10000), SigLIP-Large
SO400M con sondeo lineal sobre los \emph{embeddings} de
$1\,152$ dimensiones, y M7 (clasificador unificado sobre el
\emph{merged43} del corpus armonizado, descrito en el
Anexo~\ref{cap:anexoh}).

\begin{table}[H]
\centering
\footnotesize
\setlength{\tabcolsep}{4pt}
\renewcommand{\arraystretch}{0.95}
\caption{Ensemble sobre HAM10000 ($N = 1\,232$, $70$ melanomas).
\emph{Recall} \emph{top-3}: la clase melanoma aparece entre las
tres primeras del ranking. FP\,mel: falsos positivos melanoma
(imágenes no-melanoma marcadas).}
\label{tab:ensemble}
\begin{tabular}{lrrrrrr}
\hline
Configuración                   & Acc           & BAcc          & AUROC         & Mel rec.\ top-1 & Mel rec.\ top-3 & FP\,mel \\
\hline
M1 (PanDerm FT)                 & $0{,}919$     & $0{,}813$     & $0{,}986$     & $0{,}714$ (50/70) & $0{,}971$ (68/70) & $36$ \\
SigLIP-L LP                     & $0{,}900$     & $0{,}776$     & $0{,}971$     & $0{,}614$ (43/70) & $0{,}986$ (69/70) & $31$ \\
Avg M1 + SigLIP                 & $0{,}917$     & $0{,}836$     & $0{,}986$     & $0{,}686$         & $0{,}986$         & $32$ \\
MaxMel M1 + SigLIP              & $0{,}909$     & $0{,}832$     & $0{,}985$     & $\mathbf{0{,}771}$ & $\mathbf{1{,}000}$ & $50$ \\
OR top-3 M1 + SigLIP            & ---           & ---           & ---           & ---               & $\mathbf{1{,}000}$ (70/70) & $885$ \\
\hline
\end{tabular}
\end{table}

La combinación OR \emph{top-3} de M1 y SigLIP alcanza
\emph{recall} de melanoma del $100\%$ ($70$ de $70$)
sobre el conjunto de test de HAM10000. El análisis de errores
identifica que los dos melanomas que M1 no captura en
\emph{top-3} (ISIC\_0024886.jpg e ISIC\_0030360.jpg) sí son
capturados por SigLIP, lo que confirma la complementariedad de los
\emph{backbones} y motiva la elección del par. M7, evaluado en el
subconjunto de $198$ imágenes con predicción simultánea de los
tres modelos ($12$ melanomas), no añade melanomas adicionales
respecto al par M1$+$SigLIP, pero su valor se mantiene en escenarios
multidataset por su cobertura sobre las $43$ clases
ontológicas del corpus armonizado.

Tres caveats acompañan la interpretación de este resultado: (i)~el
\emph{recall} del $100\%$ es \emph{top-3}, no \emph{top-1};
(ii)~el tamaño muestral de melanomas es $N = 70$, con intervalos
de confianza al $95\%$ amplios; (iii)~no existe validación
prospectiva sobre cohorte hospitalaria. El resultado se interpreta
como prueba de concepto para un escenario de cribado oportunista,
no como sustituto del juicio clínico.

\section{Validación en castellano sobre DermapixelAI 1.0}
\label{sec:dermapixel-results}

A diferencia de los datasets externos contrastados en
secci\'on~\ref{sec:rep-panderm}, tres de los cuales presentan
solapamiento documentado con Derm1M
(secci\'on~\ref{sec:auditoria-resultados}), el corpus DermapixelAI~1.0
descrito en el Anexo~\ref{cap:anexoe} es completamente independiente
del corpus de preentrenamiento de la familia PanDerm/DermLIP. La
auditoría por hash MD5 sobre $507\,657$ imágenes catalogadas
(Derm1M $415\,451$ ítems, HAM10000, BCN20000, ISIC2017/2018,
PH2, PAD-UFES-20, DDI, Dermnet, Fitzpatrick17k, HIBA, MSKCC y
Derm7pt) reporta intersección nula: $0$ coincidencias sobre las
$1\,062$ imágenes evaluables del corpus. La evaluación reportada en
esta sección no está, por tanto, sujeta a la fuga de información
caracterizada en secci\'on~\ref{sec:auditoria-resultados}.

Se aplica el protocolo de sondeo lineal descrito en
secci\'on~\ref{sec:lp} sobre las $1\,062$ imágenes con
\texttt{label\_source=ontology} y \texttt{ontology\_l1} no vacío,
particionadas en train ($N = 874$), validación ($N = 152$) y test
($N = 36$). PanDerm Base (768-D) y PanDerm Large (1024-D) operan
como extractores congelados; la cabeza es un clasificador
\emph{LogisticRegression} con solver L-BFGS, $C = 1{,}0$ y
$\mathrm{max\_iter} = 5\,000$. Los intervalos de confianza al
$95\%$ se calculan por \emph{bootstrap} estratificado por clase
con $1\,000$ remuestreos. La evaluación L2 se reporta sobre las
$38$ subcategorías clínicas efectivamente representadas en el
corpus tras consolidar las dos variantes ortográficas de
\emph{Trastornos de la queratinización} (Anexo~\ref{cap:anexoc},
secci\'on~\ref{sec:anexo-l2-l3}), de las $43$ entradas del
vocabulario ontológico original.

\begin{table}[H]
\centering
\footnotesize
\setlength{\tabcolsep}{4pt}
\renewcommand{\arraystretch}{0.95}
\caption{Sondeo lineal con PanDerm Base y Large sobre
DermapixelAI~1.0 en los tres niveles de la ontología jerárquica
($N_{\mathrm{train}} = 874$, $N_{\mathrm{test}} = 36$ en L1/L2 y
$N_{\mathrm{test}} = 28$ en L3 tras filtrar clases ausentes en
train). Entre corchetes, IC95\% por \emph{bootstrap} estratificado
($1\,000$ remuestreos). Mejor valor por columna y nivel en
negrita.}
\label{tab:dermapixel-lp}
\begin{tabular}{llrrrrr}
\hline
Nivel & Modelo & Acc & BAcc & AUROC & W-F1 & Kappa \\
\hline
L1 (4 clases)   & Base  & $0{,}639$ [$0{,}472$--$0{,}778$]
                        & $0{,}570$ [$0{,}371$--$0{,}751$]
                        & $0{,}697$ [$0{,}527$--$0{,}841$]
                        & $0{,}624$ [$0{,}457$--$0{,}760$]
                        & $0{,}296$ \\
L1 (4 clases)   & Large & $\mathbf{0{,}750}$ [$0{,}639$--$0{,}861$]
                        & $\mathbf{0{,}735}$ [$0{,}637$--$0{,}835$]
                        & $\mathbf{0{,}796}$ [$0{,}656$--$0{,}930$]
                        & $\mathbf{0{,}722}$ [$0{,}596$--$0{,}845$]
                        & $\mathbf{0{,}520}$ \\
L2 (38 clases)  & Base  & $0{,}222$ [$0{,}139$--$0{,}333$]
                        & $0{,}221$ [$0{,}147$--$0{,}294$]
                        & $0{,}707$ [$0{,}666$--$0{,}747$]
                        & $0{,}212$ [$0{,}092$--$0{,}302$]
                        & $0{,}170$ \\
L2 (38 clases)  & Large & $\mathbf{0{,}250}$ [$0{,}167$--$0{,}333$]
                        & $\mathbf{0{,}265}$ [$0{,}206$--$0{,}338$]
                        & $\mathbf{0{,}793}$ [$0{,}749$--$0{,}837$]
                        & $\mathbf{0{,}223}$ [$0{,}139$--$0{,}291$]
                        & $\mathbf{0{,}199}$ \\
L3 (224 clases) & Base  & $0{,}143$ [$0{,}107$--$0{,}179$]
                        & $0{,}132$ [$0{,}105$--$0{,}158$]
                        & $0{,}735$ [$0{,}693$--$0{,}779$]
                        & $0{,}107$ [$0{,}064$--$0{,}141$]
                        & $0{,}112$ \\
L3 (224 clases) & Large & $\mathbf{0{,}179}$ [$0{,}143$--$0{,}214$]
                        & $\mathbf{0{,}184}$ [$0{,}158$--$0{,}211$]
                        & $\mathbf{0{,}813}$ [$0{,}769$--$0{,}856$]
                        & $\mathbf{0{,}143}$ [$0{,}100$--$0{,}184$]
                        & $\mathbf{0{,}154}$ \\
\hline
\end{tabular}
\end{table}

PanDerm Large supera a PanDerm Base en las cinco métricas y en los
tres niveles. La mejora en AUROC es de $+9{,}9$\,pp en L1
($0{,}697 \to 0{,}796$), $+8{,}6$\,pp en L2
($0{,}707 \to 0{,}793$) y $+7{,}8$\,pp en L3
($0{,}735 \to 0{,}813$). La mejora en BAcc es de
$+16{,}5$\,pp en L1, $+4{,}4$\,pp en L2 y $+5{,}2$\,pp en
L3. La cifra L1 BAcc $= 0{,}735$ IC95\% [$0{,}637$--$0{,}835$]
se sitúa dentro del rango reportado para PanDerm Large bajo sondeo
lineal sobre los diez datasets de la
tabla~\ref{tab:lp-panderm} (BAcc entre $0{,}382$ y $0{,}792$),
lo que ubica DermapixelAI~1.0 como un \emph{benchmark} exigente
en castellano sin dependencia de marcas dermatoscópicas
estandarizadas. Frente al azar (BAcc esperada de $0{,}25 / 0{,}026 / 0{,}004$
para L1/L2/L3), PanDerm Large multiplica el rendimiento por
$2{,}9\times / 10{,}2\times / 40{,}9\times$.

\subsection*{Subestudio \emph{rosa\_verified}: imágenes verificadas
por la dermatóloga}

Se compara el sondeo lineal restringido a las imágenes verificadas
por la Dra.~R.~Taberner (\texttt{rosa\_verified=True}, $N = 919$,
train $= 744$, test $= 32$) frente al corpus completo
(\texttt{all\_verified}, $N = 1\,062$, train $= 874$, test $= 36$).
La tabla~\ref{tab:dermapixel-rosa} reporta los resultados sobre los
tres niveles ontológicos.

\begin{table}[H]
\centering
\footnotesize
\setlength{\tabcolsep}{4pt}
\renewcommand{\arraystretch}{0.95}
\caption{Subestudio \emph{rosa\_verified}: comparación entre el
subconjunto verificado caso a caso por la dermatóloga
($N_{\mathrm{train}} = 744$, $N_{\mathrm{test}} = 32$) y el corpus
completo ($N_{\mathrm{train}} = 874$, $N_{\mathrm{test}} = 36$).
Entre corchetes, IC95\% por \emph{bootstrap} estratificado.}
\label{tab:dermapixel-rosa}
\begin{tabular}{lllrrr}
\hline
Modelo & Subconjunto & Nivel & Acc & BAcc & W-F1 \\
\hline
PanDerm Large & \emph{rosa\_true\_only} & L1
  & $0{,}719$ [$0{,}594$--$0{,}844$]
  & $0{,}747$ [$0{,}646$--$0{,}865$]
  & $0{,}700$ [$0{,}550$--$0{,}841$] \\
PanDerm Large & \emph{all\_verified}    & L1
  & $0{,}750$ [$0{,}639$--$0{,}861$]
  & $0{,}735$ [$0{,}637$--$0{,}835$]
  & $0{,}722$ [$0{,}596$--$0{,}845$] \\
PanDerm Large & \emph{rosa\_true\_only} & L2
  & $0{,}250$ [$0{,}188$--$0{,}344$]
  & $0{,}250$ [$0{,}203$--$0{,}313$]
  & $0{,}227$ [$0{,}150$--$0{,}304$] \\
PanDerm Large & \emph{all\_verified}    & L2
  & $0{,}250$ [$0{,}167$--$0{,}333$]
  & $0{,}265$ [$0{,}206$--$0{,}338$]
  & $0{,}223$ [$0{,}139$--$0{,}291$] \\
PanDerm Large & \emph{rosa\_true\_only} & L3
  & $0{,}167$ [$0{,}167$--$0{,}167$]
  & $0{,}176$ [$0{,}176$--$0{,}176$]
  & $0{,}125$ [$0{,}117$--$0{,}139$] \\
PanDerm Large & \emph{all\_verified}    & L3
  & $0{,}179$ [$0{,}143$--$0{,}214$]
  & $0{,}184$ [$0{,}158$--$0{,}211$]
  & $0{,}143$ [$0{,}100$--$0{,}184$] \\
PanDerm Base  & \emph{rosa\_true\_only} & L1
  & $0{,}625$ [$0{,}469$--$0{,}781$]
  & $0{,}523$ [$0{,}326$--$0{,}735$]
  & $0{,}618$ [$0{,}451$--$0{,}773$] \\
PanDerm Base  & \emph{all\_verified}    & L1
  & $0{,}639$ [$0{,}472$--$0{,}778$]
  & $0{,}570$ [$0{,}371$--$0{,}751$]
  & $0{,}624$ [$0{,}457$--$0{,}760$] \\
\hline
\end{tabular}
\end{table}

Los intervalos de confianza al $95\%$ se solapan en todos los
pares de comparación, por lo que la diferencia entre evaluar sobre
el subconjunto verificado caso a caso y el corpus completo no es
estadísticamente significativa al tamaño muestral disponible. La
imputación etiológica derivada de la ontología
(\texttt{label\_source=ontology}) sobre las $170$ imágenes no
revisadas individualmente por la dermatóloga no introduce
degradación apreciable sobre las métricas reportadas, lo que
respalda operativamente el flujo de etiquetado descrito en el
Anexo~\ref{cap:anexoe}.

\subsection*{Cabezas alternativas sobre embeddings congelados}

Sobre los mismos \emph{embeddings} de PanDerm Large utilizados en la
tabla~\ref{tab:dermapixel-lp} se evalúan tres cabezas adicionales:
$k$-vecinos más cercanos con distancia coseno ($k = 10$, valor
óptimo dentro de $k \in \{1, 5, 10\}$), perceptrón multicapa con
una capa oculta de $512$ unidades y \emph{dropout} $0{,}3$ (MLP
512), y regresión logística con \texttt{class\_weight='balanced'}
(LogReg \emph{balanced}). El protocolo de partición es el del split
fijo de la tabla~\ref{tab:dermapixel-lp}
($N_{\mathrm{train}} = 874$, $N_{\mathrm{test}} = 36$ en L1/L2,
$N_{\mathrm{test}} = 28$ en L3). La tabla~\ref{tab:dermapixel-heads}
recoge la AUROC \emph{one-vs-rest}, métrica discriminativa robusta
al desbalance, como criterio comparativo.

\begin{table}[H]
\centering
\footnotesize
\setlength{\tabcolsep}{4pt}
\renewcommand{\arraystretch}{0.95}
\caption{Cabezas alternativas sobre \emph{embeddings} congelados de
PanDerm Large (test fijo de DermapixelAI~1.0). AUROC
\emph{one-vs-rest} por nivel ontológico. LP: regresión logística
estándar (referencia tabla~\ref{tab:dermapixel-lp}); kNN $k=10$:
$k$-vecinos con distancia coseno; MLP 512: perceptrón multicapa
con $512$ unidades ocultas y \emph{dropout} $0{,}3$; LogReg
\emph{balanced}: regresión logística con \texttt{class\_weight =
balanced}. Mejor valor por nivel en negrita.}
\label{tab:dermapixel-heads}
\begin{tabular}{lrrr}
\hline
Cabeza                & L1 (4 cls)         & L2 (38 cls)        & L3 (224 cls) \\
\hline
LP (referencia)       & $0{,}796$          & $0{,}793$          & $0{,}813$ \\
kNN $k=10$            & $0{,}752$          & $0{,}698$          & $0{,}694$ \\
MLP 512               & $\mathbf{0{,}817}$ & $\mathbf{0{,}812}$ & $\mathbf{0{,}827}$ \\
LogReg \emph{balanced}& $0{,}797$          & $0{,}800$          & $0{,}810$ \\
\hline
\end{tabular}
\end{table}

El MLP de $512$ unidades alcanza la mejor AUROC sobre los tres
niveles, con incrementos modestos pero consistentes sobre la cabeza
lineal de referencia: $+2{,}1$\,pp en L1
($0{,}796 \to 0{,}817$), $+1{,}9$\,pp en L2
($0{,}793 \to 0{,}812$) y $+1{,}4$\,pp en L3
($0{,}813 \to 0{,}827$). La regresión logística con re-ponderación
de clases mantiene Acc@1 idéntico al LP estándar y mejora la BAcc
sobre L1 ($0{,}703$ frente a $0{,}735$ del LP estándar dentro de
los IC95\% solapados), sin diferencias relevantes sobre L2 y L3.
La cabeza $k$-NN queda por debajo de las cabezas paramétricas en
los tres niveles, lo que se interpreta como insuficiente densidad
local en el espacio de \emph{embeddings} para $N_{\mathrm{train}}
= 874$ frente a un vocabulario de $224$ clases L3. La elección de
la regresión logística estándar como cabeza primaria en la
tabla~\ref{tab:dermapixel-lp} se mantiene defendible: la mejora
absoluta del MLP no compensa la pérdida de interpretabilidad
lineal de los coeficientes ni la mayor sensibilidad a la
inicialización aleatoria, especialmente con el tamaño muestral del
test.

\subsection*{Función de pérdida: Cross-Entropy frente a Focal}
\label{sec:dermapixel-focal}

La cola larga severa del corpus, particularmente acusada en el
nivel L3 ($77\%$ de las clases con $\leq 5$ ejemplos de
entrenamiento, $3{,}9$ muestras por clase de media en
\emph{train}, ver Anexo~\ref{cap:anexoc}), motiva contrastar la
pérdida de entropía cruzada estándar con \emph{Focal Loss} de
Lin et~al.\ 2017, diseñada explícitamente para mitigar el efecto
del desbalance reduciendo el peso de los ejemplos fáciles vía el
factor $(1-p_t)^\gamma$. Sobre los \emph{embeddings} congelados de
PanDerm Large del split fijo de la tabla~\ref{tab:dermapixel-lp}
se reentrena la cabeza lineal con \emph{Focal Loss} bajo dos
parámetros de focalización ($\gamma \in \{1{,}0,\,2{,}0,\,5{,}0\}$)
y dos regímenes de ponderación de clase
(\texttt{class\_weight = None} frente a
\texttt{class\_weight = balanced}), optimizada con LBFGS y
evaluada sobre el mismo conjunto de test que el sondeo lineal
estándar. La tabla~\ref{tab:dermapixel-focal} reporta la mejor
variante por nivel junto a las dos referencias previas (LP CE
estándar, LogReg \emph{balanced}) de la
tabla~\ref{tab:dermapixel-heads}.

\begin{table}[H]
\centering
\footnotesize
\setlength{\tabcolsep}{4pt}
\renewcommand{\arraystretch}{0.95}
\caption{Función de pérdida sobre \emph{embeddings} congelados de
PanDerm Large (test fijo de DermapixelAI~1.0,
$N_{\mathrm{test}} = 36$ en L1/L2 y $N_{\mathrm{test}} = 28$ en
L3). LP CE: regresión logística con entropía cruzada estándar
(referencia tabla~\ref{tab:dermapixel-lp}); LogReg
\emph{balanced}: re-ponderación inversa por frecuencia de clase
(tabla~\ref{tab:dermapixel-heads}); \emph{Focal} $\gamma$: pérdida
focal sin re-ponderación con parámetro de focalización indicado.
Mejor valor por columna y nivel en negrita.}
\label{tab:dermapixel-focal}
\begin{tabular}{llrrr}
\hline
Nivel        & Pérdida                                       & Acc@1              & BAcc               & AUROC \\
\hline
L1 (4 cls)   & LP CE (referencia)                            & $0{,}750$          & $0{,}735$          & $0{,}796$ \\
L1 (4 cls)   & LogReg \emph{balanced}                        & $0{,}694$          & $0{,}703$          & $0{,}797$ \\
L1 (4 cls)   & \emph{Focal} $\gamma = 2{,}0$                 & $\mathbf{0{,}778}$ & $\mathbf{0{,}768}$ & $\mathbf{0{,}800}$ \\
\hline
L2 (38 cls)  & LP CE (referencia)                            & $0{,}250$          & $0{,}265$          & $\mathbf{0{,}793}$ \\
L2 (38 cls)  & LogReg \emph{balanced}                        & $0{,}250$          & $0{,}265$          & $\mathbf{0{,}800}$ \\
L2 (38 cls)  & \emph{Focal} $\gamma = 2{,}0$                 & $\mathbf{0{,}278}$ & $\mathbf{0{,}324}$ & $0{,}778$ \\
L2 (38 cls)  & \emph{Focal} $\gamma = 5{,}0$                 & $\mathbf{0{,}278}$ & $\mathbf{0{,}324}$ & $\mathbf{0{,}793}$ \\
\hline
L3 (224 cls) & LP CE (referencia)                            & $\mathbf{0{,}143}$ & $\mathbf{0{,}132}$ & $\mathbf{0{,}813}$ \\
L3 (224 cls) & \emph{Focal} $\gamma = 2{,}0$                 & $\mathbf{0{,}143}$ & $\mathbf{0{,}132}$ & $0{,}771$ \\
\hline
\end{tabular}
\end{table}

Tres hallazgos cuantitativos estructuran la interpretación.
\textit{Primero}, \emph{Focal Loss} con $\gamma = 2{,}0$ y sin
re-ponderación de clase supera marginalmente a la entropía cruzada
estándar sobre L1 ($+2{,}8$\,pp Acc@1, $+3{,}3$\,pp BAcc,
$+0{,}4$\,pp AUROC) y sobre L2 ($+2{,}8$\,pp Acc@1,
$+5{,}9$\,pp BAcc), pero degrada la AUROC L2 en $-1{,}5$\,pp
($0{,}793 \to 0{,}778$). \textit{Segundo}, la combinación de
\emph{Focal Loss} con \texttt{class\_weight = balanced} no aporta
sobre la variante sin re-ponderación: el factor $(1-p_t)^\gamma$
modula ya el peso de los ejemplos según su dificultad y la
re-ponderación adicional no añade información discriminativa.
\textit{Tercero}, sobre L3 ningún ajuste de la función de pérdida
mejora a la entropía cruzada estándar: tanto Acc@1 como BAcc
permanecen idénticos ($0{,}143$ y $0{,}132$) y la AUROC
desciende $-4{,}2$\,pp ($0{,}813 \to 0{,}771$). Como en el caso
del FT denso (secci\'on~\ref{sec:dermapixel-ft-l1}), la cola larga
estructural del nivel L3 no admite mitigación mediante
modificación de la función de pérdida sobre el corpus actual.

El patrón observado es análogo al documentado en el aumento en
tiempo de inferencia (TTA, secci\'on~\ref{sec:dermapixel-tta}):
\emph{Focal Loss} desplaza el régimen Acc@1/BAcc-AUROC
favoreciendo la primera dimensión sobre la segunda, pero no
resuelve el cuello de botella estructural impuesto por la
cardinalidad y la cola larga del corpus. La elección del LP CE
estándar como cabeza primaria de la tabla~\ref{tab:dermapixel-lp}
se mantiene defendible: la mejora de \emph{Focal Loss} sobre L1 y
L2 está dentro de los IC95\% solapados y no compensa la pérdida
de interpretabilidad probabilística calibrada de la entropía
cruzada estándar.

\subsection*{Aumento en tiempo de inferencia}
\label{sec:dermapixel-tta}

La secci\'on~\ref{sec:ensemble-results} documenta la aplicación de
\emph{Test-Time Augmentation} (TTA) sobre el clasificador M1
(PanDerm Large \emph{fine-tuned}) del servidor de inferencia, con un
patrón de cinco augmentaciones deterministas (original, simetría
horizontal, simetría vertical, rotación de $90^\circ$ y rotación de
$270^\circ$, omitiendo la rotación de $180^\circ$). El presente
experimento replica el protocolo sobre la cabeza lineal entrenada
sobre los \emph{embeddings} de PanDerm Large del corpus
DermapixelAI~1.0: para cada imagen de test se extraen los cinco
\emph{embeddings} correspondientes a las cinco augmentaciones, cada
uno se evalúa con la cabeza lineal entrenada sobre el conjunto de
\emph{train} original y las cinco distribuciones de probabilidad se
promedian antes del \emph{argmax} final. La
tabla~\ref{tab:dermapixel-tta} compara el sondeo lineal sin TTA
(referencia tabla~\ref{tab:dermapixel-lp}) frente al sondeo lineal
con TTA sobre los tres niveles ontológicos.

\begin{table}[H]
\centering
\footnotesize
\setlength{\tabcolsep}{4pt}
\renewcommand{\arraystretch}{0.95}
\caption{Sondeo lineal con PanDerm Large sobre DermapixelAI~1.0:
referencia sin TTA frente a TTA con cinco augmentaciones
deterministas (original, simetría horizontal, simetría vertical,
rotación $90^\circ$, rotación $270^\circ$). Promedio de
probabilidades sobre las cinco augmentaciones antes del
\emph{argmax}. Mejor valor por par en negrita.}
\label{tab:dermapixel-tta}
\begin{tabular}{llrrrr}
\hline
Nivel & Modo & Acc@1 & Acc@3 & BAcc & AUROC \\
\hline
L1 (4 cls)   & LP                & $\mathbf{0{,}750}$ & $0{,}944$          & $\mathbf{0{,}735}$ & $\mathbf{0{,}796}$ \\
L1 (4 cls)   & LP $+$ TTA        & $0{,}694$          & $\mathbf{1{,}000}$ & $0{,}703$          & $0{,}779$ \\
L2 (38 cls)  & LP                & $\mathbf{0{,}250}$ & $0{,}278$          & $\mathbf{0{,}265}$ & $\mathbf{0{,}793}$ \\
L2 (38 cls)  & LP $+$ TTA        & $0{,}222$          & $\mathbf{0{,}500}$ & $0{,}250$          & $0{,}787$ \\
L3 (224 cls) & LP                & $\mathbf{0{,}179}$ & $0{,}214$          & $\mathbf{0{,}184}$ & $\mathbf{0{,}813}$ \\
L3 (224 cls) & LP $+$ TTA        & $0{,}179$          & $\mathbf{0{,}250}$ & $0{,}184$          & $0{,}807$ \\
\hline
\end{tabular}
\end{table}

El patrón es consistente sobre los tres niveles ontológicos: TTA
degrada Acc@1, BAcc y AUROC en magnitudes pequeñas ($-5{,}6$,
$-2{,}8$ y $0{,}0$\,pp Acc@1 sobre L1, L2 y L3 respectivamente;
$-3{,}2$, $-1{,}5$ y $0{,}0$\,pp BAcc; $-1{,}7$, $-0{,}6$ y
$-0{,}6$\,pp AUROC) y mejora sistemáticamente Acc@3
($+5{,}6$\,pp en L1 hasta alcanzar $1{,}000$, $+22{,}2$\,pp en
L2 y $+3{,}6$\,pp en L3). El comportamiento es coherente con la
observación reportada en secci\'on~\ref{sec:ensemble-results} para
M1 sobre HAM10000, donde TTA produce $-2{,}6$\,pp BAcc y
$+4{,}3$\,pp \emph{recall} de melanoma. La interpretación es que el
promediado sobre cinco augmentaciones suaviza la distribución de
probabilidad y reduce la confianza en la clase \emph{argmax} pero
distribuye masa de probabilidad entre clases visualmente cercanas,
lo que penaliza la precisión \emph{top-1} y mejora el conjunto
\emph{top-3}. PanDerm Base presenta la misma dirección de efecto
(cifras completas en el material reproducible). La recomendación
operativa derivada es que TTA es preferible en escenarios de cribado
con presentación \emph{top-k} de diagnósticos diferenciales al
clínico (Acc@3 L2 $= 0{,}500$ frente a $0{,}278$ sin TTA), y el
sondeo lineal sin TTA es preferible para decisiones automatizadas
\emph{top-1}.

\subsection*{Combinación de codificadores (\emph{ensemble})}
\label{sec:dermapixel-ensemble}

La secci\'on~\ref{sec:ensemble-results} documenta la combinación
de M1 (PanDerm Large \emph{fine-tuned}) y SigLIP-Large bajo sondeo
lineal sobre HAM10000 con orientación a maximizar el \emph{recall}
de melanoma. La presente subsección extiende esta línea sobre el
corpus DermapixelAI~1.0 y evalúa la combinación de los tres
codificadores fundacionales pre-entrenados disponibles en el
trabajo, operados como extractores congelados de
\emph{embeddings}: PanDerm Base ($768$\,d), PanDerm Large
($1\,024$\,d) y DermLIP v2 ($512$\,d). Sobre los
\emph{embeddings} ya cacheados de PanDerm Base y Large
(tabla~\ref{tab:dermapixel-lp}) y los \emph{embeddings} de
DermLIP v2 extraídos sobre las $1\,062$ imágenes evaluables del
corpus, se entrena una cabeza de regresión logística
independiente por codificador con el protocolo de
secci\'on~\ref{sec:lp} y se evalúan cinco estrategias de
combinación: \emph{(i)}~promedio de las tres distribuciones de
probabilidad (avg de 3), \emph{(ii)}~promedio restringido a
PanDerm Large y DermLIP v2 (avg L$+$D), \emph{(iii)}~máximo por
clase entre los tres codificadores (máx 3), \emph{(iv)}~apilado
(\emph{stacking}) con meta-cabeza de regresión logística sobre
las probabilidades concatenadas (stack 3) y \emph{(v)}~LP
individual por codificador como referencia. La
tabla~\ref{tab:dermapixel-ensemble} reporta los resultados sobre
los tres niveles ontológicos en el mismo conjunto de test fijo de
la tabla~\ref{tab:dermapixel-lp}.

\begin{table}[H]
\centering
\footnotesize
\setlength{\tabcolsep}{4pt}
\renewcommand{\arraystretch}{0.95}
\caption{Combinación de codificadores sobre DermapixelAI~1.0. LP
individual sobre PanDerm Base, PanDerm Large y DermLIP v2; tres
estrategias de combinación: promedio Large$+$DermLIP, máximo por
clase entre los tres, y apilado con meta-cabeza. Mejor valor por
columna y nivel en negrita.}
\label{tab:dermapixel-ensemble}
\begin{tabular}{llrrrr}
\hline
Nivel & Estrategia                         & Acc@1              & Acc@3              & BAcc               & AUROC \\
\hline
L1 (4 cls)   & LP PanDerm Base             & $0{,}639$          & ---                & $0{,}570$          & $0{,}697$ \\
L1 (4 cls)   & LP PanDerm Large            & $\mathbf{0{,}750}$ & $\mathbf{1{,}000}$ & $\mathbf{0{,}735}$ & $0{,}796$ \\
L1 (4 cls)   & LP DermLIP v2               & $0{,}694$          & $\mathbf{1{,}000}$ & $0{,}680$          & $0{,}789$ \\
L1 (4 cls)   & Avg L$+$D                   & $\mathbf{0{,}750}$ & $\mathbf{1{,}000}$ & $0{,}702$          & $\mathbf{0{,}817}$ \\
L1 (4 cls)   & Avg 3                       & $0{,}722$          & $\mathbf{1{,}000}$ & $0{,}694$          & $0{,}811$ \\
L1 (4 cls)   & Máx 3                       & $0{,}694$          & $\mathbf{1{,}000}$ & $0{,}652$          & $0{,}805$ \\
L1 (4 cls)   & Stack 3                     & $0{,}694$          & $\mathbf{1{,}000}$ & $0{,}703$          & $0{,}790$ \\
\hline
L2 (38 cls)  & LP PanDerm Base             & $0{,}222$          & ---                & $0{,}221$          & $0{,}707$ \\
L2 (38 cls)  & LP PanDerm Large            & $0{,}250$          & $0{,}500$          & $0{,}265$          & $0{,}793$ \\
L2 (38 cls)  & LP DermLIP v2               & $\mathbf{0{,}333}$ & $0{,}528$          & $\mathbf{0{,}279}$ & $\mathbf{0{,}860}$ \\
L2 (38 cls)  & Avg L$+$D                   & $0{,}278$          & $\mathbf{0{,}556}$ & $\mathbf{0{,}279}$ & $0{,}832$ \\
L2 (38 cls)  & Avg 3                       & $0{,}278$          & $0{,}528$          & $0{,}268$          & $0{,}821$ \\
L2 (38 cls)  & Máx 3                       & $0{,}222$          & $0{,}528$          & $0{,}235$          & $0{,}797$ \\
L2 (38 cls)  & Stack 3                     & $0{,}250$          & $0{,}528$          & $0{,}258$          & $0{,}804$ \\
\hline
L3 (224 cls) & LP PanDerm Base             & $0{,}143$          & ---                & $0{,}132$          & $0{,}735$ \\
L3 (224 cls) & LP PanDerm Large            & $\mathbf{0{,}179}$ & $\mathbf{0{,}286}$ & $\mathbf{0{,}184}$ & $\mathbf{0{,}813}$ \\
L3 (224 cls) & LP DermLIP v2               & $0{,}107$          & $0{,}214$          & $0{,}079$          & $0{,}809$ \\
L3 (224 cls) & Avg L$+$D                   & $\mathbf{0{,}179}$ & $0{,}250$          & $\mathbf{0{,}184}$ & $0{,}802$ \\
L3 (224 cls) & Avg 3                       & $0{,}143$          & $0{,}250$          & $0{,}152$          & $0{,}795$ \\
L3 (224 cls) & Máx 3                       & $0{,}143$          & $0{,}250$          & $0{,}132$          & $0{,}780$ \\
L3 (224 cls) & Stack 3                     & $0{,}143$          & $0{,}250$          & $0{,}158$          & $0{,}791$ \\
\hline
\end{tabular}
\end{table}

Tres hallazgos cuantitativos estructuran la interpretación.
\textit{Primero}, sobre L1 la combinación avg Large$+$DermLIP v2
mejora la AUROC en $+2{,}1$\,pp respecto a la mejor cabeza
individual (LP PanDerm Large, $0{,}796 \to 0{,}817$) sin mejorar
Acc@1 (se mantiene en $0{,}750$) ni BAcc ($0{,}735 \to
0{,}702$). El comportamiento es coherente con la literatura sobre
\emph{ensembles} probabilísticos: el promedio reduce la varianza
del ranking de probabilidades y mejora la métrica discriminativa
agregada (AUROC) sin trasladar la mejora a la decisión \emph{top-1}
en regímenes con tamaño muestral reducido
($N_{\mathrm{test}} = 36$).
\textit{Segundo}, sobre L2 DermLIP v2 individual lidera todas las
métricas: $+8{,}3$\,pp Acc@1 ($0{,}250 \to 0{,}333$),
$+1{,}4$\,pp BAcc ($0{,}265 \to 0{,}279$) y, especialmente,
$+6{,}7$\,pp AUROC ($0{,}793 \to 0{,}860$) respecto a LP
PanDerm Large. Ninguno de los \emph{ensembles} mejora el LP
individual de DermLIP v2 sobre L2 AUROC; el promedio avg
Large$+$DermLIP queda $-2{,}8$\,pp por debajo
($0{,}860 \to 0{,}832$). Este hallazgo es coherente con la
secci\'on~\ref{sec:dermapixel-zs}, donde DermLIP v2 lideraba
también la AUROC L2 en el experimento \emph{zero-shot}
($0{,}794$ en inglés frente a $0{,}655 / 0{,}697$ de BiomedCLIP
y SigLIP). La lectura conjunta sugiere que el preentrenamiento
contrastivo texto-imagen sobre el millón de pares de
Derm1M~\cite{derm1m} captura la granularidad de subcategoría
clínica L2 con mayor fidelidad que el preentrenamiento
auto-supervisado puro de la familia PanDerm.
\textit{Tercero}, sobre L3 ningún \emph{ensemble} mejora el LP
individual de PanDerm Large. La cola larga del nivel
($224$ clases efectivas, $3{,}9$ muestras por clase de media
en \emph{train}) limita la información complementaria que los
codificadores aportan entre sí: el codificador con mejor
representación lineal de la cola larga (PanDerm Large, AUROC
$0{,}813$) domina el comportamiento del \emph{ensemble} sin que
las cabezas adicionales aporten señal incremental.

La recomendación operativa derivada de la
tabla~\ref{tab:dermapixel-ensemble} depende del nivel ontológico:
sobre L1 conviene el \emph{ensemble} avg Large$+$DermLIP por su
ventaja en AUROC; sobre L2 conviene la cabeza individual de
DermLIP v2; sobre L3 conviene la cabeza individual de PanDerm
Large. La jerarquía PanDerm Large $>$ PanDerm Base $>$ DermLIP
v2 que sostiene la mayor parte de las cifras de
secci\'on~\ref{sec:rep-panderm} y de la
tabla~\ref{tab:dermapixel-lp} no es por tanto uniforme entre
niveles: el preentrenamiento contrastivo texto-imagen invierte
la jerarquía sobre la subcategoría clínica L2 de un corpus
independiente al preentrenamiento.

\subsection*{Ajuste fino parcial del codificador sobre los tres niveles}
\label{sec:dermapixel-ft}
\label{sec:dermapixel-ft-l1}

La secci\'on~\ref{sec:ft-results} cuantifica el efecto del ajuste
fino supervisado (FT) sobre PanDerm Base aplicado a cuatro datasets
externos contrastados, con incrementos de hasta $+47{,}1$\,pp BAcc
sobre HAM10000 y $+19{,}4$\,pp sobre PAD-UFES-20. La presente
subsección replica el protocolo sobre DermapixelAI~1.0 al nivel de
categoría etiológica L1 ($4$ clases: patología inflamatoria,
infecciosa, tumoral y genodermatosis). El experimento se restringe
al nivel L1 por una razón muestral: la cola larga de L3 con
$3{,}9$ muestras/clase de media en \emph{train} no soporta FT denso
sin saturar el régimen de sobreajuste, mientras que L1 admite la
prueba de concepto al menor coste experimental. El encoder es
PanDerm Large con las dos últimas capas descongeladas
($25{,}2$\,M parámetros entrenables, equivalente al $8{,}3\%$
del total del encoder), seguido de una cabeza
\emph{fully-connected} $1\,024 \to 4$. La receta de optimización
es AdamW con tasa de aprendizaje diferenciada (cabeza $10^{-3}$,
encoder $10^{-5}$), \emph{weight decay} $10^{-4}$, \emph{batch
size} $16$, programación coseno con calentamiento de una época y
$10$ épocas totales. Las augmentaciones de entrenamiento son
\emph{RandomResizedCrop}, simetría horizontal, simetría vertical
con probabilidad $0{,}3$ y \emph{ColorJitter}. La pérdida es
entropía cruzada con \texttt{class\_weight='balanced'}, requerida
por el fuerte desbalance de la clase \emph{Genodermatosis}
($8 / 874$ imágenes de \emph{train}). La selección del
\emph{checkpoint} se realiza por BAcc máxima sobre validación; el
tiempo total de entrenamiento sobre NVIDIA DGX Spark GB10 es de
$2{,}4$ minutos.

\begin{table}[H]
\centering
\footnotesize
\setlength{\tabcolsep}{4pt}
\renewcommand{\arraystretch}{0.95}
\caption{Curva de entrenamiento del ajuste fino parcial sobre
DermapixelAI~1.0 nivel L1 ($10$ épocas, PanDerm Large con dos
últimas capas descongeladas, $25{,}2$\,M parámetros
entrenables). La fila en negrita marca el \emph{best val
checkpoint} seleccionado para la evaluación final.}
\label{tab:dermapixel-ft-train}
\begin{tabular}{rrrrr}
\hline
Época & Loss \emph{train} & BAcc val. & BAcc test & AUROC test \\
\hline
$1$  & $1{,}32$ & $0{,}574$ & $0{,}590$ & $0{,}681$ \\
$2$  & $0{,}99$ & $0{,}606$ & $0{,}625$ & $0{,}722$ \\
$3$  & $0{,}87$ & $0{,}588$ & $0{,}654$ & $0{,}717$ \\
$\mathbf{4}$  & $\mathbf{0{,}79}$ & $\mathbf{0{,}630}$ & $\mathbf{0{,}622}$ & $\mathbf{0{,}715}$ \\
$5$  & $0{,}74$ & $0{,}592$ & $0{,}638$ & $0{,}709$ \\
$6$  & $0{,}69$ & $0{,}630$ & $0{,}640$ & $0{,}715$ \\
$7$  & $0{,}64$ & $0{,}610$ & $0{,}608$ & $0{,}719$ \\
$8$--$10$ & $0{,}60$--$0{,}63$ & $0{,}60$--$0{,}62$ & $0{,}59$--$0{,}62$ & $\sim 0{,}72$ \\
\hline
\end{tabular}
\end{table}

\begin{table}[H]
\centering
\footnotesize
\setlength{\tabcolsep}{4pt}
\renewcommand{\arraystretch}{0.95}
\caption{Métricas finales del ajuste fino parcial sobre
DermapixelAI~1.0 nivel L1 evaluadas sobre el conjunto de test
($N_{\mathrm{test}} = 36$) con el \emph{best val checkpoint}
(época $4$). Entre corchetes, IC95\% por \emph{bootstrap}
estratificado ($1\,000$ remuestreos).}
\label{tab:dermapixel-ft-final}
\begin{tabular}{lrr}
\hline
Métrica   & Valor      & IC95\% \emph{bootstrap} \\
\hline
Acc@1     & $0{,}583$  & [$0{,}444$--$0{,}722$] \\
Acc@3     & $0{,}972$  & [$0{,}917$--$1{,}000$] \\
BAcc      & $0{,}622$  & [$0{,}510$--$0{,}735$] \\
AUROC     & $0{,}715$  & [$0{,}603$--$0{,}830$] \\
W-F1      & $0{,}568$  & [$0{,}421$--$0{,}700$] \\
Kappa     & $0{,}254$  & [$0{,}027$--$0{,}484$] \\
\hline
\end{tabular}
\end{table}

Tres hallazgos cuantitativos estructuran la interpretación.
\textit{Primero}, el ajuste fino parcial mejora la BAcc en
$+18{,}2$\,pp respecto a la cifra de validación cruzada agrupada
por caso del sondeo lineal congelado sobre el mismo corpus
(BAcc $0{,}440 \pm 0{,}058$ en la
tabla~\ref{tab:dermapixel-cv} frente a $0{,}622$ con FT). El
IC95\% del FT ([$0{,}510$--$0{,}735$]) se sitúa claramente por
encima del intervalo correspondiente al sondeo lineal bajo
validación cruzada, lo que confirma que el incremento no es
atribuible a la realización muestral del split fijo y que el
ajuste fino añade valor real más allá del ruido de partición.
\textit{Segundo}, el FT no supera a la combinación lineal de
codificadores pre-entrenados \texttt{avg\_large\_dermlip}
documentada en la secci\'on~\ref{sec:dermapixel-ensemble}: BAcc
$0{,}622$ frente a $0{,}702$ del \emph{ensemble} y AUROC
$0{,}715$ frente a $0{,}817$. La interpretación es que, para
$N_{\mathrm{train}} = 874$, una combinación lineal de dos
codificadores fundacionales congelados rinde competitivamente con
el FT parcial denso, lo que sugiere un protocolo híbrido en
producción basado en \emph{ensemble} de codificadores
pre-entrenados con cabezas ligeras antes que descongelar el
encoder. \textit{Tercero}, la curva de entrenamiento
(tabla~\ref{tab:dermapixel-ft-train}) muestra que la pérdida de
\emph{train} desciende monótonamente de $1{,}32$ a $0{,}60$
durante las diez épocas, mientras que la BAcc de validación y de
test satura a partir de la época $4$ ($0{,}630$ y $0{,}622$
respectivamente). El \emph{best val checkpoint} no es el último,
patrón sintomático de sobreajuste suave compatible con el régimen
de $N_{\mathrm{train}}$ reducido y $25$\,M parámetros
entrenables. Sin ampliación del corpus o regularización más
fuerte, el FT parcial alcanza un techo bajo sobre L1.

La replicación del mismo protocolo sobre los niveles L2 y L3 cierra
el caso del FT denso parcial sobre el corpus. El encoder, las dos
últimas capas descongeladas, la receta de optimización y el
criterio de selección por mejor BAcc de validación son idénticos a
los de L1; varía únicamente la cabeza supervisada ($1\,024 \to 37$
para L2, $1\,024 \to 224$ para L3) y el coste computacional sobre
NVIDIA DGX Spark GB10 (suma de los tres niveles inferior a $5$
minutos). La tabla~\ref{tab:dermapixel-ft-multilevel} compara las
métricas finales sobre los tres niveles bajo el mismo protocolo.

\begin{table}[H]
\centering
\footnotesize
\setlength{\tabcolsep}{4pt}
\renewcommand{\arraystretch}{0.95}
\caption{Ajuste fino parcial del codificador (PanDerm Large, dos
últimas capas descongeladas) sobre los tres niveles ontológicos de
DermapixelAI~1.0. Test fijo $N_{\mathrm{test}} = 36$ en L1/L2 y
$N_{\mathrm{test}} = 28$ en L3. Entre corchetes, IC95\% por
\emph{bootstrap} estratificado ($1\,000$ remuestreos). Mejor valor
por columna en negrita.}
\label{tab:dermapixel-ft-multilevel}
\begin{tabular}{lrrrr}
\hline
Nivel        & Acc@1                            & Acc@3                            & BAcc                             & AUROC \\
\hline
L1 (4 cls)   & $\mathbf{0{,}583}$ [$0{,}444$--$0{,}722$] & $\mathbf{0{,}972}$ [$0{,}917$--$1{,}000$] & $\mathbf{0{,}622}$ [$0{,}510$--$0{,}735$] & $0{,}715$ [$0{,}603$--$0{,}830$] \\
L2 (37 cls)  & $0{,}250$ [$0{,}167$--$0{,}333$] & $0{,}500$ [$0{,}389$--$0{,}611$] & $0{,}324$ [$0{,}265$--$0{,}397$] & $\mathbf{0{,}852}$ [$0{,}814$--$0{,}888$] \\
L3 (224 cls) & $0{,}179$ [$0{,}143$--$0{,}214$] & $0{,}286$ [$0{,}214$--$0{,}357$] & $0{,}184$ [$0{,}158$--$0{,}211$] & $0{,}804$ [$0{,}773$--$0{,}835$] \\
\hline
\end{tabular}
\end{table}

Dos hallazgos cuantitativos estructuran la lectura de los tres
niveles. El primero se refiere a L2 y se documenta de forma
explícita en la tabla~\ref{tab:dermapixel-ft-vs-lora}: el FT
parcial denso sobre L2 con $25{,}2$\,M parámetros entrenables
($8{,}3\%$ del encoder) alcanza BAcc $0{,}324$ y AUROC
$0{,}852$, cifras que el protocolo SpanDerm v0 con adaptación LoRA
(secci\'on~\ref{sec:dermapixel-spanderm-v0}) supera con
$524\,K$ parámetros entrenables ($0{,}17\%$ del encoder, $48$
veces menos parámetros): BAcc $0{,}363 \pm 0{,}007$
($+4{,}0$\,pp) y AUROC $0{,}855 \pm 0{,}006$ (empate dentro de
la desviación estándar entre semillas).

\begin{table}[H]
\centering
\footnotesize
\setlength{\tabcolsep}{4pt}
\renewcommand{\arraystretch}{0.95}
\caption{Comparativa sobre el nivel L2 de DermapixelAI~1.0 entre el
sondeo lineal congelado (referencia), el ajuste fino parcial denso
con dos últimas capas descongeladas y el protocolo SpanDerm v0 con
adaptación LoRA. Mismo conjunto de test ($N_{\mathrm{test}} = 36$,
$37$--$38$ clases según consolidación de variantes ortográficas
de queratinización descrita en
secci\'on~\ref{sec:dermapixel-spanderm-v0}). Mejor valor por
columna en negrita.}
\label{tab:dermapixel-ft-vs-lora}
\begin{tabular}{lrrr}
\hline
Protocolo                       & BAcc L2                         & AUROC L2                       & Params entren. \\
\hline
LP test fijo (referencia, tabla~\ref{tab:dermapixel-lp}) & $0{,}265$                       & $0{,}793$                      & $39\,K$ \\
FT parcial denso (este trabajo) & $0{,}324$                       & $0{,}852$                      & $25{,}2$\,M \\
\textbf{SpanDerm v0 (LoRA)}     & $\mathbf{0{,}363 \pm 0{,}007}$  & $\mathbf{0{,}855 \pm 0{,}006}$ & $\mathbf{524\,K}$ \\
\hline
\end{tabular}
\end{table}

La lectura cuantitativa de la tabla~\ref{tab:dermapixel-ft-vs-lora}
refuerza la decisión arquitectónica de SpanDerm v0 documentada en
secci\'on~\ref{sec:dermapixel-spanderm-v0}: la adaptación de bajo
rango sobre un encoder fundacional congelado supera al ajuste fino
denso parcial sobre el mismo encoder bajo un régimen de
$N_{\mathrm{train}} = 874$ y cardinalidad L2 moderada. El segundo
hallazgo se refiere a L3: el FT denso parcial alcanza BAcc
$0{,}184$ y AUROC $0{,}804$ sobre las $224$ clases efectivas,
cifras que no superan al sondeo lineal congelado sobre el mismo
\emph{split} fijo (BAcc $0{,}184$ idéntica, AUROC $0{,}813$,
$-0{,}9$\,pp). La cola larga severa del nivel L3 ($3{,}9$
muestras por clase de media en \emph{train}, $77\%$ de las
clases con $\leq 5$ ejemplos) anula el beneficio del ajuste fino
denso: la varianza muestral por clase domina el régimen de
entrenamiento, lo que reserva el nivel L3 al protocolo de ranking
por prototipos coseno SpanDerm v0
(tabla~\ref{tab:dermapixel-spanderm-l3}) o, alternativamente, a la
ampliación del corpus vía banco hospitalario.

La utilidad práctica del ajuste fino parcial sobre
DermapixelAI~1.0 queda por tanto condicionada a tres factores. El
primero es la ampliación del corpus prevista vía banco
hospitalario del HUSLL, recogida como técnica pendiente en la
secci\'on~\ref{sec:dermapixel-futuro}. El segundo es la aplicación
de técnicas de FT más sofisticadas (adaptación LoRA, regularización
con \emph{label smoothing}, \emph{Mixup} y \emph{CutMix}) ya
planificadas en el diseño de SpanDerm v0 referido en la
secci\'on~\ref{sec:linea-spanderm} y ejecutadas con resultado
favorable en la secci\'on~\ref{sec:dermapixel-spanderm-v0}. El
tercero es la naturaleza estructural de la cola larga L3, que
limita en términos absolutos el rendimiento alcanzable bajo
cualquier protocolo de ajuste denso sobre el corpus actual.

\subsection*{Validación cruzada agrupada por caso}

El conjunto de test fijo $N = 36$ utilizado en la
tabla~\ref{tab:dermapixel-lp} corresponde a un único \emph{split}
publicado del corpus. Para estimar el rendimiento esperado del
sondeo lineal congelado con menor dependencia de la realización
muestral concreta, se aplica una validación cruzada de cinco
\emph{folds} agrupada por \texttt{case\_id}
(\emph{StratifiedGroupKFold} estratificado por L1 sobre el corpus
completo de $1\,062$ imágenes, garantizando que ninguna imagen del
mismo caso clínico aparezca simultáneamente en \emph{train} y
\emph{test}). Cada \emph{fold} retiene aproximadamente
$N \approx 210$ imágenes con representación de las cuatro
categorías etiológicas L1 y de la cola larga L2/L3. La
tabla~\ref{tab:dermapixel-cv} reporta la media y la desviación
estándar sobre los cinco \emph{folds} para PanDerm Large bajo
sondeo lineal.

\begin{table}[H]
\centering
\footnotesize
\setlength{\tabcolsep}{4pt}
\renewcommand{\arraystretch}{0.95}
\caption{Validación cruzada agrupada por caso (cinco \emph{folds},
estratificado por L1 sobre \texttt{case\_id}) con PanDerm Large
sobre DermapixelAI~1.0. Cifras: media $\pm$ desviación estándar
sobre los cinco \emph{folds}. La cardinalidad L3 reportada aquí
(250 clases) corresponde a las clases efectivas en el corpus
completo; las tablas anteriores del split fijo
(tabla~\ref{tab:dermapixel-lp}, tabla~\ref{tab:dermapixel-heads},
tabla~\ref{tab:dermapixel-tta}, tabla~\ref{tab:dermapixel-ensemble},
tabla~\ref{tab:dermapixel-ft-final}) usan 224 clases L3, las
únicas representadas en el split \texttt{train}.}
\label{tab:dermapixel-cv}
\begin{tabular}{lrrrrrr}
\hline
Nivel           & Acc@1                 & Acc@3                 & BAcc                  & AUROC                 & W-F1                  & Kappa \\
\hline
L1 (4 cls)      & $0{,}574 \pm 0{,}029$ & $0{,}994 \pm 0{,}004$ & $0{,}440 \pm 0{,}058$ & $0{,}753 \pm 0{,}054$ & $0{,}569 \pm 0{,}025$ & $0{,}314 \pm 0{,}035$ \\
L2 (38 cls)     & $0{,}227 \pm 0{,}015$ & $0{,}418 \pm 0{,}037$ & $0{,}168 \pm 0{,}020$ & $0{,}749 \pm 0{,}019$ & $0{,}221 \pm 0{,}022$ & $0{,}186 \pm 0{,}015$ \\
L3 (250 cls)    & $0{,}124 \pm 0{,}007$ & $0{,}232 \pm 0{,}028$ & $0{,}108 \pm 0{,}014$ & $0{,}800 \pm 0{,}018$ & $0{,}124 \pm 0{,}005$ & $0{,}110 \pm 0{,}005$ \\
\hline
\end{tabular}
\end{table}

La validación cruzada agrupada por caso reduce la cifra puntual
del split fijo y proporciona la estimación robusta del rendimiento
esperado de un sondeo lineal congelado sobre DermapixelAI~1.0. La
sobrestimación observada en el test fijo respecto a la media de
los cinco \emph{folds} es de $+29{,}5$\,pp BAcc en L1,
$+9{,}7$\,pp en L2 y $+7{,}6$\,pp en L3, lo que se atribuye a la
distribución favorable del split publicado (test sin representación
de la categoría etiológica \emph{Genodermatosis} y baja densidad de
la cola larga L3). La AUROC, en cambio, se mantiene estable entre
las dos estimaciones: $0{,}796 \to 0{,}753$ en L1
($-4{,}3$\,pp), $0{,}793 \to 0{,}749$ en L2
($-4{,}4$\,pp) y $0{,}813 \to 0{,}800$ en L3
($-1{,}3$\,pp), magnitudes muy inferiores a la caída observada
sobre BAcc. La cifra defendible como estimación del rendimiento
esperado del \emph{transfer learning} congelado sobre el corpus en
castellano es la de la tabla~\ref{tab:dermapixel-cv}; las cifras
del test fijo en la tabla~\ref{tab:dermapixel-lp} se mantienen
como referencia del punto operativo del split publicado y como
base de la comparación con el experimento \emph{zero-shot} de la
secci\'on~\ref{sec:dermapixel-zs}, que utiliza el mismo conjunto
de test para garantizar la trazabilidad cruzada. El experimento
de combinación de codificadores
(secci\'on~\ref{sec:dermapixel-ensemble}) confirma además que
DermLIP v2 supera a PanDerm Large sobre L2 sin requerir
entrenamiento adicional ($+6{,}7$\,pp AUROC), lo que matiza la
jerarquía PanDerm Large $>$ PanDerm Base $>$ DermLIP v2 que se
sostiene en los demás niveles ontológicos y en la mayoría de los
datasets externos contrastados: la elección del codificador
óptimo depende del nivel ontológico evaluado, con
\emph{ensemble} avg Large$+$DermLIP sobre L1, DermLIP v2
individual sobre L2 y PanDerm Large individual sobre L3. El
ajuste fino parcial del codificador
(secci\'on~\ref{sec:dermapixel-ft-l1}) recupera
$+18{,}2$\,pp BAcc sobre la cifra de validación cruzada del
sondeo lineal congelado en el nivel L1, confirmando que el FT
añade valor real más allá del ruido del \emph{split} fijo. La
combinación de codificadores pre-entrenados (\emph{ensemble}
\texttt{avg\_large\_dermlip}, secci\'on~\ref{sec:dermapixel-ensemble})
supera no obstante al FT parcial en L1 (BAcc $0{,}702$ vs
$0{,}622$, AUROC $0{,}817$ vs $0{,}715$), lo que sugiere que
para corpus de tamaño moderado el protocolo óptimo combina
codificadores múltiples sin descongelar el encoder.

\subsection*{Caveats}

Cinco \emph{caveats} acompañan la interpretación de las cifras
reportadas en esta sección. \textit{Primero}, el tamaño del conjunto
de test ($N = 36$ en L1/L2, $N = 28$ en L3) es reducido, lo que
ensancha los intervalos de confianza al $95\%$ y exige cautela
sobre las cifras puntuales. \textit{Segundo}, la categoría
etiológica \emph{Genodermatosis} no aparece en el conjunto de test
de L1 ($0 / 36$); la AUROC \emph{one-vs-rest} se promedia sobre
las tres clases con soporte y las métricas Acc/BAcc reflejan
efectivamente una clasificación de tres clases.
\textit{Tercero}, $26$ clases de L3 visibles en train no
aparecen en test, por lo que el clasificador se entrena sobre ellas
pero su discriminación no es evaluable; el tamaño efectivo del test
L3 desciende a $N = 28$. \textit{Cuarto}, la cola larga de L3
($224$ clases sobre $874$ imágenes de entrenamiento,
$3{,}9$ muestras/clase de media) limita el rendimiento absoluto
alcanzable bajo sondeo lineal y reserva el régimen de ajuste fino
o de cabezas jerárquicas como línea futura. \textit{Quinto}, la
AUROC \emph{one-vs-rest} resulta sustancialmente más robusta al
cambio entre split fijo y validación cruzada agrupada que la
accuracy y la BAcc: la caída entre la tabla~\ref{tab:dermapixel-lp}
y la tabla~\ref{tab:dermapixel-cv} es de $-4{,}3$\,pp en L1,
$-4{,}4$\,pp en L2 y $-1{,}3$\,pp en L3 para AUROC, frente a
caídas de $-29{,}5$\,pp, $-9{,}7$\,pp y $-7{,}6$\,pp para BAcc
sobre los mismos niveles. La AUROC se interpreta por tanto como
métrica primaria de comparación entre configuraciones de
\emph{transfer learning} congelado sobre corpus de tamaño moderado.

\section{Validación \emph{zero-shot} multimodal sobre DermapixelAI 1.0}
\label{sec:dermapixel-zs}

La sección~\ref{sec:dermapixel-results} reporta el rendimiento de
PanDerm Base y Large bajo sondeo lineal sobre el corpus
DermapixelAI~1.0. La sección~\ref{sec:zs-results} caracteriza el
comportamiento \emph{zero-shot} multimodal de la familia DermLIP
sobre HAM10000 y la sensibilidad documentada al tokenizador y a la
formulación del \emph{prompt}. La presente sección completa ambas
líneas con un experimento orientado a cuantificar el rendimiento
\emph{zero-shot} sobre el mismo conjunto de test de
DermapixelAI~1.0, sin entrenamiento de ninguna cabeza
clasificadora, y a contrastar la penalización inducida por la
formulación del \emph{prompt} en castellano frente a la
formulación equivalente en inglés. El experimento opera sobre el
mismo conjunto de test que la sección~\ref{sec:dermapixel-results}
($N_{\mathrm{test}} = 36$ en L1/L2, $N_{\mathrm{test}} = 28$ en L3
tras filtrar clases ausentes en train), garantizando la
trazabilidad cruzada con el sondeo lineal y conservando la garantía
de no solapamiento con Derm1M reportada en
secci\'on~\ref{sec:dermapixel-results} y
secci\'on~\ref{sec:auditoria-resultados} ($0 / 507\,657$
coincidencias por hash MD5).

Se evalúan tres modelos vision-lenguaje contrastivos:
\emph{(i)}~DermLIP v2 (encoder visual PanDerm-base con texto
PubMedBERT-256, preentrenado sobre Derm1M);
\emph{(ii)}~SigLIP-SO400M (Google, ViT-SO400M-14-384 sobre WebLI,
generalista); y \emph{(iii)}~BiomedCLIP (Microsoft, ViT-B/16 con
PubMedBERT-256 sobre PubMed-15M). La predicción se obtiene como
\emph{argmax} sobre la similitud coseno entre la representación
visual y un \emph{ensemble} de tres plantillas textuales por nivel
ontológico, en castellano y en inglés
(\textit{p.~ej.}, ``una fotografía clínica de \{c\}'', ``una
imagen dermatológica de \{c\}'', ``una lesión cutánea de tipo
\{c\}''). Los \emph{embeddings} de las tres plantillas se promedian
y re-normalizan antes de la predicción. El nivel L3 se evalúa
únicamente en castellano, dado que el vocabulario ontológico no
dispone de traducción inglesa validada para sus $224$ clases. Los
intervalos de confianza al $95\%$ se calculan por
\emph{bootstrap} estratificado sobre $1\,000$ remuestreos. La
tabla~\ref{tab:dermapixel-zs} reporta los resultados.

\begin{table}[H]
\centering
\footnotesize
\setlength{\tabcolsep}{4pt}
\renewcommand{\arraystretch}{0.95}
\caption{Clasificación \emph{zero-shot} multimodal sobre el
conjunto de test de DermapixelAI~1.0 ($N_{\mathrm{test}} = 36$ en
L1/L2, $N_{\mathrm{test}} = 28$ en L3). Entre corchetes, IC95\%
por \emph{bootstrap} estratificado ($1\,000$ remuestreos) para
las dos métricas primarias (Acc@1 y AUROC). Las métricas BAcc y
Acc@3 se reportan como valor puntual por concisión; los IC95\%
correspondientes están disponibles en el material reproducible.
L3 se evalúa únicamente en castellano. Mejor valor por columna,
nivel e idioma en negrita.}
\label{tab:dermapixel-zs}
\begin{tabular}{lllrrrr}
\hline
Nivel & Modelo & Lang & Acc@1 & Acc@3 & BAcc & AUROC \\
\hline
L1 (4 cls)   & BiomedCLIP    & es & $\mathbf{0{,}389}$ [$0{,}250$--$0{,}528$]
                                  & $\mathbf{0{,}944}$
                                  & $0{,}375$
                                  & $0{,}679$ [$0{,}562$--$0{,}781$] \\
L1 (4 cls)   & BiomedCLIP    & en & $0{,}389$ [$0{,}250$--$0{,}528$]
                                  & $\mathbf{0{,}972}$
                                  & $0{,}324$
                                  & $0{,}639$ [$0{,}524$--$0{,}738$] \\
L1 (4 cls)   & SigLIP-SO400M & es & $0{,}222$ [$0{,}111$--$0{,}361$]
                                  & $0{,}917$
                                  & $0{,}178$
                                  & $0{,}650$ [$0{,}571$--$0{,}735$] \\
L1 (4 cls)   & SigLIP-SO400M & en & $0{,}333$ [$0{,}194$--$0{,}472$]
                                  & $0{,}889$
                                  & $0{,}292$
                                  & $0{,}663$ [$0{,}595$--$0{,}719$] \\
L1 (4 cls)   & DermLIP v2    & es & $0{,}361$ [$0{,}222$--$0{,}500$]
                                  & $0{,}861$
                                  & $0{,}379$
                                  & $0{,}699$ [$0{,}600$--$0{,}794$] \\
L1 (4 cls)   & DermLIP v2    & en & $\mathbf{0{,}500}$ [$0{,}333$--$0{,}667$]
                                  & $0{,}917$
                                  & $\mathbf{0{,}441}$
                                  & $\mathbf{0{,}709}$ [$0{,}565$--$0{,}841$] \\
\hline
L2 (38 cls)  & BiomedCLIP    & es & $0{,}111$ [$0{,}111$--$0{,}111$]
                                  & $0{,}278$
                                  & $0{,}118$
                                  & $0{,}647$ [$0{,}604$--$0{,}688$] \\
L2 (38 cls)  & BiomedCLIP    & en & $0{,}083$ [$0{,}056$--$0{,}111$]
                                  & $0{,}250$
                                  & $0{,}088$
                                  & $0{,}655$ [$0{,}601$--$0{,}713$] \\
L2 (38 cls)  & SigLIP-SO400M & es & $0{,}139$ [$0{,}139$--$0{,}139$]
                                  & $0{,}250$
                                  & $0{,}235$
                                  & $0{,}636$ [$0{,}589$--$0{,}680$] \\
L2 (38 cls)  & SigLIP-SO400M & en & $0{,}194$ [$0{,}111$--$0{,}278$]
                                  & $0{,}278$
                                  & $0{,}221$
                                  & $0{,}697$ [$0{,}640$--$0{,}747$] \\
L2 (38 cls)  & DermLIP v2    & es & $\mathbf{0{,}222}$ [$0{,}167$--$0{,}278$]
                                  & $0{,}333$
                                  & $0{,}265$
                                  & $0{,}762$ [$0{,}721$--$0{,}805$] \\
L2 (38 cls)  & DermLIP v2    & en & $\mathbf{0{,}250}$ [$0{,}194$--$0{,}306$]
                                  & $\mathbf{0{,}417}$
                                  & $\mathbf{0{,}294}$
                                  & $\mathbf{0{,}794}$ [$0{,}747$--$0{,}837$] \\
\hline
L3 (224 cls) & BiomedCLIP    & es & $0{,}107$ [$0{,}036$--$0{,}179$]
                                  & $0{,}214$
                                  & $0{,}079$
                                  & $0{,}703$ [$0{,}673$--$0{,}732$] \\
L3 (224 cls) & SigLIP-SO400M & es & $0{,}071$ [$0{,}000$--$0{,}143$]
                                  & $0{,}214$
                                  & $0{,}053$
                                  & $0{,}696$ [$0{,}667$--$0{,}723$] \\
L3 (224 cls) & DermLIP v2    & es & $\mathbf{0{,}107}$ [$0{,}071$--$0{,}143$]
                                  & $\mathbf{0{,}214}$
                                  & $\mathbf{0{,}079}$
                                  & $\mathbf{0{,}740}$ [$0{,}719$--$0{,}763$] \\
\hline
\end{tabular}
\end{table}

Tres hallazgos cuantitativos estructuran la interpretación.

\textit{Primero}, DermLIP v2 obtiene el mejor o equivalente
rendimiento en las cinco columnas métricas sobre los tres niveles
ontológicos, con excepción de Acc@3 L1 (BiomedCLIP en castellano
e inglés alcanza $0{,}944$ y $0{,}972$ respectivamente, frente a
$0{,}861$ y $0{,}917$ de DermLIP v2). La superioridad
sistemática es coherente con su preentrenamiento sobre Derm1M, el
único corpus de los tres específicamente dermatológico (un millón
de pares imagen-texto curados sobre la familia
PanDerm)~\cite{derm1m}. Sobre L2, DermLIP v2 alcanza AUROC
$0{,}794$ en inglés y $0{,}762$ en castellano, magnitudes
comparables a la AUROC reportada en
secci\'on~\ref{sec:rep-panderm} para PanDerm Large bajo sondeo
lineal sobre conjuntos clínicos heterogéneos (\textit{p.~ej.},
$0{,}746$ sobre MSKCC).

\textit{Segundo}, la formulación del \emph{prompt} en castellano
penaliza el rendimiento de los modelos preentrenados sobre corpus
mayoritariamente anglófonos. Sobre L1 Acc@1, DermLIP v2 cae de
$0{,}500$ (inglés) a $0{,}361$ (castellano), una pérdida de
$13{,}9$\,pp; SigLIP-SO400M cae de $0{,}333$ a $0{,}222$
($-11{,}1$\,pp). BiomedCLIP, en cambio, mantiene
$0{,}389$ en ambos idiomas (\emph{gap} nulo). La interpretación
es coherente con la composición de los corpus textuales de
preentrenamiento: PubMedBERT incorpora literatura biomédica en
inglés con presencia anecdótica de otras lenguas, mientras que
WebLI (SigLIP) y Derm1M presentan sesgo aún más marcado hacia
contenidos anglófonos. La penalización cuantificada motiva la
construcción de modelos vision-lenguaje contrastivos con rama
textual en castellano, línea futura abierta por el corpus
DermapixelAI~1.0.

\textit{Tercero}, la comparación con el sondeo lineal sobre el
mismo conjunto de test (tabla~\ref{tab:dermapixel-lp}) cuantifica
la brecha entre un sistema entrenado y un sistema
\emph{zero-shot}. Sobre L1, PanDerm Large bajo LP alcanza
BAcc $= 0{,}735$ frente a $0{,}441$ del mejor \emph{zero-shot}
(DermLIP v2 en inglés), una diferencia de $+29{,}4$\,pp a favor
del sondeo lineal. Sobre L3, la brecha es de $+10{,}5$\,pp BAcc
($0{,}184$ vs $0{,}079$) y $+7{,}3$\,pp AUROC ($0{,}813$ vs
$0{,}740$). El comportamiento se invierte parcialmente sobre L2:
LP PanDerm Large obtiene AUROC $0{,}793$ y \emph{zero-shot}
DermLIP v2 en inglés alcanza $0{,}794$, un empate efectivo dentro
del margen de los intervalos de confianza. La interpretación es
relevante: las \emph{features} pre-entrenadas de la familia
PanDerm/DermLIP son ya competitivas a nivel de subcategoría
clínica L2 sin necesidad de cabeza entrenada, mientras que la
discriminación etiológica de grano grueso (L1) y de cola larga
diagnóstica (L3) sí requiere un clasificador supervisado.
Adicionalmente, Acc@3 sobre L1 oscila entre $0{,}861$ (DermLIP v2
en castellano) y $0{,}972$ (BiomedCLIP en inglés), magnitudes
operativamente relevantes para escenarios de cribado en los que
el sistema presente las tres categorías etiológicas más probables
al clínico antes del juicio diagnóstico.

\subsection*{Caveats}

Cuatro \emph{caveats} acompañan la interpretación de las cifras
reportadas. \textit{Primero}, el tamaño del conjunto de test
($N = 36$ en L1/L2, $N = 28$ en L3) produce intervalos de
confianza al $95\%$ amplios, lo que exige cautela sobre las
diferencias puntuales entre configuraciones cercanas (\textit{p.~ej.},
el empate L2 AUROC entre LP PanDerm y \emph{zero-shot} DermLIP v2
no descarta diferencias del orden de unos pocos puntos
porcentuales en una réplica con mayor tamaño muestral).
\textit{Segundo}, el \emph{ensemble} de tres plantillas textuales
por nivel constituye una formulación mínima de \emph{prompt
engineering}; la literatura sobre \emph{zero-shot} multimodal
emplea típicamente entre diez y ochenta plantillas generadas con
asistencia de LLMs, lo que sugiere que las cifras reportadas
constituyen una cota inferior del rendimiento alcanzable.
\textit{Tercero}, la evaluación L3 únicamente en castellano impide
cuantificar la penalización por idioma sobre la cola larga
diagnóstica, dado que el vocabulario ontológico no dispone de
traducción inglesa validada para sus $224$ clases.
\textit{Cuarto}, la auditoría sistemática de no solapamiento por
hash MD5 documentada en secci\'on~\ref{sec:auditoria-resultados}
y reiterada en secci\'on~\ref{sec:dermapixel-results} garantiza la
independencia del corpus de test frente al preentrenamiento de
DermLIP v2 sobre Derm1M, condición que aporta validez a la
comparación entre modelos sin caveats de fuga de información.

\section{Cobertura experimental sobre DermapixelAI~1.0}
\label{sec:dermapixel-cobertura}

Las secciones precedentes evalúan los codificadores sobre
DermapixelAI~1.0 como extractores de \emph{features} congelados con
cabezas paramétricas ligeras o clasificación contrastiva
\emph{zero-shot}, siguiendo el protocolo estándar de la literatura
sobre modelos fundacionales en imagen médica, que prioriza la
cuantificación limpia del valor del preentrenamiento sobre datos del
dominio sin contaminación por hiperparámetros específicos de ajuste
fino. La síntesis de las técnicas de \emph{transfer learning}
ejecutadas y pendientes sobre el corpus, junto con la estimación
prospectiva de un clasificador desplegable en castellano, se recoge
en la secci\'on~\ref{sec:dermapixel-futuro} del
capítulo~\ref{cap:conclusiones}.

\section{Ajuste fino con LoRA: SpanDerm v0}
\label{sec:dermapixel-spanderm-v0}

La secci\'on~\ref{sec:dermapixel-ft-l1} documenta el ajuste fino
parcial denso sobre L1, que mejora la BAcc en $+18{,}2$\,pp sobre
la cifra de validación cruzada del sondeo lineal congelado pero
no supera al \emph{ensemble} multi-codificador
(secci\'on~\ref{sec:dermapixel-ensemble}). La
secci\'on~\ref{sec:dermapixel-futuro} identifica la adaptación de
bajo rango (\emph{Low-Rank Adaptation}, LoRA)~\cite{lora} como
línea pendiente de prioridad alta, en línea con el diseño SpanDerm
v0 documentado en la secci\'on~\ref{sec:linea-spanderm}. La
presente sección reporta la ejecución de SpanDerm v0 sobre el
nivel L2 ($38$ clases efectivas tras la consolidación de
variantes ortográficas de queratinización ya descrita en la
secci\'on~\ref{sec:dermapixel-ft-l1}) y el ranking L3 por
prototipos sobre las $250$ clases efectivas del corpus. La
elección de L2 como objetivo de la cabeza supervisada responde a
la cardinalidad operativa del nivel (rendimiento clínicamente
informativo sin alcanzar el régimen de cola larga extrema de L3,
con $3{,}5$ muestras/clase en $N_{\mathrm{train}} = 874$). El
ranking L3 sigue la estrategia prevista en el diseño v0: prototipos
calculados como la media de las representaciones del conjunto de
entrenamiento por clase L3 y predicción por similitud coseno.

\begin{table}[H]
\centering
\footnotesize
\setlength{\tabcolsep}{4pt}
\renewcommand{\arraystretch}{0.95}
\caption{Configuración del protocolo SpanDerm v0 sobre PanDerm
Large con adaptación LoRA. Los $524\,288$ parámetros
entrenables corresponden a $39\,K$ en la cabeza FC
$1024 \to 38$ y $485\,K$ en los ocho módulos LoRA insertados en
las capas lineales de los dos últimos bloques de transformer.}
\label{tab:dermapixel-spanderm-config}
\begin{tabular}{ll}
\hline
Componente                  & Configuración \\
\hline
Encoder base                & PanDerm Large ($1\,024$\,dim, $303{,}85$\,M params, congelado) \\
LoRA                        & $r=16$, $\alpha=32$, \emph{dropout} $0{,}1$ \\
Capas adaptadas             & \texttt{blocks.22} y \texttt{blocks.23}: \texttt{attn.qkv}, \texttt{attn.proj}, \texttt{mlp.fc1}, \texttt{mlp.fc2} (8 capas) \\
Cabeza supervisada          & FC $1024 \to 38$ (L2 con queratinización consolidada) \\
Parámetros entrenables      & $524\,288$ ($0{,}17\%$ del total) \\
Optimizador                 & AdamW, $lr_{\mathrm{head}} = 10^{-3}$, $lr_{\mathrm{LoRA}} = 5 \times 10^{-4}$, $\mathrm{wd} = 10^{-4}$ \\
\emph{Scheduler}            & Cosine con \emph{warmup} de una época \\
Pérdida                     & \emph{Cross-entropy} con \texttt{class\_weight=balanced} \\
\emph{Batch size}           & $16$ \\
Épocas                      & $15$ \\
Augmentaciones \emph{train} & RandomResizedCrop($0{,}7$--$1{,}0$), hflip, vflip ($p=0{,}3$), ColorJitter ($0{,}1$) \\
Selección                   & Mejor BAcc de validación; reporte adicional de la media de las últimas cinco épocas \\
Semillas                    & $\{42, 43, 44\}$ \\
\hline
\end{tabular}
\end{table}

Las cifras reportadas a continuación corresponden a la media y
la desviación estándar sobre las tres semillas
($\{42, 43, 44\}$). La tabla~\ref{tab:dermapixel-spanderm-l2}
reporta las dos estimaciones complementarias sobre el conjunto de
test L2 ($N_{\mathrm{test}} = 36$): la cifra del \emph{best val
checkpoint} y la media de las últimas cinco épocas sin selección
de \emph{checkpoint}. La segunda estimación caracteriza el
comportamiento asintótico del entrenamiento sin la varianza
asociada al criterio de selección.

\begin{table}[H]
\centering
\footnotesize
\setlength{\tabcolsep}{4pt}
\renewcommand{\arraystretch}{0.95}
\caption{Métricas de SpanDerm v0 sobre el conjunto de test L2 de
DermapixelAI~1.0 ($N_{\mathrm{test}} = 36$, $38$ clases con
queratinización consolidada). Cifras: media $\pm$ desviación
estándar sobre tres semillas ($\{42, 43, 44\}$). La columna
\emph{best val} aplica selección por máxima BAcc de validación;
la columna \emph{últimas $5$ ép.} promedia las cinco épocas
finales sin selección. Mejor por fila en negrita.}
\label{tab:dermapixel-spanderm-l2}
\begin{tabular}{lrr}
\hline
Métrica   & \emph{Best val \emph{checkpoint}} & Media últimas $5$ ép. \\
\hline
Acc@1     & $\mathbf{0{,}250 \pm 0{,}000}$ & $0{,}209 \pm 0{,}005$ \\
Acc@3     & $\mathbf{0{,}500 \pm 0{,}023}$ & $0{,}457 \pm 0{,}011$ \\
BAcc      & $\mathbf{0{,}363 \pm 0{,}007}$ & $0{,}309 \pm 0{,}006$ \\
AUROC     & $\mathbf{0{,}855 \pm 0{,}006}$ & $0{,}837 \pm 0{,}008$ \\
W-F1      & $\mathbf{0{,}218 \pm 0{,}013}$ & --- \\
Kappa     & $\mathbf{0{,}220 \pm 0{,}001}$ & --- \\
\hline
\end{tabular}
\end{table}

La estabilidad entre semillas es alta sobre todas las métricas
reportadas (desviación estándar $\leq 0{,}013$ en valor absoluto
y $\leq 0{,}008$ sobre AUROC y BAcc). La diferencia entre el
\emph{best val checkpoint} y la media de las últimas cinco épocas
($+5{,}4$\,pp BAcc, $+1{,}8$\,pp AUROC, $+4{,}1$\,pp Acc@1)
indica que la curva de entrenamiento exhibe un patrón de
sobreajuste suave análogo al observado en el FT parcial denso de
la secci\'on~\ref{sec:dermapixel-ft-l1}, pero atenuado por la
restricción al $0{,}17\%$ de parámetros entrenables que impone
la adaptación LoRA.

La tabla~\ref{tab:dermapixel-spanderm-comp} sitúa las cifras de
SpanDerm v0 frente a las dos referencias internas sobre el mismo
conjunto de test y el mismo nivel ontológico: el sondeo lineal
congelado bajo validación cruzada agrupada por caso
(tabla~\ref{tab:dermapixel-cv}), el sondeo lineal sobre el split
fijo (tabla~\ref{tab:dermapixel-lp}) y el codificador DermLIP v2
individual reportado en la secci\'on~\ref{sec:dermapixel-ensemble}
como mejor opción sobre L2 sin entrenamiento de cabeza.

\begin{table}[H]
\centering
\footnotesize
\setlength{\tabcolsep}{4pt}
\renewcommand{\arraystretch}{0.95}
\caption{Comparativa sobre el nivel L2 de DermapixelAI~1.0
($38$ clases con queratinización consolidada,
$N_{\mathrm{test}} = 36$) entre SpanDerm v0 y las referencias
internas del trabajo. La fila LP CV se considera la cifra honesta
del sondeo lineal congelado por su independencia del \emph{split}
fijo (secci\'on~\ref{sec:dermapixel-results}). Mejor valor por
columna en negrita.}
\label{tab:dermapixel-spanderm-comp}
\begin{tabular}{lllrr}
\hline
Método                                              & Encoder                       & Params entren.    & BAcc L2                         & AUROC L2 \\
\hline
LP test fijo (tabla~\ref{tab:dermapixel-lp})        & PanDerm Large $1\,024$ (cong.) & $39\,K$           & $0{,}265$                       & $0{,}793$ \\
LP $5$-fold CV (tabla~\ref{tab:dermapixel-cv})      & PanDerm Large $1\,024$ (cong.) & $39\,K$           & $0{,}168 \pm 0{,}020$           & $0{,}749 \pm 0{,}019$ \\
DermLIP v2 individual (secci\'on~\ref{sec:dermapixel-ensemble}) & DermLIP $512$ (cong.)        & $19{,}5\,K$       & $0{,}279$                       & $\mathbf{0{,}860}$ \\
\textbf{SpanDerm v0} (este trabajo)                 & PanDerm Large $1\,024$ $+$ LoRA & $524\,K$          & $\mathbf{0{,}363 \pm 0{,}007}$  & $0{,}855 \pm 0{,}006$ \\
\hline
\end{tabular}
\end{table}

La comparativa documenta dos hallazgos cuantitativos sobre el
mismo conjunto de test. \textit{Primero}, SpanDerm v0 mejora la
BAcc L2 en $+19{,}5$\,pp respecto a la cifra de validación
cruzada agrupada por caso del sondeo lineal congelado
($0{,}168 \to 0{,}363$) y en $+8{,}4$\,pp respecto a DermLIP
v2 individual ($0{,}279 \to 0{,}363$), siendo la mejor BAcc L2
documentada sobre el corpus. \textit{Segundo}, la AUROC L2 de
SpanDerm v0 ($0{,}855 \pm 0{,}006$) empata estadísticamente con
la de DermLIP v2 individual ($0{,}860$) dentro del margen de la
desviación estándar entre semillas, mientras que mejora en
$+10{,}6$\,pp la cifra de validación cruzada del sondeo lineal
congelado. El compromiso entre BAcc y AUROC entre SpanDerm v0 y
DermLIP v2 individual se interpreta como complementariedad: la
cabeza supervisada con LoRA sobre PanDerm Large alcanza mayor
discriminación de las clases minoritarias (mejor BAcc), mientras
que el preentrenamiento contrastivo PubMedBERT-DermLIP retiene
una ligera ventaja sobre el ranking global de la AUROC
\emph{one-vs-rest}.

El ranking L3 por prototipos sobre las $250$ clases del
vocabulario de entrenamiento reporta las cifras de la
tabla~\ref{tab:dermapixel-spanderm-l3}. Los prototipos se calculan
como la media de las representaciones del conjunto de
entrenamiento por clase L3, y la predicción se obtiene por
similitud coseno entre la representación de la imagen de test y
los $250$ prototipos. Esta estrategia corresponde a la
formulación documentada en la secci\'on~\ref{sec:linea-spanderm}
para mitigar la cola larga del nivel L3 sin requerir una cabeza
supervisada directa, inviable bajo $3{,}5$ muestras/clase de
media en $N_{\mathrm{train}} = 874$.

\begin{table}[H]
\centering
\footnotesize
\setlength{\tabcolsep}{4pt}
\renewcommand{\arraystretch}{0.95}
\caption{Ranking L3 por prototipos coseno de SpanDerm v0 sobre el
conjunto de test de DermapixelAI~1.0 ($250$ clases efectivas
del vocabulario de entrenamiento). Cifras: media $\pm$ desviación
estándar sobre tres semillas. La columna $\times$~azar reporta el
factor de mejora respecto a la asignación aleatoria uniforme
sobre $250$ clases ($1/250 = 0{,}004$).}
\label{tab:dermapixel-spanderm-l3}
\begin{tabular}{lrrr}
\hline
Métrica & Media                 & Std            & $\times$~azar \\
\hline
Acc@1   & $0{,}167$             & $0{,}000$      & $\times 42$ \\
Acc@3   & $0{,}259$             & $0{,}013$      & $\times 22$ \\
Acc@5   & $0{,}296$             & $0{,}026$      & $\times 15$ \\
BAcc    & $0{,}148$             & $0{,}000$      & --- \\
\hline
\end{tabular}
\end{table}

El ranking L3 por prototipos funciona pero queda limitado por la
cola larga del corpus: Acc@1 $= 0{,}167$ y Acc@5 $= 0{,}296$
mejoran la asignación aleatoria en órdenes de magnitud de
$\times 42$ y $\times 15$ respectivamente, pero permanecen lejos
del régimen operativo clínico necesario para despliegue
asistencial sobre el vocabulario completo. La cifra de BAcc
($0{,}148$) refleja la penalización adicional por las clases
ausentes en validación: $178$ clases L3 con muestras efectivas
en el rango $1$--$5$ imágenes en el \emph{split} de entrenamiento
(cola larga estricta del vocabulario representado), lo que limita la calidad del prototipo
correspondiente al promedio de muy pocas muestras y produce
prototipos ruidosos para gran parte del vocabulario.

Tres hallazgos cuantitativos cierran la sección. \textit{Primero},
SpanDerm v0 con sólo $0{,}17\%$ de parámetros entrenables iguala
o supera a las tres referencias internas sobre BAcc L2
($0{,}363 \pm 0{,}007$ frente a $0{,}168 \pm 0{,}020$ del LP
CV, $0{,}265$ del LP fijo y $0{,}279$ de DermLIP v2 individual),
lo que valida la línea de adaptación de bajo rango como protocolo
preferente sobre el corpus en castellano y materializa la
implicación de diseño anticipada en la
secci\'on~\ref{sec:linea-spanderm}. \textit{Segundo}, la AUROC L2
de SpanDerm v0 ($0{,}855 \pm 0{,}006$) empata estadísticamente
con la de DermLIP v2 nativo ($0{,}860$): la adaptación textual
ligera de PanDerm Large mediante LoRA alcanza el mismo techo
discriminativo que el preentrenamiento contrastivo
PubMedBERT-DermLIP sobre Derm1M, lo que sugiere que la elección
del codificador y la estrategia de adaptación son ortogonales
sobre la cardinalidad L2 del corpus. \textit{Tercero}, el ranking
L3 por prototipos coseno es funcional ($\times 42$ sobre azar en
Acc@1 respecto al baseline uniforme de $1/250 \approx 0{,}004$,
con factores análogos de $\times 22$ y $\times 15$ sobre Acc@3 y
Acc@5 respecto a sus baselines uniformes correspondientes) pero
está acotado por la cola larga: $178$ de las $250$ clases L3
efectivas en entrenamiento contienen $\leq 5$ imágenes. La progresión hacia SpanDerm v1
(cabeza L3 directa) queda condicionada a la incorporación del
banco hospitalario del HUSLL, con umbral de activación de
$\geq 30$ imágenes/clase L3 ya documentado en la
secci\'on~\ref{sec:dermapixel-futuro} y en la
secci\'on~\ref{sec:linea-spanderm}.

\section{Interpretabilidad del SAE: alineamiento con conceptos
dermatoscópicos clásicos}
\label{sec:sae-derm7pt}

La línea de interpretabilidad mediante Sparse Autoencoders (SAE)
descrita en el Anexo~\ref{cap:anexof} introduce un SAE Large de
$16\,384$ \emph{features} entrenado sobre las activaciones de
PanDerm Large y valida la base conceptual del modelo sobre los
$48$ conceptos visuales de SkinCon~\cite{skincon}. La presente
sección extiende esa validación al corpus
Derm7pt~\cite{kawahara2019} ($1\,011$ lesiones dermatoscópicas) y
explota por primera vez en este trabajo las anotaciones del
\emph{Seven-Point Checklist} (siete criterios diagnósticos
estructurados por imagen). El diseño contempla tres experimentos
ejecutados sobre DGX Spark GB10 el 2026-06-02: E1 alineamiento
\emph{single-feature} entre las $16\,384$ activaciones del SAE y
los siete criterios binarizados, E2 \emph{Concept Bottleneck Model}
(CBM) jerárquico con cuello de botella sobre los siete conceptos
para predicción melanoma/no-melanoma, y E3 cruce con los $16$
conceptos dermatoscópicos del Grupo A propuestos a la
Dra.~R.~Taberner para anotación sobre DermapixelAI~1.0. Los tres
experimentos comparten el corpus completo Derm7pt
($N = 1\,011$, $24{,}5\%$ melanomas), con \emph{splits} oficiales
\emph{train}~$=413$, \emph{val}~$=203$, \emph{test}~$=395$
($100$ melanomas en test) cuando aplica división. La activación
del SAE sobre el corpus muestra $3\,214$ \emph{features} activas en
más del $5\%$ de las imágenes, $2\,978$ \emph{features}
inactivas (\emph{dead}) y $2\,717$ \emph{features} muy activas
($>50\%$), distribución coherente con la sparsidad nominal del
SAE.

El experimento E1 evalúa, para cada uno de los siete criterios
binarizados como \emph{present}/\emph{absent}, la AUROC efectiva
$\max(\mathrm{AUROC}, 1-\mathrm{AUROC})$ de cada una de las
$16\,384$ \emph{features} individuales como predictor binario
del criterio. La tabla~\ref{tab:sae-d7pt-e1} reporta el AUROC
efectivo de la mejor \emph{feature} por criterio
(\emph{top-1}) y el promedio sobre las diez mejores
(\emph{top-10}). Todas las cifras se calculan sobre el corpus
completo Derm7pt ($N = 1\,011$).

\begin{table}[H]
\centering
\footnotesize
\setlength{\tabcolsep}{4pt}
\renewcommand{\arraystretch}{0.95}
\caption{Alineamiento \emph{single-feature} entre las
$16\,384$ \emph{features} del SAE Large y los siete criterios
del \emph{Seven-Point Checklist} sobre Derm7pt
($N = 1\,011$). AUROC efectiva
$\max(\mathrm{AUROC}, 1-\mathrm{AUROC})$ de la mejor
\emph{feature} individual (\emph{top-1}) y promedio sobre las
diez mejores (\emph{top-10}). Mejor cifra agregada en negrita.}
\label{tab:sae-d7pt-e1}
\begin{tabular}{lrrrr}
\hline
Criterio                 & $N_{+}$ & $N_{-}$ & AUROC \emph{top-1} & AUROC \emph{top-10} \\
\hline
pigment\_network         & $611$ & $400$ & $0{,}754$ & $0{,}737$ \\
streaks                  & $358$ & $653$ & $0{,}688$ & $0{,}684$ \\
pigmentation             & $423$ & $588$ & $0{,}700$ & $0{,}657$ \\
regression\_structures   & $253$ & $758$ & $0{,}687$ & $0{,}673$ \\
dots\_and\_globules      & $782$ & $229$ & $0{,}681$ & $0{,}663$ \\
blue\_whitish\_veil      & $195$ & $816$ & $0{,}746$ & $0{,}721$ \\
vascular\_structures     & $188$ & $823$ & $\mathbf{0{,}823}$ & $\mathbf{0{,}782}$ \\
\hline
\end{tabular}
\end{table}

Las siete \emph{features} más alineadas individualmente alcanzan
AUROC efectivo en el intervalo $[0{,}681,\,0{,}823]$, con el
máximo sobre \emph{vascular structures} ($0{,}823$) y mínimo sobre
\emph{dots and globules} ($0{,}681$). La estabilidad entre el
\emph{top-1} y el \emph{top-10} (diferencia $\leq 0{,}043$ en
valor absoluto) indica que el alineamiento no se concentra en una
única \emph{feature} aislada sino que el SAE codifica los
conceptos mediante un pequeño subconjunto coherente de
activaciones.

El experimento E2 construye un CBM jerárquico de dos etapas:
una primera cabeza de regresión logística sobre las
$16\,384$ activaciones del SAE estima la probabilidad de cada
uno de los siete criterios; un meta-clasificador logístico toma
estas siete estimaciones como vector de entrada y predice
melanoma/no-melanoma. La comparativa con el sondeo lineal directo
sobre las $16\,384$ \emph{features} sin cuello de botella se
reporta en la tabla~\ref{tab:sae-d7pt-cbm}. El AUROC binario por
concepto de la cabeza intermedia se reporta en la
tabla~\ref{tab:sae-d7pt-cbm-coefs} junto con el peso aprendido
por el meta-clasificador sobre cada concepto.

\begin{table}[H]
\centering
\footnotesize
\setlength{\tabcolsep}{4pt}
\renewcommand{\arraystretch}{0.95}
\caption{Comparativa entre sondeo lineal directo sobre las
$16\,384$ \emph{features} del SAE Large y CBM jerárquico con
los siete criterios \emph{Seven-Point Checklist} como cuello de
botella. Conjunto de test Derm7pt
($N_{\mathrm{test}} = 395$, $100$ melanomas). AUROC en
negrita.}
\label{tab:sae-d7pt-cbm}
\begin{tabular}{lrrrrr}
\hline
Arquitectura                                & Acc@1     & BAcc      & AUROC              & W-F1      & Kappa     \\
\hline
Sondeo lineal directo (SAE $\to$ melanoma)  & $0{,}841$ & $0{,}735$ & $\mathbf{0{,}890}$ & $0{,}829$ & $0{,}527$ \\
CBM (SAE $\to$ $7$ conceptos $\to$ melanoma) & $0{,}825$ & $0{,}681$ & $\mathbf{0{,}832}$ & $0{,}801$ & $0{,}440$ \\
\hline
\end{tabular}
\end{table}

\begin{table}[H]
\centering
\footnotesize
\setlength{\tabcolsep}{4pt}
\renewcommand{\arraystretch}{0.95}
\caption{Cabeza intermedia del CBM: AUROC binario sobre el conjunto
de test Derm7pt ($N_{\mathrm{test}} = 395$) por concepto del
\emph{Seven-Point Checklist}, y peso aprendido por el
meta-clasificador logístico para la predicción binaria
melanoma/no-melanoma. Pesos ordenados de mayor a menor
importancia relativa.}
\label{tab:sae-d7pt-cbm-coefs}
\begin{tabular}{lrr}
\hline
Concepto                & AUROC concepto & Peso meta-LP \\
\hline
blue\_whitish\_veil     & $0{,}895$ & $+2{,}54$ \\
regression\_structures  & $0{,}818$ & $+2{,}07$ \\
streaks                 & $0{,}829$ & $+1{,}78$ \\
vascular\_structures    & $0{,}906$ & $+1{,}65$ \\
pigmentation            & $0{,}792$ & $+1{,}44$ \\
dots\_and\_globules     & $0{,}723$ & $+1{,}41$ \\
pigment\_network        & $0{,}901$ & $+1{,}04$ \\
\hline
\end{tabular}
\end{table}

La AUROC binaria de la cabeza intermedia se sitúa en el intervalo
$[0{,}723,\,0{,}906]$ sobre los siete conceptos, con
\emph{vascular structures} ($0{,}906$), \emph{pigment network}
($0{,}901$) y \emph{blue whitish veil} ($0{,}895$) como las
predicciones conceptuales más fiables. El paso de las activaciones
crudas del SAE a través del cuello de botella conceptual reduce
la AUROC binaria melanoma de $0{,}890$ a $0{,}832$
($-5{,}8$\,pp), pérdida cuantificada que documenta el compromiso
entre interpretabilidad y precisión bruta. Los pesos del
meta-clasificador identifican \emph{blue whitish veil} y
\emph{regression structures} como los conceptos con mayor peso
relativo para la decisión final melanoma/no-melanoma
($+2{,}54$ y $+2{,}07$ respectivamente), seguidos de
\emph{streaks} y \emph{vascular structures}. Esta jerarquía
coincide cualitativamente con el peso clínico asignado a estos
criterios en el \emph{Seven-Point Checklist} original
descrito por Kawahara et~al.~\cite{kawahara2019}, donde el velo
azul-blanquecino y las estructuras de regresión figuran entre los
criterios mayores ($2$ puntos cada uno) de la escala
diagnóstica.

El experimento E3 cruza los $16$ conceptos dermatoscópicos del
Grupo~A propuestos a la Dra.~R.~Taberner para anotación sobre las
$49$ imágenes dermatoscópicas de DermapixelAI~1.0
(secci\'on~\ref{sec:dermapixel}) con los siete criterios del
\emph{Seven-Point Checklist} sobre Derm7pt. Para cada concepto
Rosa con \emph{mapping} a un criterio Derm7pt (directo o por
subcategoría) se reporta la AUROC efectiva de la mejor
\emph{feature} SAE sobre la binarización correspondiente. El
resultado se sintetiza en la tabla~\ref{tab:sae-d7pt-cross-rosa}.

\begin{table}[H]
\centering
\footnotesize
\setlength{\tabcolsep}{4pt}
\renewcommand{\arraystretch}{0.95}
\caption{Cruce de los $16$ conceptos dermatoscópicos del
Grupo~A propuestos a la Dra.~R.~Taberner con los siete criterios del
\emph{Seven-Point Checklist} sobre Derm7pt ($N = 1\,011$).
Para los conceptos con \emph{mapping} se reporta la AUROC efectiva
de la mejor \emph{feature} SAE Large sobre la binarización
correspondiente. Conceptos cubiertos (AUROC $\geq 0{,}70$) en
negrita.}
\label{tab:sae-d7pt-cross-rosa}
\begin{tabular}{rlllrl}
\hline
\# & Concepto Rosa (Grupo~A)                & \emph{Mapping} Derm7pt       & $N_{+}$ & AUROC eff           & Cobertura \\
\hline
$1$  & Red dermatoscópica regular           & pigment\_network=typical            & $381$ & $\mathbf{0{,}748}$  & $\checkmark$ \\
$2$  & Red dermatoscópica atípica           & pigment\_network=atypical           & $230$ & $\mathbf{0{,}732}$  & $\checkmark$ \\
$3$  & Glóbulos pigmentarios                & dots\_and\_globules                 & $782$ & $0{,}681$           & --- \\
$4$  & Puntos pigmentarios                  & dots\_and\_globules                 & $782$ & $0{,}681$           & --- \\
$5$  & Estrías radiales                     & streaks                             & $358$ & $0{,}688$           & --- \\
$6$  & Pseudopodios                         & streaks=irregular                   & $251$ & $\mathbf{0{,}729}$  & $\checkmark$ \\
$7$  & Velo azul-blanquecino                & blue\_whitish\_veil=present         & $195$ & $\mathbf{0{,}746}$  & $\checkmark$ \\
$8$  & Estructuras de regresión             & regression\_structures $\neq$ absent& $253$ & $0{,}687$           & --- \\
$9$  & Patrón empedrado                     & ---                                  & ---   & ---                 & --- \\
$10$ & Patrón paralelo de surcos            & ---                                  & ---   & ---                 & --- \\
$11$ & Patrón paralelo de crestas           & ---                                  & ---   & ---                 & --- \\
$12$ & Lagunas vasculares                   & ---                                  & ---   & ---                 & --- \\
$13$ & Vasos polimorfos                     & vascular\_structures (agregado)     & $140$ & $\mathbf{0{,}810}$  & $\checkmark$ \\
$14$ & Vasos en horquilla                   & vascular\_structures=hairpin        & $15$  & $\mathbf{0{,}926}$  & $\checkmark$ \\
$15$ & Vasos en corona                      & ---                                  & ---   & ---                 & --- \\
$16$ & Patrón homogéneo desestructurado     & pigmentation (\emph{diffuse})       & $380$ & $\mathbf{0{,}704}$  & $\checkmark$ \\
\hline
\end{tabular}
\end{table}

De los $16$ conceptos del Grupo~A, $11$ admiten \emph{mapping}
operativo a un criterio o subcategoría del \emph{Seven-Point
Checklist} (directo en tres casos, por subcategoría en ocho), y
los $5$ restantes (patrón empedrado, patrón paralelo de surcos,
patrón paralelo de crestas, lagunas vasculares, vasos en corona)
corresponden a estructuras dermatoscópicas no contempladas en la
escala original. Sobre los $11$ conceptos mapeados, $7$ quedan
cubiertos por al menos una \emph{feature} SAE con AUROC efectivo
$\geq 0{,}70$. El máximo absoluto sobre vasos en horquilla
(AUROC $= 0{,}926$ con sólo $15$ positivos sobre $1\,011$
imágenes) ilustra que la representación SAE captura conceptos
poco frecuentes en el corpus sin contaminación por la
distribución mayoritaria. La cobertura sobre vasos polimorfos
($0{,}810$, $N_{+} = 140$) y sobre la red atípica
($0{,}732$, $N_{+} = 230$) sostiene el alineamiento entre el
vocabulario operativo del consenso clínico propuesto por la
Dra.~R.~Taberner y la representación aprendida por el SAE Large.

\subsection*{Ajuste fino multitarea: cabezas conceptuales y
melanoma}
\label{sec:sae-d7pt-e4}

El experimento E4 replica conceptualmente el protocolo multitarea
descrito por Kawahara et~al.~\cite{kawahara2019} sustituyendo la
red convolucional original por el codificador fundacional PanDerm
Large. El \emph{fine-tuning} se realiza mediante adaptación de
bajo rango (LoRA, $r = 16$) aplicada exclusivamente a los dos
últimos bloques del codificador (\emph{blocks.22-23}), de forma
idéntica al protocolo SpanDerm v0 reportado en la
secci\'on~\ref{sec:dermapixel-spanderm-v0}. Sobre la
representación adaptada se conectan ocho cabezas paralelas: una
cabeza multi-clase por cada uno de los siete criterios del
\emph{Seven-Point Checklist} y una cabeza binaria para
melanoma/no-melanoma. La función objetivo es la suma
equiponderada de entropías cruzadas sobre las ocho cabezas,
$\mathcal{L} = \sum_{i=1}^{7} \mathrm{CE}(\hat{c}_i, c_i) +
\mathrm{CE}(\hat{m}, m)$, con \emph{class weight} balanceado por
clase. La configuración de optimización emplea AdamW con
\emph{learning rate} diferenciado
($\eta_{\mathrm{heads}} = 10^{-3}$,
$\eta_{\mathrm{LoRA}} = 5 \times 10^{-4}$) durante $15$ épocas,
con selección del \emph{checkpoint} de mejor AUROC melanoma sobre
validación (época $13$, AUROC val $= 0{,}861$). El número total
de parámetros entrenables es $0{,}56$\,M ($0{,}18\%$ del total
del codificador), magnitud equivalente a la de SpanDerm v0.

La tabla~\ref{tab:sae-d7pt-e4-comp} compara el rendimiento de
E4 frente al sondeo lineal directo (E2 Direct LP sobre las
$16\,384$ activaciones del SAE) y al CBM jerárquico de siete
conceptos (E2 CBM) sobre el conjunto de test Derm7pt
($N_{\mathrm{test}} = 395$, $n_{\mathrm{mel}} \approx 98$). Los
intervalos de confianza al $95\%$ se estiman por \emph{bootstrap}
no paramétrico sobre el conjunto de test.

\begin{table}[H]
\centering
\footnotesize
\setlength{\tabcolsep}{4pt}
\renewcommand{\arraystretch}{0.95}
\caption{Comparativa sobre la tarea binaria melanoma/no-melanoma
entre sondeo lineal directo sobre las $16\,384$ \emph{features}
del SAE Large (E2 Direct LP), CBM jerárquico de siete conceptos
(E2 CBM) y \emph{fine-tuning} multitarea LoRA con ocho cabezas
paralelas (E4). Conjunto de test Derm7pt
($N_{\mathrm{test}} = 395$, $n_{\mathrm{mel}} \approx 98$), mismo
\emph{split} oficial. AUROC en negrita.}
\label{tab:sae-d7pt-e4-comp}
\begin{tabular}{lrrrrrr}
\hline
Método                                       & Params entr. & Acc@1     & BAcc      & AUROC              & W-F1      & Kappa     \\
\hline
E2 Direct LP (SAE $\to$ melanoma)             & $16$\,K       & $0{,}841$ & $0{,}735$ & $0{,}890$          & $0{,}829$ & $0{,}527$ \\
E2 CBM (SAE $\to$ $7$ conceptos $\to$ mel)    & $\sim 120$    & $0{,}825$ & $0{,}681$ & $0{,}832$          & $0{,}801$ & $0{,}440$ \\
\textbf{E4 FT multitarea LoRA}                & $0{,}56$\,M   & $\mathbf{0{,}841}$ & $\mathbf{0{,}804}$ & $\mathbf{0{,}891}$ & $\mathbf{0{,}843}$ & $\mathbf{0{,}590}$ \\
\hline
\end{tabular}
\end{table}

La tabla~\ref{tab:sae-d7pt-e4-concepts} reporta la AUROC binaria
\emph{present}/\emph{absent} por concepto sobre el mismo conjunto
de test y la compara con la cabeza intermedia del CBM (regresión
logística sobre activaciones SAE, una por concepto) reportada en
la tabla~\ref{tab:sae-d7pt-cbm-coefs}.

\begin{table}[H]
\centering
\footnotesize
\setlength{\tabcolsep}{4pt}
\renewcommand{\arraystretch}{0.95}
\caption{AUROC binaria \emph{present}/\emph{absent} por concepto
del \emph{Seven-Point Checklist} sobre el conjunto de test
Derm7pt ($N_{\mathrm{test}} = 395$). Comparativa entre sondeo
lineal separado por concepto sobre las activaciones SAE Large
(E2 LP separado, una regresión logística por concepto) y las
siete cabezas conceptuales del \emph{fine-tuning} multitarea
LoRA (E4). $\Delta$ en puntos porcentuales.}
\label{tab:sae-d7pt-e4-concepts}
\begin{tabular}{lrrr}
\hline
Concepto                & E2 LP separado & E4 multitarea & $\Delta$ \\
\hline
pigment\_network        & $0{,}901$ & $0{,}889$ & $-1{,}2$\,pp \\
vascular\_structures    & $0{,}906$ & $0{,}880$ & $-2{,}6$\,pp \\
blue\_whitish\_veil     & $0{,}895$ & $0{,}879$ & $-1{,}6$\,pp \\
streaks                 & $0{,}829$ & $0{,}818$ & $-1{,}1$\,pp \\
pigmentation            & $0{,}792$ & $0{,}793$ & $+0{,}1$\,pp \\
regression\_structures  & $0{,}818$ & $0{,}791$ & $-2{,}7$\,pp \\
dots\_and\_globules     & $0{,}723$ & $0{,}725$ & $+0{,}2$\,pp \\
\hline
\end{tabular}
\end{table}

Tres observaciones cuantitativas sintetizan E4. \textit{Primero},
la cabeza melanoma del \emph{fine-tuning} multitarea alcanza
AUROC test $= 0{,}891$ (IC95\% $[0{,}856,\,0{,}927]$), valor que
supera marginalmente al sondeo lineal directo sobre las
$16\,384$ \emph{features} del SAE ($0{,}890$, $+0{,}1$\,pp,
incremento dentro del ruido del \emph{bootstrap}) y mejora el
CBM jerárquico en $+5{,}9$\,pp AUROC. La diferencia operativa
respecto al Direct LP es que E4 ofrece simultáneamente predicción
melanoma e inferencia sobre los siete criterios diagnósticos en
una sola pasada, lo que preserva la interpretabilidad explícita
sin coste discriminativo apreciable sobre la tarea binaria
primaria. \textit{Segundo}, las siete cabezas conceptuales pierden
entre $1$ y $3$\,pp de AUROC respecto al sondeo lineal separado
por concepto: el coste de compartir representación entre tareas
auxiliares y tarea principal se sitúa en
$-1{,}2$\,pp (\emph{pigment network}), $-1{,}6$\,pp (\emph{blue
whitish veil}), $-2{,}6$\,pp (\emph{vascular structures}) y
$-2{,}7$\,pp (\emph{regression structures}); dos conceptos
mejoran ligeramente bajo el régimen multitarea
(\emph{pigmentation} $+0{,}1$\,pp,
\emph{dots and globules} $+0{,}2$\,pp). Este patrón es
consistente con el comportamiento documentado del aprendizaje
multitarea y cuantifica el compromiso aceptado a cambio de un
único codificador adaptado para las ocho cabezas
simultáneas. \textit{Tercero}, el hecho de que la cabeza melanoma
no sufra penalización medible respecto al sondeo lineal
especializado sugiere que el aprendizaje multitarea no degrada la
tarea principal cuando los conceptos auxiliares son
taxonómicamente coherentes con el objetivo. El
\emph{Seven-Point Checklist} es, por construcción, un instrumento
diagnóstico de melanoma~\cite{kawahara2019}, lo que justifica
empíricamente la elección del corpus de criterios.

Ocho hallazgos cuantitativos cierran la sección.
\textit{Primero}, existen \emph{features} SAE individuales que
codifican conceptos dermatoscópicos clásicos sin supervisión
explícita durante el preentrenamiento: AUROC efectivo hasta
$0{,}823$ sobre vascular structures (E1) y hasta $0{,}926$
sobre vasos en horquilla (E3, $N_{+} = 15$), lo que indica
que la base sparse aprendida sobre Derm1M no es solo un
mecanismo de reducción dimensional sino que ha capturado
representaciones interpretables alineadas con la taxonomía
dermatoscópica. \textit{Segundo}, el CBM jerárquico cuantifica
el compromiso interpretabilidad-precisión sobre la tarea binaria
melanoma/no-melanoma: el cuello de botella conceptual de siete
dimensiones reduce la AUROC en $-5{,}8$\,pp ($0{,}890 \to 0{,}832$)
respecto al sondeo lineal directo sobre las $16\,384$
activaciones, pérdida que el sistema acepta a cambio de
trazabilidad concepto-a-concepto en la decisión. \textit{Tercero},
los pesos del meta-clasificador del CBM ordenan
\emph{blue whitish veil} ($+2{,}54$) y \emph{regression
structures} ($+2{,}07$) como conceptos dominantes para la
discriminación de melanoma, jerarquía coherente con el peso
clínico asignado a estos criterios en el
\emph{Seven-Point Checklist} de
referencia~\cite{kawahara2019}. \textit{Cuarto}, $7$ de los
$11$ conceptos del Grupo~A propuestos a la Dra.~R.~Taberner
con \emph{mapping} operativo a Derm7pt quedan cubiertos por
\emph{features} SAE sin entrenamiento específico sobre el
corpus en castellano, y los $5$ conceptos sin \emph{mapping}
(empedrado, paralelos surcos/crestas, lagunas vasculares, vasos en
corona) constituyen precisamente el contenido específico de la
anotación clínica colaborativa documentada en la
secci\'on~\ref{sec:linea-derm7pt-sae} como línea de continuidad.
\textit{Quinto}, el SAE Large entrenado sobre activaciones de
PanDerm Large generaliza la interpretabilidad emergente desde
Derm1M ($507\,657$ imágenes) hasta Derm7pt
($1\,011$ lesiones dermatoscópicas) sin reentrenamiento, lo que
sostiene la línea de interpretabilidad mediante Sparse
Autoencoders documentada en el Anexo~\ref{cap:anexof}.
\textit{Sexto}, el \emph{fine-tuning} multitarea LoRA con ocho
cabezas paralelas (E4, secci\'on~\ref{sec:sae-d7pt-e4}) supera
marginalmente al sondeo lineal directo sobre las
$16\,384$ activaciones del SAE en AUROC binario melanoma
($0{,}891$ frente a $0{,}890$, incremento dentro del intervalo de
\emph{bootstrap}) y mejora el CBM jerárquico en $+5{,}9$\,pp
AUROC, preservando interpretabilidad explícita concepto-a-concepto
mediante las siete cabezas auxiliares activas
simultáneamente. \textit{Séptimo}, las cabezas conceptuales del
régimen multitarea pierden entre $1$ y $3$\,pp de AUROC respecto
al sondeo lineal separado por concepto sobre las activaciones SAE
($-1{,}2$ a $-2{,}7$\,pp sobre los cuatro conceptos más
sensibles, ganancias residuales sobre los dos restantes), patrón
consistente con el coste documentado del aprendizaje
compartido. \textit{Octavo}, la cabeza melanoma del régimen
multitarea no se degrada respecto al sondeo lineal especializado
($\Delta \mathrm{AUROC} = +0{,}1$\,pp), evidencia empírica de que
los siete criterios del \emph{Seven-Point
Checklist}~\cite{kawahara2019} constituyen auxiliares
taxonómicamente coherentes con la tarea binaria primaria de
discriminación melanoma/no-melanoma. El conjunto E1-E4 proporciona
evidencia consistente del valor del SAE Large entrenado sobre
PanDerm Large para tareas de interpretabilidad emergente en
dermatología, y refuerza la viabilidad del módulo M7 de
interpretabilidad del prototipo DermApIxel
(Anexo~\ref{cap:anexoh}).

\section{Codificadores adicionales y arquitecturas de inferencia
alternativas}
\label{sec:dermapixel-extra}

Las secciones~\ref{sec:dermapixel-results}, \ref{sec:dermapixel-zs}
y~\ref{sec:dermapixel-spanderm-v0} caracterizan codificadores
específicos de dominio (PanDerm Base/Large, DermLIP v2) y
vision-lenguaje médicos (BiomedCLIP, SigLIP-SO400M) sobre el
corpus DermapixelAI~1.0, con cabezas paramétricas ligeras y bajo
clasificación contrastiva \emph{zero-shot}. La presente sección
cierra el panorama experimental con dos extensiones complementarias
ejecutadas sobre NVIDIA DGX Spark GB10 el 2026-06-02:
\emph{(i)}~codificadores generalistas no médicos (DINOv2 ViT-L
auto-supervisado puro y CLIP-L de OpenAI contrastivo) como
referencia fuera del dominio, y \emph{(ii)}~dos arquitecturas
alternativas de inferencia sobre el mismo conjunto de
\emph{embeddings} ya cacheados: \emph{retrieval} $k$-NN sobre
índice FAISS con embeddings de PanDerm Large, y \emph{zero-shot
jerárquico} en cascada L1$\to$L2$\to$L3 condicionado sobre el
codificador DermLIP v2. El conjunto de cuatro experimentos
(referidos como N1-N4 en el material reproducible) acota el techo
discriminativo del corpus desde dos ángulos no explorados en las
secciones previas: el rendimiento de codificadores ajenos al
dominio biomédico y el efecto de modificar el clasificador o la
estrategia de inferencia sin cambiar el codificador.

\subsection*{Codificadores generalistas no médicos: DINOv2 y CLIP-L}
\label{sec:dermapixel-extra-encoders}

Se evalúan dos codificadores fuera del dominio biomédico bajo el
protocolo de sondeo lineal congelado descrito en la
secci\'on~\ref{sec:dermapixel-results}, sobre el mismo
\emph{split} fijo ($N_{\mathrm{train}} = 874$,
$N_{\mathrm{test}} = 36$ en L1/L2, $N_{\mathrm{test}} = 28$ en
L3). DINOv2 ViT-L (\texttt{vit\_large\_patch14\_dinov2.lvd142m},
embedding $1\,024$-D, CLS) corresponde al modelo auto-supervisado
visual generalista de Meta sobre LVD-142M~\cite{dinov2}; CLIP-L-336
(OpenAI, ViT-L/14 con resolución $336$) corresponde al modelo
contrastivo texto-imagen generalista entrenado sobre $400$
millones de pares web~\cite{clip2021}, con embedding proyectado de
$768$-D. La cabeza es regresión logística L-BFGS con $C = 1{,}0$ y
$\mathrm{max\_iter} = 5\,000$, idéntica a la
secci\'on~\ref{sec:dermapixel-results}.

\begin{table}[H]
\centering
\footnotesize
\setlength{\tabcolsep}{4pt}
\renewcommand{\arraystretch}{0.95}
\caption{Sondeo lineal sobre embeddings congelados de
codificadores generalistas no médicos en DermapixelAI~1.0
($N_{\mathrm{test}} = 36$ en L1/L2, $N_{\mathrm{test}} = 28$ en
L3). Referencia: PanDerm Large
(tabla~\ref{tab:dermapixel-lp}). Mejor valor por columna y nivel en
negrita.}
\label{tab:dermapixel-extra-encoders}
\begin{tabular}{lrrrrrr}
\hline
              & \multicolumn{2}{c}{L1 (4 cls)} & \multicolumn{2}{c}{L2 (38 cls)} & \multicolumn{2}{c}{L3 (224 cls)} \\
Encoder       & BAcc & AUROC & BAcc & AUROC & BAcc & AUROC \\
\hline
PanDerm Large & $\mathbf{0{,}735}$ & $\mathbf{0{,}796}$
              & $0{,}265$          & $0{,}793$
              & $0{,}184$          & $0{,}813$ \\
DINOv2 ViT-L  & $0{,}537$          & $0{,}650$
              & $0{,}250$          & $0{,}718$
              & $\mathbf{0{,}211}$ & $0{,}794$ \\
CLIP-L-336    & $0{,}556$          & $0{,}658$
              & $\mathbf{0{,}279}$ & $\mathbf{0{,}826}$
              & $0{,}184$          & $\mathbf{0{,}827}$ \\
\hline
\end{tabular}
\end{table}

CLIP-L-336 alcanza la mejor AUROC sobre L2 ($0{,}826$ frente a
$0{,}793$ de PanDerm Large, $+3{,}3$\,pp) y sobre L3
($0{,}827$ frente a $0{,}813$, $+1{,}4$\,pp) a pesar de no
incorporar material biomédico en su preentrenamiento. DINOv2 ViT-L,
auto-supervisado puro sin componente textual, gana la BAcc L3
($0{,}211$ frente a $0{,}184$ de PanDerm Large,
$+2{,}7$\,pp) pero pierde sistemáticamente sobre L1
($-19{,}8$\,pp BAcc, $-14{,}6$\,pp AUROC) y se mantiene por
debajo sobre L2 ($-1{,}5$\,pp BAcc, $-7{,}5$\,pp AUROC).
PanDerm Large conserva la primacía sobre L1 en BAcc y AUROC.
Ningún codificador domina simultáneamente en todas las métricas y
los tres niveles ontológicos, patrón coherente con la lectura de la
secci\'on~\ref{sec:dermapixel-results} sobre la complementariedad
entre el preentrenamiento auto-supervisado de PanDerm y el
contrastivo texto-imagen de DermLIP. La clasificación
\emph{zero-shot} con CLIP-L (\emph{prompts} en castellano e
inglés idénticos a la secci\'on~\ref{sec:dermapixel-zs}) resulta
catastrófica sobre los tres niveles (AUROC L1 $0{,}559$ en
castellano y $0{,}635$ en inglés, AUROC L2 $0{,}699$ en
castellano, AUROC L3 $0{,}617$ en castellano), confirmando que los
codificadores generalistas requieren adaptación supervisada para
operar sobre material clínico dermatológico.

\subsection*{\emph{Retrieval} $k$-NN con índice FAISS sobre
embeddings PanDerm Large}
\label{sec:dermapixel-extra-faiss}

Una alternativa al clasificador paramétrico sobre el espacio de
\emph{embeddings} congelados consiste en la recuperación basada en
similitud: dado un \emph{embedding} de consulta, se recuperan sus
$k$ vecinos más próximos del conjunto de entrenamiento bajo
similitud coseno y se predice por votación. Se construye un índice
\texttt{IndexFlatIP} de FAISS~\cite{faiss2019} sobre los
$874$ \emph{embeddings} L2-normalizados de PanDerm Large
correspondientes al conjunto de \emph{train} de la
secci\'on~\ref{sec:dermapixel-results} y se evalúan dos variantes
(votación por mayoría y \emph{softmax} ponderado por
similitud) cruzadas con $k \in \{1, 5, 10, 20\}$. La
tabla~\ref{tab:dermapixel-extra-faiss} reporta la mejor combinación
por nivel y métrica frente a la referencia LP paramétrica.

\begin{table}[H]
\centering
\footnotesize
\setlength{\tabcolsep}{4pt}
\renewcommand{\arraystretch}{0.95}
\caption{\emph{Retrieval} $k$-NN sobre embeddings PanDerm Large
con índice FAISS (\texttt{IndexFlatIP}, similitud coseno tras
L2-normalize). Mejor combinación $k \times $ variante por nivel y
métrica. Referencia: sondeo lineal paramétrico sobre el mismo
encoder (LP, tabla~\ref{tab:dermapixel-lp}). Mejor valor por
columna y nivel en negrita.}
\label{tab:dermapixel-extra-faiss}
\begin{tabular}{llrrr}
\hline
Nivel & Método & Acc@1 & BAcc & AUROC \\
\hline
L1 & LP paramétrico (LR)              & $0{,}750$          & $\mathbf{0{,}735}$ & $\mathbf{0{,}796}$ \\
L1 & $k$-NN $k=10$ \emph{majority}     & $0{,}722$          & $0{,}686$          & $0{,}748$ \\
\hline
L2 & LP paramétrico (LR)              & $0{,}250$          & $0{,}265$          & $\mathbf{0{,}793}$ \\
L2 & $k$-NN $k=20$ \emph{majority}     & $\mathbf{0{,}278}$ & $\mathbf{0{,}265}$ & $0{,}700$ \\
\hline
L3 & LP paramétrico (LR)              & $0{,}179$          & $0{,}184$          & $\mathbf{0{,}813}$ \\
L3 & $k$-NN $k=5$ \emph{majority}      & $\mathbf{0{,}214}$ & $\mathbf{0{,}210}$ & $0{,}615$ \\
L3 & $k$-NN $k=20$ \emph{weighted}     & $0{,}214$          & $0{,}210$          & $0{,}719$ \\
\hline
\end{tabular}
\end{table}

El \emph{retrieval} $k$-NN supera al sondeo lineal paramétrico
sobre el régimen de cola larga severa L3 ($224$ clases efectivas,
$3{,}5$ muestras/clase media en \emph{train}): $+3{,}5$\,pp
Acc@1 ($0{,}179 \to 0{,}214$) y $+2{,}6$\,pp BAcc
($0{,}184 \to 0{,}210$) con $k = 5$ y votación por mayoría. Sobre
L2, $k$-NN con $k = 20$ empata en BAcc ($0{,}265$) y mejora
Acc@1 en $+2{,}8$\,pp ($0{,}250 \to 0{,}278$). Sobre L1, en
cambio, $k$-NN pierde respecto al LP paramétrico
($-4{,}9$\,pp BAcc, $-4{,}8$\,pp AUROC), comportamiento
consistente con que la cardinalidad reducida del nivel ($4$ clases)
favorece a la regresión logística sobre la votación de vecinos. La
AUROC global del $k$-NN es sistemáticamente inferior a la del LP
(diferencias entre $-4{,}8$ y $-9{,}3$\,pp según nivel),
atribuible a que la estimación de probabilidad por frecuencia
relativa de vecinos no es densa sobre el conjunto de clases. El
hallazgo es coherente con la arquitectura del módulo M4 del
prototipo de producción (\emph{retrieval} sobre Derm1M descrito en
el Anexo~\ref{cap:anexoh}), que emplea el mismo principio de
similitud coseno sobre \emph{embeddings} PanDerm Large como
soporte conceptual al diagnóstico ofrecido por las cabezas
paramétricas.

\subsection*{\emph{Zero-shot} jerárquico: cascada L1$\to$L2$\to$L3
con propagación de errores}
\label{sec:dermapixel-extra-hierzs}

La secci\'on~\ref{sec:dermapixel-zs} evalúa la clasificación
\emph{zero-shot} sobre los tres niveles ontológicos de forma
independiente, considerando en cada caso todas las clases
candidatas del nivel. Una estrategia alternativa, motivada por el
flujo diagnóstico clínico real, consiste en una cascada
jerárquica: predecir primero L1, restringir las candidatas L2 a las
hijas del L1 predicho, y restringir las candidatas L3 a las hijas
del L2 predicho. La jerarquía L1$\to$L2 utilizada se deriva del
corpus de \emph{train} (Inflamatoria: $25$ subcategorías L2;
Infecciosa: $5$; Tumoral: $12$; Genodermatosis: $2$). El encoder
es DermLIP v2 (el mejor \emph{zero-shot} de la
secci\'on~\ref{sec:dermapixel-zs}), bajo \emph{prompts} en
castellano. Se contrastan tres variantes: el \emph{zero-shot}
plano de la secci\'on~\ref{sec:dermapixel-zs} (todas las clases del
nivel como candidatas), la cascada jerárquica estricta con el L1
$\arg\max$, y una variante \emph{top-3 L1} que mantiene las tres
categorías L1 más probables como candidatas y agrega sus hijas L2.

\begin{table}[H]
\centering
\footnotesize
\setlength{\tabcolsep}{4pt}
\renewcommand{\arraystretch}{0.95}
\caption{Comparación entre clasificación \emph{zero-shot} plana y
\emph{zero-shot jerárquico} en cascada sobre el conjunto de test de
DermapixelAI~1.0 ($N_{\mathrm{test}} = 36$ en L1/L2,
$N_{\mathrm{test}} = 28$ en L3). Encoder: DermLIP v2, prompts en
castellano. La fila ``L2 hier $|$ L1 correcto'' restringe la
evaluación a las muestras con L1 predicho correctamente. Mejor
valor por columna en negrita; la cifra condicionada se reporta
separadamente.}
\label{tab:dermapixel-extra-hierzs}
\begin{tabular}{lrrr}
\hline
Variante                                  & L1 Acc@1 & L2 Acc@1 & L3 Acc@1 \\
\hline
ZS plano (referencia, secci\'on~\ref{sec:dermapixel-zs}) & $\mathbf{0{,}361}$ & $\mathbf{0{,}222}$ & $\mathbf{0{,}107}$ \\
ZS jerárquico (cascada estricta)           & $0{,}361$          & $0{,}139$          & $0{,}036$ \\
ZS jerárquico (top-3 L1)                   & $0{,}361$          & $0{,}194$          & $0{,}036$ \\
\hline
L2 hier $|$ L1 correcto (condicionada)     & ---                & $0{,}385$          & --- \\
\hline
\end{tabular}
\end{table}

La cascada jerárquica estricta degrada el rendimiento plano sobre
L2 en $-8{,}3$\,pp Acc@1 ($0{,}222 \to 0{,}139$) y sobre L3 en
$-7{,}1$\,pp Acc@1 ($0{,}107 \to 0{,}036$). La causa
principal es la propagación de errores desde el nivel L1: dado que
DermLIP v2 acierta L1 sólo en el $36{,}1\%$ de los casos, el
$63{,}9\%$ restante recibe en L2 un conjunto de candidatas
desalineado con la verdad de terreno, lo que invalida toda
predicción posterior independientemente del codificador. La
variante \emph{top-3 L1} amortigua parcialmente la propagación al
mantener las tres categorías L1 más probables como hipótesis activa
y recupera $0{,}194$ Acc@1 sobre L2 ($-2{,}8$\,pp respecto a
plano), sin efecto sobre L3. Condicionado a L1 correcto, la
cascada jerárquica L2 alcanza Acc@1 $= 0{,}385$, casi el doble
del plano, evidencia de que la restricción del espacio de
candidatas a las subcategorías clínicamente coherentes con la
categoría etiológica acertada incrementa la discriminación de
forma sustancial. La implicación clínica es directa: la arquitectura
jerárquica funciona si la decisión inicial sobre la familia
etiológica es robusta; con \emph{zero-shot} modesto, la cascada
acumula errores y degrada el rendimiento plano.

\subsection*{Síntesis: panorama completo de codificadores y
arquitecturas sobre DermapixelAI 1.0}
\label{sec:dermapixel-extra-sintesis}

La tabla~\ref{tab:dermapixel-extra-sintesis} consolida las cifras
de BAcc y AUROC sobre los tres niveles ontológicos para el conjunto
completo de codificadores y arquitecturas evaluados sobre el corpus
DermapixelAI~1.0 a lo largo del presente capítulo. La síntesis
permite identificar el mejor método por nivel y métrica sin
dispersión entre las secciones de origen.

\begin{table}[H]
\centering
\footnotesize
\setlength{\tabcolsep}{4pt}
\renewcommand{\arraystretch}{0.95}
\caption{Síntesis de codificadores y arquitecturas evaluados sobre
DermapixelAI~1.0 (test fijo: $N = 36$ en L1/L2, $N = 28$ en L3).
La columna \emph{Ref.} remite a la sección de origen de las cifras
completas con sus intervalos de confianza. Mejor valor por columna
y nivel en negrita.}
\label{tab:dermapixel-extra-sintesis}
\begin{tabular}{llrrrrrr}
\hline
                                  &     & \multicolumn{2}{c}{L1} & \multicolumn{2}{c}{L2} & \multicolumn{2}{c}{L3} \\
Método                            & Ref.& BAcc & AUROC           & BAcc & AUROC           & BAcc & AUROC \\
\hline
PanDerm Base LP                   & §4.11 & $0{,}570$          & $0{,}697$
                                          & $0{,}221$          & $0{,}707$
                                          & $0{,}132$          & $0{,}735$ \\
PanDerm Large LP                  & §4.11 & $\mathbf{0{,}735}$ & $\mathbf{0{,}796}$
                                          & $0{,}265$          & $0{,}793$
                                          & $0{,}184$          & $0{,}813$ \\
DermLIP v2 LP                     & §4.11 & ---                & ---
                                          & ---                & $0{,}860$
                                          & ---                & --- \\
DINOv2 ViT-L LP                   & §4.16 & $0{,}537$          & $0{,}650$
                                          & $0{,}250$          & $0{,}718$
                                          & $0{,}211$          & $0{,}794$ \\
CLIP-L-336 LP                     & §4.16 & $0{,}556$          & $0{,}658$
                                          & $0{,}279$          & $\mathbf{0{,}826}$
                                          & $0{,}184$          & $\mathbf{0{,}827}$ \\
SpanDerm v0 LoRA                  & §4.14 & ---                & ---
                                          & $\mathbf{0{,}363}$ & $0{,}855$
                                          & ---                & --- \\
FT parcial denso L1$+$L2$+$L3      & §4.11 & $0{,}622$          & $0{,}715$
                                          & $0{,}324$          & $0{,}852$
                                          & $0{,}184$          & $0{,}804$ \\
$k$-NN FAISS PanDerm Large         & §4.16 & $0{,}686$          & $0{,}748$
                                          & $0{,}265$          & $0{,}700$
                                          & $\mathbf{0{,}210}$ & $0{,}719$ \\
\hline
\end{tabular}
\end{table}

El mejor método sobre DermapixelAI~1.0 depende del nivel
ontológico y de la métrica. Sobre L1, el sondeo lineal sobre
PanDerm Large encabeza simultáneamente BAcc ($0{,}735$) y
AUROC ($0{,}796$); sobre L2, SpanDerm v0 LoRA encabeza la BAcc
($0{,}363$) y CLIP-L-336 LP encabeza la AUROC ($0{,}826$); sobre
L3, el \emph{retrieval} $k$-NN FAISS sobre embeddings de PanDerm
Large encabeza la BAcc ($0{,}210$) y CLIP-L-336 LP encabeza la
AUROC ($0{,}827$). La discriminación binaria melanoma/no-melanoma
sobre Derm7pt (secci\'on~\ref{sec:sae-derm7pt}, fuera de
DermapixelAI~1.0) alcanza AUROC $0{,}891$ con el \emph{fine-tuning}
multitarea LoRA E4 sobre las activaciones del SAE. La ausencia de
un método único óptimo refuerza la lectura interpretativa de la
secci\'on~\ref{sec:disc-especializacion}: la elección del
codificador y del clasificador debe condicionarse al nivel
ontológico objetivo y al criterio operativo predominante
(discriminación robusta, ranking probabilístico o recuperación
para soporte explicativo).

\section{SpanDerm-CLIP: zero-shot multimodal castellano sobre
DermapixelAI~1.0}
\label{sec:spanderm-clip}

Las secciones previas caracterizan la familia DermLIP en
\emph{zero-shot} sobre datasets externos
(secci\'on~\ref{sec:zs-results}) y sobre DermapixelAI~1.0 con
plantillas en castellano e inglés
(secci\'on~\ref{sec:dermapixel-zs}), y SpanDerm v0 como cabeza
supervisada con LoRA sobre L2
(secci\'on~\ref{sec:dermapixel-spanderm-v0}). Las dos primeras
operan en régimen multimodal con encoder textual fijo
(PubMedBERT/CLIP) y la tercera en régimen supervisado con
vocabulario cerrado de $38$ clases L2. La presente sección
cierra una tercera configuración no cubierta por las anteriores:
un sistema vision-lenguaje contrastivo con encoder textual
multilingüe entrenado sobre el propio corpus DermapixelAI~1.0,
evaluado en \emph{zero-shot} \emph{open-class} sobre el vocabulario
completo de la ontología (L1, L2 y L3) en castellano. La
diferencia operativa frente a SpanDerm v0 es la capacidad de
mapear texto clínico libre en castellano sobre el espacio de
imágenes sin necesidad de fijar las clases por anticipado, condición
operacional relevante para escenarios de búsqueda
documental clínica sobre el archivo de la Dra.~R.~Taberner. La
diferencia operativa frente a DermLIP es la sustitución del encoder
textual anglosajón (PubMedBERT-256) por un encoder multilingüe
adaptado al lenguaje clínico castellano del propio corpus. La
justificación cuantitativa frente a un \emph{wrapper} alternativo
con modelo fundacional inglés más traducción ES--EN se desarrolla
en el Anexo~\ref{cap:anexoi}.

\subsection*{Arquitectura, datos y protocolo experimental}

El sistema, denominado SpanDerm-CLIP, sigue el patrón dual-encoder
estilo CLIP~\cite{clip2021}: un encoder visual de
$1\,024$ dimensiones (PanDerm Large, ViT-L/16, CLS \emph{pooling}),
un encoder textual multilingüe de $768$ dimensiones
(\texttt{intfloat/multilingual-e5-base}), capas de proyección
lineales con normalización L2 a un espacio compartido de
$512$ dimensiones y entrenamiento por pérdida InfoNCE bidireccional
(imagen$\to$texto y texto$\to$imagen) con temperatura aprendible.
El corpus de entrenamiento es DermapixelAI~1.0
(Anexo~\ref{cap:anexoc}) bajo el \emph{split} per-case estratificado
por L1 ($854$ \emph{train} / $107$ \emph{val} / $101$ \emph{test}
imágenes; $521 / 66 / 66$ casos), idéntico al utilizado en las
secciones~\ref{sec:dermapixel-results} y~\ref{sec:dermapixel-zs}.
El manifest SHA-256 del \emph{split} train/val/test
(\texttt{5dcd9b52\dots cd394538}) queda fijado en cada
\emph{checkpoint} y los conjuntos de validación y test no se tocan
en ninguna iteración del experimento; el manifest de los pares
cross-modal se registra por separado en el repositorio de
provenance.

Para amplificar la señal textual del corpus se generan, mediante
asistencia de un modelo de lenguaje grande partiendo del texto original
de cada caso del blog y de sus metadatos, ocho perspectivas
clínicas por imagen (morfología, dermatoscopia, criterios ABCDE,
diagnóstico diferencial, informe clínico, estilo blog,
\emph{case\_summary} y \emph{key\_findings}). El total alcanza
$7\,680$ captions sobre $960$ imágenes (corpus de la fase v1--v4);
$102$ imágenes fallan la
generación automática y conservan únicamente el texto base. El
coste real de generación es de $32{,}54$ USD ($0{,}030$ USD por
imagen completa) y la longitud media es de $87{,}5$ palabras por
caption; coste y longitud media proceden del registro de
provenance del repositorio. Combinadas con las nueve variantes textuales base
(plantillas L1/L2/L3, jerárquicas \texttt{L1L2L3}, título de blog,
tres \emph{chunks} de \texttt{case\_text} y
\texttt{diagnosis\_raw}) el corpus de entrenamiento alcanza del
orden de $22\,000$ pares (imagen, texto). Cada \emph{epoch} muestra
cada imagen una vez con una variante textual seleccionada por hash
determinista \texttt{(seed, epoch, image\_id)}, lo que garantiza
reproducibilidad bajo reanudación.

Se entrenan cinco iteraciones del sistema bajo configuraciones
crecientes de capacidad arquitectónica y, en la última, bajo
ampliación del corpus de entrenamiento. La progresión articula tres
palancas independientes ---adaptación del lado textual, adaptación
del lado visual y escala del corpus--- cuya contribución marginal
se aísla por construcción. La
tabla~\ref{tab:spanderm-config} resume las cinco configuraciones.

\begin{table}[H]
\centering
\footnotesize
\setlength{\tabcolsep}{4pt}
\renewcommand{\arraystretch}{0.95}
\caption{Configuración de las cinco iteraciones SpanDerm-CLIP
entrenadas sobre DermapixelAI~1.0. Las versiones v1--v4 comparten
el mismo \emph{split} per-case ($854$ \emph{train} / $107$ \emph{val} /
$101$ \emph{test}) y el mismo manifest SHA-256. La versión v5
amplía el conjunto de entrenamiento con $15\,128$ imágenes externas
(PAD-UFES-20~\cite{padufes} y DermNet~\cite{dermnet}) reponderadas
mediante \emph{Weighted Random Sampler} con ratio efectivo $6\times$
en favor del subconjunto castellano original, manteniendo los
conjuntos de validación y test idénticos.}
\label{tab:spanderm-config}
\begin{tabular}{llll}
\hline
Versión & Image encoder                                                  & Text encoder              & Train \\
\hline
v1      & PanDerm-Large frozen                                            & E5-base \emph{full FT}    & $854$ castellano \\
v2      & PanDerm-Large frozen                                            & E5-base $+$ LoRA $r=16$, $\alpha=32$ & $854$ castellano \\
v3      & PanDerm-Large $+$ LoRA $r=16$ (últimos $4/24$ bloques)           & E5-base $+$ LoRA $r=16$, $\alpha=32$ & $854$ castellano \\
v4      & PanDerm-Large $+$ LoRA $r=32$ (últimos $12/24$ bloques)          & E5-base $+$ LoRA $r=32$, $\alpha=64$ & $854$ castellano \\
v5      & PanDerm-Large $+$ LoRA $r=32$ (últimos $12/24$ bloques)          & E5-base $+$ LoRA $r=32$, $\alpha=64$ & $854$ cast. $+\,15\,128$ ext. \\
\hline
\end{tabular}
\end{table}

El paso v1$\to$v2 verifica si los $277$\,M parámetros del E5
\emph{full fine-tune} sobreajustan sobre $854$ imágenes y si una
adaptación LoRA del lado textual captura mejor el lenguaje rico
clínico castellano sin destruir el preentrenamiento multilingüe.
El paso v2$\to$v3 testea la hipótesis de que el cuello de botella
está en el lado visual, mediante inyección de LoRA en los
últimos cuatro bloques ViT (\texttt{attn.qkv}, \texttt{attn.proj},
\texttt{mlp.fc1}, \texttt{mlp.fc2}) del encoder visual. El paso
v3$\to$v4 amplía la capacidad arquitectónica del lado visual
duplicando el rango ($r=16\to32$) y triplicando la cobertura de
bloques ($4\to12$ últimos bloques del ViT), bajo la hipótesis ---
formulada tras la medición del techo \emph{Linear Probing}
discutida más adelante en esta misma sección--- de que el ajuste
confinado de v3 constituía capacidad insuficiente para reesculpir
las \emph{features} visuales en el régimen \emph{fine-grained}
castellano. El paso v4$\to$v5 disocia la palanca arquitectónica de
la palanca de escala: bajo idéntica configuración LoRA que v4 se
amplía el corpus de entrenamiento con $15\,128$ imágenes externas
de habla portuguesa y vocabulario clínico afín, acompañadas de
\emph{captions} sintéticas generadas mediante un modelo de lenguaje
grande con el mismo protocolo aplicado al subconjunto castellano,
para evaluar si la escala externa genérica resuelve la cola larga
($327$ clases L3 con $\leq 5$ imágenes considerando el vocabulario
ontológico completo Taberner, de las cuales $178$ disponen de
muestras $1$--$5$ en entrenamiento y $149$ carecen de imágenes en
el corpus actual) o degrada la alineación específica del registro
lingüístico original. Los tiempos de entrenamiento sobre NVIDIA DGX Spark GB10
son $\sim 18$\,min (v1), $\sim 3$\,min (v2), $\sim 5$\,min (v3),
$\sim 9$\,min (v4) y $\sim 54$\,min (v5). El objetivo experimental,
fijado en el diseño antes de observar resultados, es L3
\emph{accuracy}@1 $\geq 0{,}25$ en \emph{zero-shot} sobre las $54$
clases L3 representadas en el conjunto de test.

Una lección metodológica derivada de la generación de
\emph{captions} sintéticas del corpus externo de v5 (coste de
generación $42{,}33$ USD) resultó de calibrar la longitud
objetivo respecto al umbral de validación del verificador
automático posterior. Un primer intento con rango $30$--$40$
palabras produjo una tasa de rechazo del $54\%$, dado que las
\emph{captions} se agrupaban en el suelo del rango y el
verificador descartaba sistemáticamente las más cortas;
recalibrar el objetivo a $40$--$55$ palabras redujo el rechazo a
cero. La regla operativa derivada es dimensionar la longitud
objetivo de generación con margen explícito sobre el umbral
mínimo aceptable, no sobre la media deseada.

\subsection*{Resultados \emph{zero-shot}: progresión metodológica
v1$\to$v4 y ablación de escala v5}

La tabla~\ref{tab:spanderm-zs} reporta la \emph{accuracy}@1 sobre
las $101$ imágenes del conjunto de test bajo régimen híbrido
\emph{H6} homogéneo (combinación lineal de logits visuales y de
logits textuales sobre plantillas ricas con pesos $w_v = 0{,}3$ y
$w_t = 0{,}7$ respectivamente, calibrado sobre el conjunto de
validación independientemente del conjunto de test), aplicado de
forma uniforme a las cinco iteraciones para garantizar la
comparativa cross-checkpoint directa. Para v1--v3 los prototipos
visuales empleados en \emph{H6} derivan del espacio embebido del
checkpoint sin LoRA visual; para v4--v5 derivan del espacio
embebido del checkpoint con \emph{image-LoRA} expandida. La
comparativa permanece interpretable porque los prototipos derivan
del propio espacio embebido del \emph{checkpoint} evaluado en cada
caso. Los niveles L1 y L2 \emph{zero-shot} sólo se reportan para
v4 y v5 por constituir éstos los \emph{checkpoints} relevantes
para los regímenes operativos del prototipo.

\begin{table}[H]
\centering
\footnotesize
\setlength{\tabcolsep}{3pt}
\renewcommand{\arraystretch}{0.95}
\caption{Clasificación \emph{zero-shot} SpanDerm-CLIP sobre el
conjunto de test de $101$ imágenes (per-case estratificado por
L1) bajo régimen híbrido \emph{H6} homogéneo
($w_v = 0{,}3$, $w_t = 0{,}7$) aplicado a las cinco iteraciones.
La columna \emph{Objetivo} corresponde al umbral fijado en el
diseño experimental antes de observar resultados. Mejor cifra
por fila excluyendo la columna \emph{Objetivo} en negrita.
Provenance: commit \texttt{d0190c3} del repositorio
\texttt{spanderm-clip}.}
\label{tab:spanderm-zs}
\begin{tabular}{lrrrrrrr}
\hline
Tarea                        & Random    & v1                  & v2                  & v3                  & v4                  & v5                  & Objetivo \\
\hline
L1 acc@1 ($4$ cls)            & $0{,}250$ & ---                 & ---                 & ---                 & $0{,}4257$          & $\mathbf{0{,}4554}$ & ---           \\
L2 acc@1 ($23$ cls)           & $0{,}043$ & ---                 & ---                 & ---                 & $0{,}1837$          & $\mathbf{0{,}2041}$ & ---           \\
\textbf{L3 acc@1 ($54$ cls)}  & $0{,}019$ & $0{,}1782$          & $0{,}1980$          & $0{,}1980$          & $\mathbf{0{,}2475}$ & $0{,}2079$          & $\geq 0{,}25$ \\
\hline
\end{tabular}
\end{table}

La progresión v1$\to$v4 articula la respuesta cuantitativa al
diseño experimental sobre el nivel L3 bajo régimen \emph{H6}
homogéneo. Las tres primeras iteraciones quedan acotadas en el
intervalo $[0{,}1782,\,0{,}1980]$ sobre L3 \emph{accuracy}@1 con
dispersión del orden de $\pm 0{,}02$, equivalente a $\pm 2$
imágenes sobre $N = 101$, lo que sitúa a v1, v2 y v3
estadísticamente empatadas. La
figura~\ref{fig:spanderm-epoch-sweep} muestra que v3 oscila en el
intervalo $[0{,}04,\,0{,}12]$ a lo largo de los $10$ \emph{epochs}
sin aproximarse al objetivo. La cuarta iteración alcanza
$0{,}2475$ sobre L3 \emph{accuracy}@1, esencialmente el umbral
paper-grade fijado a priori: la diferencia de $-0{,}25$\,pp
respecto al objetivo $0{,}25$ corresponde a una unidad de
\emph{accuracy} sobre $N = 101$ ($0{,}0099$) y se interpreta como
ruido de muestreo sobre el conjunto de test. El salto $+4{,}95$\,pp
L3 entre v3 ($0{,}1980$) y v4 ($0{,}2475$) bajo idéntico régimen
\emph{H6} se atribuye exclusivamente a la ampliación de la
capacidad arquitectónica del image encoder (LoRA $r=32$ sobre los
últimos doce bloques frente a $r=16$ sobre cuatro), aislada del
efecto del régimen de inferencia. La diferencia v1$\to$v4
($0{,}1782 \to 0{,}2475$, $+6{,}93$\,pp, mejora relativa $\sim 39\%$)
articula la respuesta global del experimento bajo el mismo
régimen evaluativo. El salto resuelve la lectura cualitativa
que las tres primeras iteraciones permitían anticipar: el cuello
de botella no estaba en la formulación del lado textual ni en la
falta de \emph{captions} ricas, sino en la capacidad del lado
visual para reesculpir las \emph{features} \emph{fine-grained}
castellanas. La evidencia complementaria es el lift retrieval
$14{,}0\times$ sobre el baseline aleatorio ya alcanzado por v2
con LoRA textual ligera
(reportado más adelante en esta misma sección), que indica
alineación cross-modal lograda sobre captions ricas sin que ello
trasladase el \emph{zero-shot} L3 por encima del techo del
encoder visual congelado.

La quinta iteración constituye una ablación deliberada de la
hipótesis de escala. Bajo idéntica configuración LoRA que v4 y
con el mismo régimen \emph{H6} en inferencia, la ampliación del
corpus de entrenamiento con $15\,128$ imágenes externas
genéricas (ratio efectivo $6\times$ en favor del subconjunto
castellano original) reduce el L3 \emph{accuracy}@1 a $0{,}2079$
($-4{,}0$\,pp respecto a v4). El resultado se interpreta como
\emph{flat-to-negative} sobre el régimen híbrido y se discute
mecanísticamente en
secci\'on~\ref{sec:spanderm-escala-especificidad}: el encoder
visual mejora aisladamente sobre régimen visual puro
(\emph{H1} $0{,}168 \to 0{,}238$ entre v4 y v5, $+41{,}2\%$ relativo),
el techo \emph{Linear Probing} L2 supervisado se eleva
$+54{,}4\%$ relativo, y los niveles L1 y L2 \emph{zero-shot}
también mejoran bajo \emph{H6} (L1 acc@1 $0{,}4257 \to 0{,}4554$,
L2 acc@1 $0{,}1837 \to 0{,}2041$), pero el régimen híbrido sobre
texto rico castellano para el nivel L3 \emph{fine-grained} se
degrada porque las captions externas erosionan la alineación
específica del registro lingüístico del archivo original. La
$\emph{long tail}$ no se rellena: las $327$ clases L3 con $\leq 5$
imágenes en el subconjunto castellano permanecen sin cobertura
adicional tras la incorporación del corpus externo.

\begin{figure}[H]
\centering
\includegraphics[width=0.88\textwidth]{fig_spanderm_epoch_sweep.png}
\caption{Barrido de \emph{epochs} de SpanDerm-CLIP v3 sobre el
conjunto de test de $101$ imágenes. \emph{Accuracy}@1 L1, L2 y L3
y \texttt{val i2t\_R@5} frente al paso de entrenamiento. L3
oscila en el intervalo $[0{,}04,\,0{,}12]$ a lo largo de los diez
\emph{epochs} sin aproximarse al objetivo $\geq 0{,}25$, mientras
que \texttt{val i2t\_R@5} alcanza su máximo en \emph{epoch} $2$ y
desciende con el entrenamiento conjunto adicional del codificador
visual.}
\label{fig:spanderm-epoch-sweep}
\end{figure}

\subsection*{Recuperación cross-modal: resultado favorable
paper-grade}

El segundo eje de evaluación es la recuperación cross-modal
per-imagen sobre los $107$ pares (imagen, caption rica) del
conjunto de validación, con script idéntico para las tres
versiones. La tabla~\ref{tab:spanderm-retrieval} reporta el
régimen \emph{rich-caption} (campo \texttt{llm\_case\_summary},
el más representativo del lenguaje clínico real de la
Dra.~R.~Taberner).

\begin{table}[H]
\centering
\footnotesize
\setlength{\tabcolsep}{4pt}
\renewcommand{\arraystretch}{0.95}
\caption{Recuperación cross-modal i2t R@$5$ de SpanDerm-CLIP sobre
los $107$ pares (imagen, caption rica) de validación, régimen
\emph{rich-caption} (\texttt{llm\_case\_summary}). La columna
\emph{Lift} reporta el factor multiplicativo sobre el baseline
aleatorio (\emph{pool} $= 107$, baseline $R@5 = 5/107 = 0{,}0467$).
v2 marca el pico histórico de alineación cross-modal castellano;
v5 muestra descenso coherente con la degradación documentada en
la tabla~\ref{tab:spanderm-zs} sobre el régimen híbrido. v1 y v3
no se reportan por insuficiencia de datos en el régimen
homogéneo cross-checkpoint. El \emph{lift} sobre random se
calcula como $R@5_{\mathrm{abs}} / (k / N_{\mathrm{pool}})$, con
$k = 5$ y $N_{\mathrm{pool}} = 107$; normaliza respecto a la
dificultad intrínseca del \emph{pool} y constituye la métrica
honesta para comparar entre \emph{checkpoints}.}
\label{tab:spanderm-retrieval}
\begin{tabular}{lrrr}
\hline
Versión & i2t R@$5$ absoluto  & Random baseline & Lift sobre random        \\
\hline
v2      & $\mathbf{0{,}654}$  & $0{,}047$       & $\mathbf{14{,}0\times}$  \\
v4      & $0{,}579$           & $0{,}047$       & $12{,}4\times$           \\
v5      & $0{,}505$           & $0{,}047$       & $10{,}8\times$           \\
\hline
\end{tabular}
\end{table}

La alineación cross-modal castellano-imagen alcanza su pico
histórico en v2: i2t R@$5 = 0{,}654$ frente al baseline aleatorio
$0{,}047$ sobre \emph{pool} de $107$, equivalente a
$\mathbf{14{,}0\times}$ lift sobre random. Esta cifra constituye
la contribución empírica favorable del experimento bajo régimen
\emph{rich-caption}. Las iteraciones posteriores conservan lift
sustancial (v4 $12{,}4\times$, v5 $10{,}8\times$) pero descienden
respecto al pico, en coherencia con la degradación documentada en
la tabla~\ref{tab:spanderm-zs} sobre régimen híbrido: el
entrenamiento conjunto con LoRA visual añadida (v4) y la posterior
incorporación de \emph{captions} externas (v5) desplazan el
espacio embebido respecto al óptimo \emph{rich-caption} castellano
puro alcanzado por v2. El descenso v2$\to$v5 del lift
($14{,}0\times \to 10{,}8\times$) acompaña cuantitativamente la
hipótesis de erosión de la \emph{fingerprint} clínica castellana
discutida en
secci\'on~\ref{sec:spanderm-escala-especificidad}.

\begin{figure}[H]
\centering
\includegraphics[width=0.88\textwidth]{fig_spanderm_retrieval.png}
\caption{Recuperación cross-modal i2t R@$5$ de SpanDerm-CLIP
sobre $107$ pares de validación, régimen \emph{rich-caption}.
Pico v2 de $\mathbf{14{,}0\times}$ lift sobre random
($0{,}047 \to 0{,}654$) atribuible al conjunto LoRA textual
$+$ captions sintéticas ricas generadas mediante un modelo de
lenguaje grande. v4 y v5 conservan lift sustancial
($12{,}4\times$ y $10{,}8\times$ respectivamente) pero descienden
respecto al pico, en coherencia con la degradación documentada
sobre el régimen híbrido.}
\label{fig:spanderm-retrieval}
\end{figure}

\subsection*{Techo de \emph{Linear Probing} y diagnóstico
mecanístico}

La pregunta clave del experimento es por qué ningún cambio en el
lado textual mueve la aguja \emph{zero-shot} L3. La respuesta
surge de la medición rigurosa del techo \emph{Linear Probing}
sobre las \emph{features} $1\,024$-D del encoder visual, ejecutada
con script idéntico para v2 (encoder visual frozen) y v3 (LP en
\emph{epoch} $2$ y \emph{epoch} $6$). La
tabla~\ref{tab:spanderm-lp-ceiling} reúne los tres puntos de la
comparación.

\begin{table}[H]
\centering
\footnotesize
\setlength{\tabcolsep}{4pt}
\renewcommand{\arraystretch}{0.95}
\caption{Techo \emph{Linear Probing} supervisado sobre
\emph{features} $1\,024$-D del image encoder. Comparativa
\emph{frozen} (v2 e$7$) frente a image-LoRA tras entrenamiento
adicional (v3 \emph{epoch} $6$). Métrica: \emph{accuracy}@1. Las
cifras de v3 corresponden a la evaluación intermedia conservada
del Sprint97; la reejecución posterior arroja desviaciones
$< 0{,}01$ absolutas, consistentes con ruido de muestreo
($N = 101$) y sin alterar la lectura ($+0{,}01 \approx$ ruido).}
\label{tab:spanderm-lp-ceiling}
\begin{tabular}{lrrrr}
\hline
Nivel              & LP v2-e$7$ frozen & LP v3 best.pt (e$2$) & LP v3 e$6$ ($+$image train) & $\Delta$ frozen$\to$trained \\
\hline
L1 ($4$ vías)       & $0{,}584$         & $0{,}584$            & $0{,}594$                   & $+0{,}01$                   \\
L2 ($38$ vías)      & $0{,}214$         & $0{,}214$            & $0{,}225$                   & $+0{,}01$                   \\
L3 ($225$ vías)     & $0{,}079$         & $0{,}079$            & $0{,}089$                   & $+0{,}01$                   \\
\hline
\end{tabular}
\end{table}

Dos observaciones estructuran el diagnóstico paper-grade. En
primer lugar, el \emph{zero-shot} L3 alcanzado por v1/v2/v3
($\sim 0{,}13$) supera ya el techo LP supervisado
($0{,}079$--$0{,}089$) sobre el mismo encoder visual frozen.
Las \emph{features} $1\,024$-D del PanDerm-Large congelado no
separan linealmente los diagnósticos castellanos finos al nivel
necesario para superar L3 $\sim 0{,}15$. Ningún ajuste del lado
textual puede extraer información discriminativa que las
\emph{features} visuales no contienen. En segundo lugar, la
adaptación LoRA visual de v3 (últimos cuatro bloques,
$1{,}05$\,M parámetros) eleva el techo LP sólo $+0{,}01$ en
cada nivel, del orden de una imagen de test, dentro del ruido.
La LoRA confinada a los últimos cuatro bloques constituye
capacidad insuficiente para reesculpir las \emph{features} de
PanDerm-Large en el régimen \emph{fine-grained} castellano del
archivo. El veredicto técnico es \emph{``image LoRA did not
help''} (no destruyó el preentrenamiento, el LP nunca cayó, pero
tampoco lo mejoró sustancialmente).

El mecanismo por el cual el \emph{zero-shot} de v3 cae respecto
a v2 sin que el LP suba es interpretable: el entrenamiento
conjunto desplaza los embeddings visuales hacia el espacio
contrastivo de captions ricas y los aleja del espacio de
plantillas de nombre de clase que el \emph{zero-shot} usa
internamente. La evidencia directa es que L1 \emph{zero-shot}
\emph{accuracy}@1 cae $0{,}426$ (e$2$) $\to 0{,}287$ (e$7$)
mientras el \emph{train loss} sigue descendiendo a $0{,}29$, lo
que indica sobreajuste sobre el régimen \emph{rich-caption}.

La cuarta y quinta iteraciones aportan la verificación cuantitativa
del diagnóstico mecanístico anterior bajo las dos rutas
candidatas de solución: capacidad arquitectónica adicional sobre
el corpus original (v4) y escala externa genérica bajo capacidad
idéntica a v4 (v5). La ruta arquitectónica se confirma: la
extensión de LoRA visual de $r=16$ sobre cuatro bloques a $r=32$
sobre los últimos doce bloques eleva el L3 \emph{accuracy}@1
\emph{zero-shot} desde $0{,}119$ (v3 \emph{best-epoch}) a
$0{,}2475$ (v4 régimen \emph{H6}), un salto de $+12{,}9$\,pp que
sitúa el sistema esencialmente en el umbral paper-grade
$0{,}25$. La lectura es directa: el cuello de botella
identificado en v3 ---LoRA confinada a cuatro bloques con
$1{,}05$\,M parámetros--- estaba infradotado, y la palanca real
era la capacidad arquitectónica del image encoder, no la
formulación textual ni el volumen de \emph{captions}. La ruta de
escala, en cambio, no replica la mejora: bajo configuración LoRA
idéntica a v4, la ampliación del corpus con $15\,128$ imágenes
externas degrada el L3 \emph{accuracy}@1 \emph{zero-shot} a
$0{,}2079$, $-4{,}0$\,pp respecto a v4, pese a que el techo
\emph{Linear Probing} L2 supervisado se eleva $+54{,}4\%$ relativo
en BAcc sobre el mismo encoder visual evaluado de forma aislada.
La conjunción de ambas
observaciones refina el diagnóstico paper-grade: en corpora
clínicos pequeños de alta especificidad lingüística, la palanca
correcta para superar el techo $\sim 0{,}13$ del encoder
congelado es ampliar capacidad arquitectónica del lado visual
sobre el propio corpus, no incorporar imágenes externas con
\emph{captions} sintéticas generadas mediante un modelo de
lenguaje grande sobre dominios afines pero estilísticamente
discordantes.

\subsection*{Calibración por \emph{temperature scaling}}

Los logits \emph{zero-shot} de SpanDerm-CLIP están severamente
sobreconfiados, propiedad esperada en arquitecturas CLIP-style
entrenadas con InfoNCE. La calibración con \emph{temperature
scaling} de un único parámetro escalar~\cite{guo2017} sobre los
$107$ pares de validación reduce el \emph{Expected Calibration
Error} (ECE) medido sobre los $101$ pares de test. La
tabla~\ref{tab:spanderm-calibration} sintetiza el resultado sobre
el \emph{checkpoint} paper-grade v4 y la
figura~\ref{fig:spanderm-calibration} ilustra el diagrama de
fiabilidad para L3.

\begin{table}[H]
\centering
\footnotesize
\setlength{\tabcolsep}{4pt}
\renewcommand{\arraystretch}{0.95}
\caption{Calibración post-hoc por \emph{temperature scaling} sobre
SpanDerm-CLIP v4. Ajuste del escalar $\tau$ por optimización LBFGS
sobre $107$ pares de validación; ECE medido sobre $101$ pares de
test. La calibración sobre v5 no converge (la optimización LBFGS
diverge sobre el espacio embebido \emph{image-LoRA} expandido y
$\tau \to 0$), por lo que queda registrada como limitación
metodológica abierta.}
\label{tab:spanderm-calibration}
\begin{tabular}{lrrr}
\hline
Nivel & ECE antes  & ECE después & $\tau$       \\
\hline
L3 (v4) & $0{,}7262$ & $0{,}0213$  & $0{,}4165$ \\
\hline
\end{tabular}
\end{table}

\begin{figure}[H]
\centering
\includegraphics[width=0.88\textwidth]{fig_spanderm_calibration.png}
\caption{Diagrama de fiabilidad sobre L3 antes y después de
\emph{temperature scaling} ($\tau = 0{,}4165$) sobre SpanDerm-CLIP
v4. ECE desciende de $0{,}7262$ a $0{,}0213$ ($-97\%$). La
calibración escalar de un parámetro restaura la interpretabilidad
probabilística de los logits sin alterar el \emph{ranking} entre
clases.}
\label{fig:spanderm-calibration}
\end{figure}

La reducción es del $-97\%$ sobre L3 con un único parámetro
escalar~\cite{guo2017}. La implicación se conecta con la
sensibilidad de DermLIP documentada en
secci\'on~\ref{sec:zs-results}: sin calibración post-hoc, las
probabilidades de un sistema CLIP-style \emph{zero-shot} castellano
no son interpretables como confianza clínica. Cualquier despliegue
posterior debe incorporar este escalado como capa final, ajustado
sobre datos del propio archivo de destino. La extensión del mismo
procedimiento a v5 no converge: la optimización LBFGS sobre el
espacio embebido \emph{image-LoRA} expandido lleva $\tau$ hacia
cero y el ECE empeora respecto a la baseline sin calibrar. La
limitación no afecta la integridad del resultado paper-grade
asociado a v4 y se documenta como observación metodológica para
iteraciones futuras: el procedimiento de calibración debe acotar
el rango de búsqueda o emplear búsqueda en rejilla cuando el
espacio embebido se ha expandido sustancialmente respecto al
\emph{checkpoint} sobre el que se calibró originalmente.

\subsection*{Trade-off entre escala externa y especificidad de
dominio: ablación v4 vs.~v5}
\label{sec:spanderm-escala-especificidad}

La quinta iteración disocia experimentalmente las dos palancas
candidatas para superar el techo \emph{zero-shot} L3 del encoder
visual sobre el corpus castellano: capacidad arquitectónica del
\emph{image encoder} (fijada en v4) y escala de datos externos
con \emph{captions} sintéticas (introducida en v5 manteniendo la
misma configuración LoRA). El resultado, reportado en la
tabla~\ref{tab:spanderm-zs}, es \emph{flat-to-negative} sobre el
nivel L3 \emph{zero-shot} ($0{,}2475 \to 0{,}2079$, $-4{,}0$\,pp).
La descomposición por régimen de inferencia ilumina la asimetría
del efecto: bajo régimen visual puro \emph{H1} ---logits
calculados exclusivamente sobre la similitud entre embedding
visual de test y prototipos visuales castellanos--- la transición
v4$\to$v5 mejora la \emph{accuracy}@1 L3 de $0{,}168$ a $0{,}238$,
un $+41{,}2\%$ relativo, lo que confirma que el corpus externo
fortalece el encoder visual aisladamente; pero bajo régimen
híbrido \emph{H6} ---combinación lineal con el lado textual
sobre plantillas ricas--- la cifra global cae. El techo
\emph{Linear Probing} L2 supervisado, evaluado con el mismo
script empleado para v2-frozen y v3, se eleva $+54{,}4\%$ relativo
en BAcc de v4 a v5 ($+36{,}8\%$ en \emph{accuracy}@1, métrica
auxiliar; se reporta la BAcc como métrica principal por el
desbalance severo de la distribución L2 del corpus castellà),
lo que constata que las \emph{features}
visuales sí han ganado capacidad discriminativa en el espacio
$1\,024$-D.

La interpretación que reconcilia ambas observaciones es que las
\emph{captions} sintéticas generadas sobre el corpus externo
genérico (PAD-UFES-20~\cite{padufes}, DermNet~\cite{dermnet}) no
preservan el registro discursivo del archivo
Dermapixel~1.0: el modelo aprende a alinear los \emph{embeddings}
de imagen con un espacio textual mixto que diluye la
\emph{fingerprint} clínica castellana sobre la que se calibra el
régimen \emph{H6}. La pérdida de alineación supera, en términos
de \emph{accuracy}@1 L3, la ganancia visual aislada. El hallazgo
es metodológicamente convergente con el documentado en
secci\'on~\ref{sec:zs-results} sobre Derm1M: el enriquecimiento de
\emph{prompts} con material léxico no alineado con el corpus de
preentrenamiento empeora la AUROC \emph{zero-shot}, en lugar de
mejorarla por incremento de información léxica. La regularidad
empírica que ambos resultados sugieren es que la alineación
cross-modal específica de un dominio clínico de alta
especificidad lingüística constituye un recurso frágil: pequeñas
perturbaciones estilísticas del lado textual pueden trasladar la
distribución conjunta lo suficiente como para erosionar la
discriminación final aunque las representaciones marginales
mejoren.

Una segunda lectura cuantitativa refuerza el diagnóstico. La
ampliación del corpus en v5 ($854 \to 15\,982$ imágenes, factor
$\approx 19\times$) no rellena la cola larga estructural del
nivel L3: las $327$ clases con $\leq 5$ imágenes en el
subconjunto castellano original mantienen su cardinalidad tras
la incorporación de las imágenes externas, dado que el
vocabulario externo no se mapea sobre las entradas L3 finas del
archivo Taberner. Las ocho clases L3 críticas $\emph{deepened}$
en v5 ---seis tumorales más dermatitis atópica y urticaria---
ganan cobertura significativa de muestras, pero el grueso de la
cola larga permanece inactiva. La implicación operativa para el
diseño de futuras iteraciones es directa y se desarrolla en
secci\'on~\ref{sec:linea-active-learning}: superar
\emph{zero-shot} L3 $\sim 0{,}25$ sobre vocabulario castellano
fino requiere anotación experta dirigida \emph{in-domain} sobre
las clases con peor cobertura, no escala genérica
indiscriminada.

Una ablación independiente con un modelo fundacional inglés grande
preentrenado (DermLIP v2 sobre $\approx 403\,000$ pares
Derm1M~\cite{derm1m}) confirma cuantitativamente este
diagnóstico: el \emph{wrapper} inglés con traducción de etiquetas
alcanza \emph{accuracy}@1 L3 $= 0{,}2871$ en régimen \emph{H6}
comparable, frente a $0{,}2475$ de SpanDerm-CLIP v4
($+3{,}96$\,pp), y $0{,}3861$ en su mejor régimen de plantillas
más texto rico ($+13{,}86$\,pp), evidencia consistente con que la
palanca limitante era la escala de la alineación cross-modal y no
la formulación textual ni el codificador visual. El análisis
completo, junto con la cuantificación del eje idioma puro sobre el
propio \emph{checkpoint} v4 ($0{,}2475$ castellà frente a
$0{,}2277$ inglés bajo \emph{H6}) y el \emph{trade-off} operativo
entre ambas estrategias (rendimiento bruto frente a privacidad,
especificidad lingüística y autonomía de despliegue), se
desarrolla en el Anexo~\ref{cap:anexoi}.

\subsection*{Diagnóstico cuantitativo del techo experimental
sobre L2}

La tabla~\ref{tab:spanderm-l2-ceiling} sitúa SpanDerm-CLIP v2 en
\emph{zero-shot} L2 frente a las dos referencias internas
documentadas previamente en el capítulo sobre el mismo conjunto
de test ($N_{\mathrm{test}} = 36$ en
secci\'on~\ref{sec:dermapixel-results} y
secci\'on~\ref{sec:dermapixel-spanderm-v0}).

\begin{table}[H]
\centering
\footnotesize
\setlength{\tabcolsep}{4pt}
\renewcommand{\arraystretch}{0.95}
\caption{SpanDerm-CLIP \emph{zero-shot} L2 frente a las dos
referencias internas del capítulo sobre el mismo nivel ontológico.
Las tres referencias evalúan sobre el conjunto de test de
DermapixelAI~1.0; la cifra de SpanDerm-CLIP v2 corresponde al
test de $101$ imágenes del \emph{split} per-case del experimento
SpanDerm-CLIP y las dos restantes al test fijo
$N_{\mathrm{test}} = 36$ de la secci\'on~\ref{sec:dermapixel-results}.
La comparación es indicativa por la diferencia muestral.}
\label{tab:spanderm-l2-ceiling}
\begin{tabular}{llrr}
\hline
Sistema                                                 & Régimen               & Acc@1 L2          & Observaciones \\
\hline
LP PanDerm Large (referencia, tabla~\ref{tab:dermapixel-lp}) & Supervisado L2 cerrado & $0{,}250$        & $38$ clases, $N_{\mathrm{test}} = 36$ \\
SpanDerm v0 LoRA (tabla~\ref{tab:dermapixel-spanderm-l2}) & Supervisado L2 cerrado & $0{,}250$        & $38$ clases, $N_{\mathrm{test}} = 36$ \\
\textbf{SpanDerm-CLIP v2}                                & \emph{Zero-shot} open  & $\mathbf{0{,}133}$ & $23$ clases efectivas, $N_{\mathrm{test}} = 101$ \\
\hline
\end{tabular}
\end{table}

La lectura cuantitativa es coherente con la jerarquía esperada
entre los tres regímenes: la cabeza supervisada cerrada con LoRA
sobre el mismo encoder (SpanDerm v0) alcanza Acc@1 $0{,}250$
sobre L2, mientras que el sistema \emph{zero-shot open-class}
(SpanDerm-CLIP v2) alcanza $0{,}133$ sobre un vocabulario reducido
a las $23$ clases efectivas en el test del \emph{split}
per-case del experimento. La diferencia operativa entre ambos
regímenes es la capacidad de SpanDerm-CLIP de aceptar consultas
textuales libres en castellano sin necesidad de fijar el
vocabulario por anticipado, propiedad explotada en el módulo
RAG castellano del prototipo (Anexo~\ref{cap:anexoh}).

\subsection*{Síntesis: progresión v1$\to$v4 y ablación negativa v5}

La tabla~\ref{tab:spanderm_iteraciones} sintetiza la trayectoria
completa de la familia SpanDerm-CLIP sobre el conjunto de test
de DermapixelAI~1.0, articulando para cada iteración la
configuración LoRA, el régimen de inferencia, las dos cifras de
referencia ($L3$ \emph{accuracy}@1 \emph{zero-shot} y lift de
recuperación cross-modal $i2t$ R@$5$) y la lectura cualitativa
asociada.

\begin{table}[H]
\centering
\footnotesize
\setlength{\tabcolsep}{3pt}
\renewcommand{\arraystretch}{1.05}
\caption{Comparativa de las cinco iteraciones del entrenamiento
SpanDerm-CLIP sobre el corpus DermapixelAI~1.0. Régimen evaluativo
\emph{zero-shot} L3 sobre $101$ imágenes de test ($54$ clases L3
representadas), \emph{H6 hybrid} v$30$\_r$70$ (combinación
$30\%$ prototipos visuales $+$ $70\%$ texto rico) aplicado de
forma homogénea a las cinco iteraciones. v4 alcanza el objetivo
paper-grade $0{,}25$ acc@1 L3 mediante expansión arquitectónica
(\emph{image-LoRA} $r=32$ sobre $12$ bloques $+$ \emph{ensemble}
híbrido texto-visual). v5 confirma negativamente que la escala
externa genérica con \emph{captions} sintéticas no supera v4. Para
v1--v3 los prototipos visuales derivan del espacio embebido del
\emph{checkpoint} sin LoRA visual; v4--v5 emplean prototipos
sobre el espacio embebido con \emph{image-LoRA} expandido. La
comparativa permanece interpretable porque los prototipos derivan
del propio espacio embebido del \emph{checkpoint} evaluado.
Provenance: commit \texttt{d0190c3} del repositorio
\texttt{spanderm-clip}. Las cifras Linear Probing BAcc L2 absolutas
para v4 ($0{,}1860$) y v5 ($0{,}2871$) proceden de los ficheros
JSON \texttt{eval\_v5/linear\_probe\_v4ref.json} y
\texttt{eval\_v5/linear\_probe\_v5.json}; el delta relativo
$+54{,}4\%$ queda confirmado.}
\label{tab:spanderm_iteraciones}
\begin{tabular}{lllrrll}
\hline
Versión & LoRA texto & LoRA visual & acc@1 L3 & Lift i2t R@$5$ & LP BAcc L2 & Notas \\
\hline
v1 & ---                       & ---                          & $0{,}1782$          & ---                       & ---                                 & Baseline e5 \emph{full FT} \\
v2 & $r=16$, $\alpha=32$       & ---                          & $0{,}1980$          & $\mathbf{14{,}0\times}$   & ---                                 & Pico retrieval rich \\
v3 & $r=16$, $\alpha=32$       & $r=16$, últ.~$4/24$           & $0{,}1980$          & ---                       & ---                                 & Image-LoRA infradotado \\
v4 & $r=32$, $\alpha=64$       & $r=32$, últ.~$12/24$          & $\mathbf{0{,}2475}$ & $12{,}4\times$            & $0{,}1860$   & Target paper-grade \\
v5 & $r=32$, $\alpha=64$       & $r=32$, últ.~$12/24$          & $0{,}2079$          & $10{,}8\times$            & $0{,}2871$   & FLAT: escala externa ($+54{,}4\%$ relativo sobre v4) \\
\hline
\end{tabular}

\smallskip
\footnotesize
\textit{Nota histórica.} Los valores \emph{rich-only baseline} (un
\emph{caption} clínico por clase $\sim 30$--$80$ palabras
castellano) fueron $0{,}1287$ (v1), $0{,}1386$ (v2) y $0{,}1089$
(v3) en el régimen \texttt{baseline\_5\_templates} original de la
evaluación previa. La transición al régimen \emph{H6 hybrid}
homogéneo presentado en esta tabla permite la comparativa directa
cross-\emph{checkpoint} mediante el mismo procedimiento
evaluativo (prototipos visuales $30\%$ $+$ texto rico $70\%$). El
detalle metodológico se desarrolla en
Anexo~\ref{cap:anexog}.
\normalsize
\end{table}

La lectura conjunta sobre las cinco iteraciones articula tres
ejes de mejora con valoración empírica diferenciada. El eje de
capacidad arquitectónica queda validado por v4: la transición
v3$\to$v4 (LoRA $r=16$ sobre cuatro bloques $\to$ $r=32$ sobre
doce bloques) eleva el L3 \emph{zero-shot} de $0{,}1980$ a
$0{,}2475$ bajo régimen \emph{H6} homogéneo, $+4{,}95$\,pp, y la
mejora global v1$\to$v4 alcanza el $\sim 39\%$ relativo
($0{,}1782 \to 0{,}2475$). La hipótesis formulada tras la medición
del LP en v3 ---el cuello de botella reside en la capacidad del
image encoder--- queda confirmada cuantitativamente bajo la
solución arquitectónica correcta. El eje de escala externa queda
\emph{flat-to-negative} en v5: bajo configuración LoRA idéntica a
v4, la ampliación del corpus con $15\,128$ imágenes externas
genéricas no replica la mejora pese a que las representaciones
visuales aisladas mejoran ($+41{,}2\%$ relativo H1, $+54{,}4\%$
relativo LP L2). La conclusión metodológica es que la escala
externa genérica con \emph{captions} sintéticas erosiona la
\emph{fingerprint} clínica específica del registro castellano
original, en convergencia con la sensibilidad de DermLIP al
tokenizador documentada en secci\'on~\ref{sec:zs-results}. El eje de
supervisión conceptual ---anotaciones del \emph{Seven-Point
Checklist} sobre Derm7pt~\cite{kawahara2019} o anotaciones del
Grupo~A de la Dra.~R.~Taberner (Anexo~\ref{cap:anexof})---
queda abierto a futuras iteraciones del proyecto y se desarrolla
como línea prioritaria en
secci\'on~\ref{sec:linea-active-learning} bajo la formulación
de aprendizaje activo dirigido sobre las clases L3 con peor
cobertura.

El experimento aporta cuatro contribuciones al cuerpo del
trabajo. En primer lugar, valida la utilidad del codificador
PanDerm Large para tareas castellanas más allá del LP/FT
supervisado documentado en las secciones previas: bajo la
combinación de LoRA $r=32$ sobre los últimos doce bloques del
ViT y régimen híbrido visual-textual \emph{H6} en inferencia, el
sistema alcanza esencialmente el umbral paper-grade $L3$
$0{,}25$ sobre vocabulario castellano fino. En segundo lugar,
disocia experimentalmente capacidad arquitectónica de escala de
datos sobre el mismo corpus: la primera demuestra ser palanca
efectiva (v3$\to$v4) y la segunda \emph{flat-to-negative}
(v4$\to$v5), un resultado negativo honesto que orienta el diseño
de futuras iteraciones. En tercer lugar, establece metodología
reproducible para evaluaciones CLIP-style sobre archivos
privados de pequeña escala: recuperación per-imagen (no
per-pair, sesgada sobre corpus con multiplicidad alta de
captions), medición LP consistente entre versiones, régimen de
inferencia \emph{H6} con calibración del peso visual sobre
validación independiente, \emph{ensemble} de plantillas
castellanas, calibración escalar post-hoc~\cite{guo2017},
manifest SHA-256 fijado en cada \emph{checkpoint} y
\emph{splits} per-case estratificados. En cuarto lugar, aporta
material empírico publicable: la recuperación cross-modal
$\mathbf{14{,}0\times}$ lift sobre random con captions ricas
castellanas (v2, $R@5 = 0{,}654$ frente a $0{,}047$), calibración
escalar $-97\%$ ECE sobre L3 (v4, $0{,}7262 \to 0{,}0213$ con
$\tau = 0{,}4165$), umbral L3 \emph{zero-shot} $0{,}2475$ sobre
vocabulario castellano fino bajo configuración arquitectónica
correcta (v4) y ablación negativa documentada de la hipótesis de
escala externa indiscriminada (v5) constituyen una unidad
coherente que ningún trabajo previo ha publicado en castellano
sobre archivo privado clínico. El sistema validado en v4 es
candidato a integrarse como módulo de búsqueda textual castellano
del prototipo DermApIxel (Anexo~\ref{cap:anexoh}), complementario
al módulo M4 sobre Derm1M en inglés, con la calibración escalar
$\tau = 0{,}4165$ ajustada sobre el conjunto de validación del
propio archivo de destino como capa final antes de exponer
probabilidades al usuario.
